% Editable master: glossaries/old-english.md
% Complete export; generated by scripts/build_glossary_exports.py.
\newcommand{\runningtitle}{Old English–English}
% Layout inherited from the reviewed previous course; set a running title per unit.
% This is a small hand-maintained TeX layout, not a conversion pipeline.
\RequirePackage{fix-cm}
\documentclass[11pt,twoside]{article}
\usepackage{fontspec}
\setmainfont{Cambria}[
  Renderer=HarfBuzz,
  BoldFont=cambriab.ttf,
  ItalicFont=cambriai.ttf,
  BoldItalicFont=cambriaz.ttf
]
\setmonofont{lmmono10-regular.otf}[BoldFont=lmmonolt10-bold.otf]
\usepackage[paperwidth=165mm,paperheight=241mm,inner=21mm,outer=19mm,
  top=22mm,bottom=20mm,headheight=14pt,headsep=7mm,footskip=10mm]{geometry}
\usepackage{xcolor,array,booktabs,longtable,tabularx}
\usepackage{enumitem,fancyhdr}
\usepackage[unicode,hidelinks]{hyperref}
\definecolor{ink}{HTML}{000000}
\definecolor{muted}{HTML}{4C5558}
\setlength{\parindent}{0pt}
\setlength{\parskip}{3.5pt}
\AtBeginDocument{\fontsize{11}{14.3}\selectfont}
\setlength{\tabcolsep}{4pt}
\renewcommand{\arraystretch}{1.10}
\setlist[itemize]{leftmargin=1.3em,label=\textbullet,itemsep=2pt,topsep=3pt,parsep=0pt}
\setlist[enumerate]{leftmargin=1.8em,itemsep=3pt,topsep=4pt,parsep=0pt}
\widowpenalty=10000
\clubpenalty=10000
\emergencystretch=1em
\raggedbottom
\pagestyle{fancy}
\fancyhf{}
\providecommand{\runningtitle}{Old English}
\fancyhead[LE,RO]{\small\itshape \runningtitle}
\fancyhead[LO,RE]{\small Old English}
\fancyfoot[LE,RO]{\small\thepage}
\renewcommand{\headrulewidth}{0.3pt}
\newcommand{\booktitle}[2]{
  {\fontsize{9}{11}\selectfont\color{muted}#1\par}
  \vspace{3pt}
  {\fontsize{20}{24}\selectfont\color{ink}#2\par}
  \vspace{6pt}
}
\newcommand{\source}[1]{\par{\footnotesize\color{muted}#1\par}}
\newcommand{\readingtext}[1]{
  \par\addvspace{5pt}
  \begingroup\leftskip=1.5em
    \fontsize{12}{16}\selectfont #1\par
  \endgroup\addvspace{5pt}
}
\makeatletter
\renewcommand{\section}{\@startsection{section}{1}{\z@}
  {-2.4ex \@plus -0.6ex \@minus -0.2ex}{0.8ex \@plus 0.2ex}
  {\normalfont\fontsize{12}{15}\selectfont\bfseries\color{ink}}}
\renewcommand{\subsection}{\@startsection{subsection}{2}{\z@}
  {-1.8ex \@plus -0.5ex \@minus -0.2ex}{0.5ex \@plus 0.2ex}
  {\normalfont\fontsize{11}{14}\selectfont\bfseries\color{ink}}}
\makeatother

\usepackage{multicol}
\raggedcolumns
\newcommand{\needspace}[1]{\par\begingroup\skip0=\lastskip\vskip-\skip0\dimen0=#1\relax\advance\dimen0 by\skip0\vskip0pt plus\dimen0\penalty-100\vskip0pt plus-\dimen0\vskip\dimen0\penalty9999\vskip-\dimen0\vskip\skip0\endgroup}
\setlength{\columnsep}{6mm}
\hyphenpenalty=10000\exhyphenpenalty=10000
\emergencystretch=2em
\begin{document}
\addfontfeatures{Ligatures=TeXOff}
\raggedright
\pdfbookmark[0]{Old English–English glossary}{outline-B00001}
% glossary-block:B00001:begin
\booktitle{GLOSSARY}{Old English–English glossary}
% glossary-block:B00001:end
\fontsize{10.5}{13.2}\selectfont
% glossary-block:B00002:begin
This glossary covers the Old English in all eighteen units and their keys, including tables, readings, alternatives and source notes, and in the pronunciation guide and its key. It is a reference, not a new memorization list. The exercise instructions still determine what you must know and what is supplied.\par
% glossary-block:B00002:end
% glossary-block:B00003:begin
Look under the dictionary headword, normally the infinitive for a verb. Strong verbs give the preterite first singular, preterite plural and past participle; weak verbs give the preterite and participle. Distinctive presents and irregular verbs have their forms labelled. Selected difficult forms, including vowel-changing third-person singulars, have \textbf{see} entries. Routine inflections belong to the unit paradigms.\par
% glossary-block:B00003:end
% glossary-block:B00004:begin
Nouns give gender, genitive singular and nominative plural where established. Adjectives and adverbs give selected comparison forms or useful stem clues. Words appearing only in a source note may have a limited citation rather than a complete set of principal parts. A missing cell supplies no form by implication. Meanings and construction notes retain the limits of the cited use.\par
% glossary-block:B00004:end
% glossary-block:B00005:begin
The order is a, æ, b … t, þ/ð, u … z. Length marks and hooks do not affect alphabetization; they remain essential parts of the printed spelling. Thorn and eth share a sorting position. Words beginning ge- are under g. Ordinary forms are listed in lowercase; proper-name capitalization is retained. Sentence-initial capitals do not create separate entries.\par
% glossary-block:B00005:end
% glossary-block:B00006:begin
Abbreviations: n. noun; v. verb; adj. adjective; adv. adverb; prep. preposition; conj. conjunction; pron. pronoun; interj. interjection; m., f., n. masculine, feminine, neuter; gen. genitive; nom. nominative; pres. present; pret. preterite; ptcp. participle; sg. singular; pl. plural; comp. comparative; superl. superlative. Numbers 1 and 3 identify persons. Parenthesized (ge) preserves the source's optional prefix notation; it does not license adding ge- to other words.\par
% glossary-block:B00006:end
% glossary-block:B00007:begin
Case and construction notes qualify the meaning beside them. A supplied expression is not a general word-for-word translation rule. Source spellings corrected in the course, deliberately wrong examples and reconstructed forms are identified separately at the end.\par
% glossary-block:B00007:end
% glossary-block:B00008:begin
The definitions and spelling follow the course's checked Cook sources, with the recorded Fulk, Wright and dictionary support where used. Principal parts include checked reference formations from the stated grammatical rules as well as exact source citations. This glossary is newly arranged reference material. The companion English–Old English glossary indexes exercise meanings and constructions.\par
% glossary-block:B00008:end
\pdfbookmark[0]{Entries}{outline-B00009}
\clearpage
\fontsize{10.5}{13.2}\selectfont
\setlength{\parskip}{3.8pt}
\begin{multicols*}{2}
% glossary-block:B00009:begin
\needspace{5\baselineskip}
\section*{Entries}
% glossary-block:B00009:end
\pdfbookmark[1]{A}{outline-B00010}
% glossary-block:B00010:begin
\needspace{5\baselineskip}
\subsection*{A}
% glossary-block:B00010:end
% glossary-block:B00011:begin
\textbf{ā} (adv.). ever; always.\par
% glossary-block:B00011:end
% glossary-block:B00012:begin
\needspace{3\baselineskip}
\textbf{ābēodan} (v.). announce; bid a greeting (hǣl ābēodan). hǣl ābēodan: bid someone hail / give someone greeting.\par
% glossary-block:B00012:end
% glossary-block:B00013:begin
\needspace{3\baselineskip}
\textbf{āberan} (v.). carry. Supplied carried form; retain the ā-prefix.\par
% glossary-block:B00013:end
% glossary-block:B00014:begin
\needspace{3\baselineskip}
\textbf{Abraham} (proper name). Abraham. Biblical patriarch; supplied name forms.\par
% glossary-block:B00014:end
% glossary-block:B00015:begin
\needspace{3\baselineskip}
\textbf{ābūtan} (prep.). around; about. With accusative.\par
% glossary-block:B00015:end
% glossary-block:B00016:begin
\needspace{3\baselineskip}
\textbf{āc} (n.). oak. Pronunciation example.\par
% glossary-block:B00016:end
% glossary-block:B00017:begin
\needspace{3\baselineskip}
\textbf{ac} (conj.). but.\par
% glossary-block:B00017:end
% glossary-block:B00018:begin
\needspace{3\baselineskip}
\textbf{ācęnnan} (v.); past ptcp. ācenned, ācęnned. produce; bring forth. beget; bring forth.\par
% glossary-block:B00018:end
% glossary-block:B00019:begin
\needspace{3\baselineskip}
\textbf{āceorfan} (v. strong); pret. 1 sg. ācearf; pret. pl. ācurfon; past ptcp. ācorfen; pres. 3 sg. ācierfð. cut off. Accusative thing cut; of in the assigned expression is a separate off particle.\par
% glossary-block:B00019:end
% glossary-block:B00020:begin
\needspace{3\baselineskip}
\textbf{ācierfð} — pres. ind. 3 sg.: see \textbf{āceorfan}.\par
% glossary-block:B00020:end
% glossary-block:B00021:begin
\needspace{3\baselineskip}
\textbf{adesa} (n. m.); gen. sg. adesan; nom. pl. adesan. hatchet.\par
% glossary-block:B00021:end
% glossary-block:B00022:begin
\needspace{3\baselineskip}
\textbf{ādrīfan} (v. strong); pret. 1 sg. ādrāf; pret. pl. ādrifon; past ptcp. ādrifen. drive out; expel. Accusative person expelled; fram + dative may mark the point of separation.\par
% glossary-block:B00022:end
% glossary-block:B00023:begin
\needspace{3\baselineskip}
\textbf{āfǣran} (v.); past ptcp. āfǣred. frighten.\par
% glossary-block:B00023:end
% glossary-block:B00024:begin
\needspace{3\baselineskip}
\textbf{āfēdan} (v. weak); pret. 1 sg. āfēdde; past ptcp. āfēded. support; nourish. Accusative person supported.\par
% glossary-block:B00024:end
% glossary-block:B00025:begin
\needspace{3\baselineskip}
\textbf{āfierran} (v.). remove; take away.\par
% glossary-block:B00025:end
% glossary-block:B00026:begin
\needspace{3\baselineskip}
\textbf{āflīeman} (v.). put to flight; expel.\par
% glossary-block:B00026:end
% glossary-block:B00027:begin
\needspace{3\baselineskip}
\textbf{āfyllan} (v. weak); pret. 1 sg. āfylde; past ptcp. āfylled. fill.\par
% glossary-block:B00027:end
% glossary-block:B00028:begin
\needspace{3\baselineskip}
\textbf{āgan} (v. preterite-present); pres. 1/3 sg. āh; pres. pl. āgon; pret. 1 sg. āhte; adjectival ptcp. āgen. have; possess. Accusative possession; positive present and preterite indicative in the taught scope. Āgen is adjectival: “own”.\par
% glossary-block:B00028:end
% glossary-block:B00029:begin
\needspace{3\baselineskip}
\textbf{āgān} (v.); past ptcp. āgān. go away. Supplied gone-away form; distinct from āgan possess.\par
% glossary-block:B00029:end
% glossary-block:B00030:begin
\needspace{3\baselineskip}
\textbf{āgēan} — citation/spelling variant: see \textbf{ongēan}.\par
% glossary-block:B00030:end
% glossary-block:B00031:begin
\needspace{3\baselineskip}
\textbf{āgen} (adj.). own.\par
% glossary-block:B00031:end
% glossary-block:B00032:begin
\needspace{3\baselineskip}
\textbf{āgiefan} (v.). give back; restore.\par
% glossary-block:B00032:end
% glossary-block:B00033:begin
\needspace{3\baselineskip}
\textbf{āginnan} (v.). begin. Supplied source form began.\par
% glossary-block:B00033:end
% glossary-block:B00034:begin
\needspace{3\baselineskip}
\textbf{āgyldan} (v.); pres. 3 sg. āgylt. reward; repay.\par
% glossary-block:B00034:end
% glossary-block:B00035:begin
\needspace{3\baselineskip}
\textbf{āh} — pres. ind. 1/3 sg.: see \textbf{āgan}.\par
% glossary-block:B00035:end
% glossary-block:B00036:begin
\needspace{3\baselineskip}
\textbf{āhębban} (v. strong); pret. 1 sg. āhōf; pret. pl. āhōfon; past ptcp. āhafen; pres. 3 sg. āhęfð. raise; lift.\par
% glossary-block:B00036:end
% glossary-block:B00037:begin
\needspace{3\baselineskip}
\textbf{āhęfð} — pres. ind. 3 sg.: see \textbf{āhębban}.\par
% glossary-block:B00037:end
% glossary-block:B00038:begin
\needspace{3\baselineskip}
\textbf{āhrēosan} (v. strong); pret. 1 sg. āhrēas; pret. pl. āhruron; past ptcp. āhroren; pres. 3 sg. āhrīest. fall; fall upon.\par
% glossary-block:B00038:end
% glossary-block:B00039:begin
\needspace{3\baselineskip}
\textbf{āhrīest} — pres. ind. 3 sg.: see \textbf{āhrēosan}.\par
% glossary-block:B00039:end
% glossary-block:B00040:begin
\needspace{3\baselineskip}
\textbf{āhte} — pret. sg.: see \textbf{āgan}.\par
% glossary-block:B00040:end
% glossary-block:B00041:begin
\needspace{3\baselineskip}
\textbf{āhwā} (indefinite pronoun). anyone. Reference in a source or explanatory note.\par
% glossary-block:B00041:end
% glossary-block:B00042:begin
\needspace{3\baselineskip}
\textbf{āhwæðer} (indefinite pronoun/adjective). either; one of two. Reference in a source or explanatory note.\par
% glossary-block:B00042:end
% glossary-block:B00043:begin
\needspace{3\baselineskip}
\textbf{āhyldan} (v.). lay; incline.\par
% glossary-block:B00043:end
% glossary-block:B00044:begin
\needspace{3\baselineskip}
\textbf{ālīefan} (v.); pres. 3 sg. ālīefð. permit; allow. Impersonal permission frame, with dative permitted person; consult the complete supplied expression.\par
% glossary-block:B00044:end
% glossary-block:B00045:begin
\needspace{3\baselineskip}
\textbf{ālīehtan} (v. weak); pret. 1 sg. ālīehte; past ptcp. ālīeht. illuminate; give light to. Accusative person or thing illuminated.\par
% glossary-block:B00045:end
% glossary-block:B00046:begin
\needspace{3\baselineskip}
\textbf{ālīesan} (v. weak); pret. 1 sg. ālīesde; past ptcp. ālīesed. deliver; release. Accusative person delivered or released; fram + dative marks separation.\par
% glossary-block:B00046:end
% glossary-block:B00047:begin
\needspace{3\baselineskip}
\textbf{ālīesend} (n. m.); gen. sg. ālīesendes; nom. pl. ālīesend. Redeemer; one who delivers.\par
% glossary-block:B00047:end
% glossary-block:B00048:begin
\needspace{3\baselineskip}
\textbf{ālūtan} (v. strong); pret. 1 sg. ālēat; pret. pl. āluton; past ptcp. āloten; pres. 3 sg. ālȳtt. bow; incline oneself. Reflexive meaning in the assigned use.\par
% glossary-block:B00048:end
% glossary-block:B00049:begin
\needspace{3\baselineskip}
\textbf{ālȳtt} — pres. ind. 3 sg.: see \textbf{ālūtan}.\par
% glossary-block:B00049:end
% glossary-block:B00050:begin
\needspace{3\baselineskip}
\textbf{Amen} (interj.). Amen.\par
% glossary-block:B00050:end
% glossary-block:B00051:begin
\needspace{3\baselineskip}
\textbf{ān} (numeral; adj.). one. alone. a certain. Introduces an unnamed singular individual; not an obligatory indefinite article. one; single.\par
% glossary-block:B00051:end
% glossary-block:B00052:begin
\needspace{3\baselineskip}
\textbf{āncęnned} (adj.). only-begotten.\par
% glossary-block:B00052:end
% glossary-block:B00053:begin
\needspace{3\baselineskip}
\textbf{and} (conj.). and. Simple coordination.\par
% glossary-block:B00053:end
% glossary-block:B00054:begin
\needspace{3\baselineskip}
\textbf{andgiet} (n. n.); gen. sg. andgietes. meaning; sense. understanding.\par
% glossary-block:B00054:end
% glossary-block:B00055:begin
\needspace{3\baselineskip}
\textbf{andgietfullice} (adv.); superl. andgietfullicost. In superlative andgietfullicost: most intelligibly.\par
% glossary-block:B00055:end
% glossary-block:B00056:begin
\needspace{3\baselineskip}
\textbf{andgietfullicost} — superlative adverb; exact source spelling: see \textbf{andgietfullice}.\par
% glossary-block:B00056:end
% glossary-block:B00057:begin
\needspace{3\baselineskip}
\textbf{andlīcnes} (n. f.); gen. sg. andlīcnesse; nom. pl. andlīcnessa. image.\par
% glossary-block:B00057:end
% glossary-block:B00058:begin
\needspace{3\baselineskip}
\textbf{andswarian} (v. weak); pret. 1 sg. andswarode; past ptcp. andswarod. answer; reply. Dative addressee; the reply may be given as direct speech.\par
% glossary-block:B00058:end
% glossary-block:B00059:begin
\needspace{3\baselineskip}
\textbf{andswerian} — citation/spelling variant: see \textbf{andswarian}.\par
% glossary-block:B00059:end
% glossary-block:B00060:begin
\needspace{3\baselineskip}
\textbf{andweard} (adj.). present.\par
% glossary-block:B00060:end
% glossary-block:B00061:begin
\needspace{3\baselineskip}
\textbf{andwlita} (n. m.); gen. sg. andwlitan; nom. pl. andwlitan. face; countenance.\par
% glossary-block:B00061:end
% glossary-block:B00062:begin
\needspace{3\baselineskip}
\textbf{angel} (n. m.); gen. sg. angles; nom. pl. anglas. fish-hook.\par
% glossary-block:B00062:end
% glossary-block:B00063:begin
\needspace{3\baselineskip}
\textbf{anginn} (n. n.); gen. sg. anginnes; nom. pl. anginn. beginning.\par
% glossary-block:B00063:end
% glossary-block:B00064:begin
\needspace{3\baselineskip}
\textbf{angrīslic} (adj.). fierce; raging.\par
% glossary-block:B00064:end
% glossary-block:B00065:begin
\needspace{3\baselineskip}
\textbf{ann} — pres. ind. 1/3 sg.: see \textbf{unnan}.\par
% glossary-block:B00065:end
% glossary-block:B00066:begin
\needspace{3\baselineskip}
\textbf{ansīen} (n. f.); gen. sg. ansīene. appearance.\par
% glossary-block:B00066:end
% glossary-block:B00067:begin
\needspace{3\baselineskip}
\textbf{Apollines} (proper name; source nom. sg.). Apollo. The god, distinct from Apollonius. Apollo, distinct from Apollonius.\par
% glossary-block:B00067:end
% glossary-block:B00068:begin
\needspace{3\baselineskip}
\textbf{Apollonius} (proper name). Apollonius. Supplied dative person in the assigned source expression. Apollonius.\par
% glossary-block:B00068:end
% glossary-block:B00069:begin
\needspace{3\baselineskip}
\textbf{ār} (n. f.). honour; high position. Contextual rank or office, not a new general meaning of every ār form.\par
% glossary-block:B00069:end
% glossary-block:B00070:begin
\needspace{3\baselineskip}
\textbf{Arcestrates} (proper name). Arcestrates. King of the region in the story; supplied name forms. Arcestrates. The king; supplied name forms.\par
% glossary-block:B00070:end
% glossary-block:B00071:begin
\needspace{3\baselineskip}
\textbf{ār-fæstnes} (n. f.). kindness.\par
% glossary-block:B00071:end
% glossary-block:B00072:begin
\needspace{3\baselineskip}
\textbf{ārīsan} (v. strong); pret. 1 sg. ārās; pret. pl. ārison; past ptcp. ārisen. arise. Intransitive arise; of + dative source in the bounded rise-from-death expression.\par
% glossary-block:B00072:end
% glossary-block:B00073:begin
\needspace{3\baselineskip}
\textbf{arn} — pret. sg.: see \textbf{iernan}.\par
% glossary-block:B00073:end
% glossary-block:B00074:begin
\needspace{3\baselineskip}
\textbf{āscian} (v. weak); pret. 1 sg. āscode; past ptcp. āscod. ask about. Accusative person asked, genitive matter asked about; an indirect question may express the matter. ask about. Supplied source form asked about.\par
% glossary-block:B00074:end
% glossary-block:B00075:begin
\needspace{3\baselineskip}
\textbf{āsęndan} (v. weak); pret. 1 sg. āsęnde; past ptcp. āsęnd. send.\par
% glossary-block:B00075:end
% glossary-block:B00076:begin
\needspace{3\baselineskip}
\textbf{āsęttan} (v. weak); pret. 1 sg. āsętte; past ptcp. āsętt. set; place.\par
% glossary-block:B00076:end
% glossary-block:B00077:begin
\needspace{3\baselineskip}
\textbf{āsmēan} (v.). conceive. Exact supplied infinitive āsmēan in the heart-conceiving expression; not an independently assigned full paradigm.\par
% glossary-block:B00077:end
% glossary-block:B00078:begin
\needspace{3\baselineskip}
\textbf{assa} (n. m.); gen. sg. assan; nom. pl. assan. ass; donkey.\par
% glossary-block:B00078:end
% glossary-block:B00079:begin
\needspace{3\baselineskip}
\textbf{Asser} (proper name). Asser.\par
% glossary-block:B00079:end
% glossary-block:B00080:begin
\needspace{3\baselineskip}
\textbf{āstīgan} (v. strong); pret. 1 sg. āstāg; pret. pl. āstigon; past ptcp. āstigen. go aboard; board. Accusative ship boarded in the assigned frame. descend. Intransitive descend; of + dative departure point.\par
% glossary-block:B00080:end
% glossary-block:B00081:begin
\needspace{3\baselineskip}
\textbf{āstreahte} — pret. sg.: see \textbf{āstręccean}.\par
% glossary-block:B00081:end
% glossary-block:B00082:begin
\needspace{3\baselineskip}
\textbf{āstręccean} (v. weak); pret. 1 sg. āstreahte; past ptcp. āstreaht. prostrate oneself. Reflexive accusative; the before phrase retains its own case/construal.\par
% glossary-block:B00082:end
% glossary-block:B00083:begin
\needspace{3\baselineskip}
\textbf{astyrian} (v.). move; set in motion. Supplied source sense.\par
% glossary-block:B00083:end
% glossary-block:B00084:begin
\needspace{3\baselineskip}
\textbf{āð} (n. m.); gen. sg. āðes; nom. pl. āðas. oath.\par
% glossary-block:B00084:end
% glossary-block:B00085:begin
\needspace{3\baselineskip}
\textbf{āðęnian} (v.). direct; stretch toward. Directing the mind toward what is heard; supplied figurative phrase.\par
% glossary-block:B00085:end
% glossary-block:B00086:begin
\needspace{3\baselineskip}
\textbf{āðindan} (v.). swollen. Supplied weak participial adjective describing the sea.\par
% glossary-block:B00086:end
% glossary-block:B00087:begin
\needspace{3\baselineskip}
\textbf{āwęg} — citation/spelling variant: see \textbf{onwęg}.\par
% glossary-block:B00087:end
% glossary-block:B00088:begin
\needspace{3\baselineskip}
\textbf{āwęndednes} (n. f.); gen. sg. āwęndednesse; nom. pl. āwęndednessa. translation.\par
% glossary-block:B00088:end
% glossary-block:B00089:begin
\needspace{3\baselineskip}
\textbf{āweorpan} (v. strong); pret. 1 sg. āwearp; pret. pl. āwurpon; past ptcp. āworpen; pres. 3 sg. āwierpð. cast down; cast away.\par
% glossary-block:B00089:end
% glossary-block:B00090:begin
\needspace{3\baselineskip}
\textbf{āwierpð} — pres. ind. 3 sg.: see \textbf{āweorpan}.\par
% glossary-block:B00090:end
% glossary-block:B00091:begin
\needspace{3\baselineskip}
\textbf{āwiht} (indefinite substantive). anything.\par
% glossary-block:B00091:end
% glossary-block:B00092:begin
\needspace{3\baselineskip}
\textbf{āwrītan} (v. strong); pret. 1 sg. āwrāt; pret. pl. āwriton; past ptcp. āwriten. write.\par
% glossary-block:B00092:end
% glossary-block:B00093:begin
\needspace{3\baselineskip}
\textbf{āwyrcean} (v.). do; deserve. In the source, deserving through one's conduct; retain the complete local frame.\par
% glossary-block:B00093:end
\pdfbookmark[1]{Æ}{outline-B00094}
% glossary-block:B00094:begin
\needspace{5\baselineskip}
\subsection*{Æ}
% glossary-block:B00094:end
% glossary-block:B00095:begin
\textbf{ǣ} (n. f.). law.\par
% glossary-block:B00095:end
% glossary-block:B00096:begin
\needspace{3\baselineskip}
\textbf{æcer} (n. m.); gen. sg. æceres; nom. pl. æceras. field.\par
% glossary-block:B00096:end
% glossary-block:B00097:begin
\needspace{3\baselineskip}
\textbf{æcern} (n. n.); gen. sg. æcernes; nom. pl. æcernu. acorn.\par
% glossary-block:B00097:end
% glossary-block:B00098:begin
\needspace{3\baselineskip}
\textbf{ǣfen} (n. n.); gen. sg. ǣfenes; nom. pl. ǣfenu. evening.\par
% glossary-block:B00098:end
% glossary-block:B00099:begin
\needspace{3\baselineskip}
\textbf{ǣfre} (adv.). ever; always.\par
% glossary-block:B00099:end
% glossary-block:B00100:begin
\needspace{3\baselineskip}
\textbf{æfter} (prep.; conj. in fixed expressions). according to. With dative. after. Temporal preposition with dative. after. In the supplied finite-clause conjunction æfter ðǣm ðe. after; in search of. Person sought or inquired after; keep this lexical frame distinct from temporal after. after; following. Dative person followed, including motion after that person.\par
% glossary-block:B00100:end
% glossary-block:B00101:begin
\needspace{3\baselineskip}
\textbf{ǣghwā} (indefinite pronoun). everyone; each one. Reference in a source or explanatory note.\par
% glossary-block:B00101:end
% glossary-block:B00102:begin
\needspace{3\baselineskip}
\textbf{ǣghwæðer} (indefinite pronoun/adjective). each of two; both. Reference in a source or explanatory note.\par
% glossary-block:B00102:end
% glossary-block:B00103:begin
\needspace{3\baselineskip}
\textbf{ǣghwilc} (indefinite adjective). every; each. each; every.\par
% glossary-block:B00103:end
% glossary-block:B00104:begin
\needspace{3\baselineskip}
\textbf{ǣgðer} (coordinating adverb). both. In ǣgðer ge … ge.\par
% glossary-block:B00104:end
% glossary-block:B00105:begin
\needspace{3\baselineskip}
\textbf{ǣht} (n. f.). possessions; goods.\par
% glossary-block:B00105:end
% glossary-block:B00106:begin
\needspace{3\baselineskip}
\textbf{ǣl} (n. f.); gen. sg. ǣle; nom. pl. ǣla. awl.\par
% glossary-block:B00106:end
% glossary-block:B00107:begin
\needspace{3\baselineskip}
\textbf{ǣlc} (adj.). each; every. With singular noun; substantival each follows the case required by its clause.\par
% glossary-block:B00107:end
% glossary-block:B00108:begin
\needspace{3\baselineskip}
\textbf{Ælfred} (proper name); gen. sg. Ælfredes. Alfred.\par
% glossary-block:B00108:end
% glossary-block:B00109:begin
\needspace{3\baselineskip}
\textbf{ælmihtig} (adj.). almighty.\par
% glossary-block:B00109:end
% glossary-block:B00110:begin
\needspace{3\baselineskip}
\textbf{ǣmetta} (n. m.). leisure; available time.\par
% glossary-block:B00110:end
% glossary-block:B00111:begin
\needspace{3\baselineskip}
\textbf{ǣmtig} (adj.). empty; void.\par
% glossary-block:B00111:end
% glossary-block:B00112:begin
\needspace{3\baselineskip}
\textbf{ǣnig} (indefinite adjective). any.\par
% glossary-block:B00112:end
% glossary-block:B00113:begin
\needspace{3\baselineskip}
\textbf{ǣnne} — acc. masc. sg., variant: see \textbf{ān}.\par
% glossary-block:B00113:end
% glossary-block:B00114:begin
\needspace{3\baselineskip}
\textbf{ænne} — acc. masc. sg., variant: see \textbf{ān}.\par
% glossary-block:B00114:end
% glossary-block:B00115:begin
\needspace{3\baselineskip}
\textbf{æppelbǣre} (adj.). fruit-bearing.\par
% glossary-block:B00115:end
% glossary-block:B00116:begin
\needspace{3\baselineskip}
\textbf{ǣr} (adv.). before; earlier.\par
% glossary-block:B00116:end
% glossary-block:B00117:begin
\needspace{3\baselineskip}
\textbf{ǣr} (conj.). before.\par
% glossary-block:B00117:end
% glossary-block:B00118:begin
\needspace{3\baselineskip}
\textbf{ǣr} (n. n.); gen. sg. ǣres; nom. pl. ǣr. copper.\par
% glossary-block:B00118:end
% glossary-block:B00119:begin
\needspace{3\baselineskip}
\textbf{ǣrendwreca} (n. m.). messengers; envoys.\par
% glossary-block:B00119:end
% glossary-block:B00120:begin
\needspace{3\baselineskip}
\textbf{ǣrest} (adv.). first; at first.\par
% glossary-block:B00120:end
% glossary-block:B00121:begin
\needspace{3\baselineskip}
\textbf{ǣresta} (ordinal numeral). first.\par
% glossary-block:B00121:end
% glossary-block:B00122:begin
\needspace{3\baselineskip}
\textbf{ǣrra} (adj.). former; earlier.\par
% glossary-block:B00122:end
% glossary-block:B00123:begin
\needspace{3\baselineskip}
\textbf{ǣswician} (v. weak); pret. 1 sg. ǣswicode; past ptcp. ǣswicod. offend; cause to stumble. Accusative person; ethical/metaphorical offence in the assigned use.\par
% glossary-block:B00123:end
% glossary-block:B00124:begin
\needspace{3\baselineskip}
\textbf{æt} (prep.). at. from. æt + dative person with learn. from; out of. Proceeds of work, in the assigned expression.\par
% glossary-block:B00124:end
% glossary-block:B00125:begin
\needspace{3\baselineskip}
\textbf{ætberstan} (v.). escape.\par
% glossary-block:B00125:end
% glossary-block:B00126:begin
\needspace{3\baselineskip}
\textbf{ætēowian} (v. weak); pret. 1 sg. ætēowode; past ptcp. ætēowod. appear.\par
% glossary-block:B00126:end
% glossary-block:B00127:begin
\needspace{3\baselineskip}
\textbf{ætforan} (prep.). before; in front of. Dative landmark/person in the supplied frame; retain ætforan rather than replacing it with beforan.\par
% glossary-block:B00127:end
% glossary-block:B00128:begin
\needspace{3\baselineskip}
\textbf{ætgædere} (adv.). together; at once. Paired with samod in the supplied source.\par
% glossary-block:B00128:end
% glossary-block:B00129:begin
\needspace{3\baselineskip}
\textbf{Æthelwulf} (proper name). Æthelwulf.\par
% glossary-block:B00129:end
% glossary-block:B00130:begin
\needspace{3\baselineskip}
\textbf{ætīewan} (v. weak); pret. 1 sg. ætīewde; past ptcp. ætīewed. show; display; reveal.\par
% glossary-block:B00130:end
% glossary-block:B00131:begin
\needspace{3\baselineskip}
\textbf{æðelboren} (adj.). high-born.\par
% glossary-block:B00131:end
% glossary-block:B00132:begin
\needspace{3\baselineskip}
\textbf{æðele} (adj.). noble people. Substantival adjective in the source. noble; illustrious.\par
% glossary-block:B00132:end
% glossary-block:B00133:begin
\needspace{3\baselineskip}
\textbf{Æðelðrȳð} (proper name). Etheldred; Æthelthryth. Proper name.\par
% glossary-block:B00133:end
% glossary-block:B00134:begin
\needspace{3\baselineskip}
\textbf{Æðelwald} (proper name). Æthelwald. Proper name in apposition to father.\par
% glossary-block:B00134:end
% glossary-block:B00135:begin
\needspace{3\baselineskip}
\textbf{æx} (n. f.); gen. sg. æxe; nom. pl. æxa. ax.\par
% glossary-block:B00135:end
\pdfbookmark[1]{B}{outline-B00136}
% glossary-block:B00136:begin
\needspace{5\baselineskip}
\subsection*{B}
% glossary-block:B00136:end
% glossary-block:B00137:begin
\textbf{bā} — nom./acc. fem./neut. pl.: see \textbf{bēgen}.\par
% glossary-block:B00137:end
% glossary-block:B00138:begin
\needspace{3\baselineskip}
\textbf{bām} — dat. pl.: see \textbf{bēgen}.\par
% glossary-block:B00138:end
% glossary-block:B00139:begin
\needspace{3\baselineskip}
\textbf{bān} (n. n.); gen. sg. bānes; nom. pl. bān. bone; ivory. Ivory is elpes bān, literally elephant's bone.\par
% glossary-block:B00139:end
% glossary-block:B00140:begin
\needspace{3\baselineskip}
\textbf{baðian} (v.); pret. 1 sg. baðode. bathe. Supplied source form bathed.\par
% glossary-block:B00140:end
% glossary-block:B00141:begin
\needspace{3\baselineskip}
\textbf{bǣm} — dat. pl., variant: see \textbf{bēgen}.\par
% glossary-block:B00141:end
% glossary-block:B00142:begin
\needspace{3\baselineskip}
\textbf{bæð} (n. n.); gen. sg. bæðes; nom. pl. baðu. bath.\par
% glossary-block:B00142:end
% glossary-block:B00143:begin
\needspace{3\baselineskip}
\textbf{bæðstęde} (n. m.). bathing place; gymnasium. Gymnasium is the contextual interpretation in the supplied reading.\par
% glossary-block:B00143:end
% glossary-block:B00144:begin
\needspace{3\baselineskip}
\textbf{be} (prep.). according to. With dative. concerning; about. With dative. by that; from that. Refers to the preceding proposition; supplied expression. word by word. In word be worde; literal correspondence. by. Body part by which a person is led; supplied grasping phrase. in the south. Only in supplied be sūðan. by that; on that evidence. Evidence/inference frame, not an unrestricted source preposition. by permission. Supplied permission expression.\par
% glossary-block:B00144:end
% glossary-block:B00145:begin
\needspace{3\baselineskip}
\textbf{beald} (adj.). bold; brave.\par
% glossary-block:B00145:end
% glossary-block:B00146:begin
\needspace{3\baselineskip}
\textbf{bearn} (n. n.); gen. sg. bearnes; nom. pl. bearn. son; child.\par
% glossary-block:B00146:end
% glossary-block:B00147:begin
\needspace{3\baselineskip}
\textbf{bēatan} (v. strong); pret. 1 sg. bēot; pret. pl. bēoton; past ptcp. bēaten; pres. 3 sg. bīett. beat; strike. Accusative person or thing struck.\par
% glossary-block:B00147:end
% glossary-block:B00148:begin
\needspace{3\baselineskip}
\textbf{bebēodan} (v. strong); pret. 1 sg. bebēad; pret. pl. bebudon; past ptcp. beboden; pres. 3 sg. bebīett. command; instruct. Dative addressee and a ðæt clause with subjunctive for the commanded action.\par
% glossary-block:B00148:end
% glossary-block:B00149:begin
\needspace{3\baselineskip}
\textbf{bebīett} — pres. ind. 3 sg.: see \textbf{bebēodan}.\par
% glossary-block:B00149:end
% glossary-block:B00150:begin
\needspace{3\baselineskip}
\textbf{bebod} (n. n.); gen. sg. bebodes; nom. pl. bebodu. commandment.\par
% glossary-block:B00150:end
% glossary-block:B00151:begin
\needspace{3\baselineskip}
\textbf{bēc} — gen./dat. sg.; nom./acc. pl.: see \textbf{bōc}.\par
% glossary-block:B00151:end
% glossary-block:B00152:begin
\needspace{3\baselineskip}
\textbf{bēcrǣding} (n. f.); gen. sg. bēcrǣdinge; nom. pl. bēcrǣdinga. reading books.\par
% glossary-block:B00152:end
% glossary-block:B00153:begin
\needspace{3\baselineskip}
\textbf{becuman} (v.). come; arrive. occur; come about. Of the calm in the source. come upon; befall. Retain the source's affected participant.\par
% glossary-block:B00153:end
% glossary-block:B00154:begin
\needspace{3\baselineskip}
\textbf{beēode} — pret. sg.: see \textbf{begān}.\par
% glossary-block:B00154:end
% glossary-block:B00155:begin
\needspace{3\baselineskip}
\textbf{befeallan} (v.); pres. 3 sg. befielð. fall in love. Supplied idiom makes lufu love the subject, its possessor the lover, and on + accusative the beloved.\par
% glossary-block:B00155:end
% glossary-block:B00156:begin
\needspace{3\baselineskip}
\textbf{befielð} — pres. ind. 3 sg.: see \textbf{befeallan}.\par
% glossary-block:B00156:end
% glossary-block:B00157:begin
\needspace{3\baselineskip}
\textbf{be-foran} (prep.). before; in front of. Dative landmark in the static frame. before; ahead of. Dative person, including going or sending ahead of that person. before; in front of. Accusative landmark in the positioning construal; static dative can be a different valid construal.\par
% glossary-block:B00157:end
% glossary-block:B00158:begin
\needspace{3\baselineskip}
\textbf{begān} (v. irregular); pres. 1 sg. begā; pres. 3 sg. begǣð; pres. pl. begāð; pret. 1 sg. beēode; pret. pl. beēodon; past ptcp. begān. ply; practise. Accusative trade or occupation. pursue; practise. Study is the thing pursued.\par
% glossary-block:B00158:end
% glossary-block:B00159:begin
\needspace{3\baselineskip}
\textbf{begǣð} — pres. ind. 3 sg.: see \textbf{begān}.\par
% glossary-block:B00159:end
% glossary-block:B00160:begin
\needspace{3\baselineskip}
\textbf{bēgen} (numeral). both.\par
% glossary-block:B00160:end
% glossary-block:B00161:begin
\needspace{3\baselineskip}
\textbf{begyrdan} (v.); past ptcp. begyrd. gird; clothe. Supplied participial clothing expression.\par
% glossary-block:B00161:end
% glossary-block:B00162:begin
\needspace{3\baselineskip}
\textbf{behātan} (v.); pret. 1 sg. behēt. promise.\par
% glossary-block:B00162:end
% glossary-block:B00163:begin
\needspace{3\baselineskip}
\textbf{behēafdian} (v. weak); pret. 1 sg. behēafdode; past ptcp. behēafdod. behead. Accusative person beheaded.\par
% glossary-block:B00163:end
% glossary-block:B00164:begin
\needspace{3\baselineskip}
\textbf{behealdan} (v. strong); pret. 1 sg. behēold; pret. pl. behēoldon; past ptcp. behealden; pres. 3 sg. behielt. behold; look at. Accusative object seen. observe. Supplied source form observed.\par
% glossary-block:B00164:end
% glossary-block:B00165:begin
\needspace{3\baselineskip}
\textbf{behēfe} (adj.). useful. Dative person to whom something is useful.\par
% glossary-block:B00165:end
% glossary-block:B00166:begin
\needspace{3\baselineskip}
\textbf{behielt} — pres. ind. 3 sg.: see \textbf{behealdan}.\par
% glossary-block:B00166:end
% glossary-block:B00167:begin
\needspace{3\baselineskip}
\textbf{behygdig} (adj.). attentive; thoughtful.\par
% glossary-block:B00167:end
% glossary-block:B00168:begin
\needspace{3\baselineskip}
\textbf{bęlg} (n. m.); gen. sg. bęlges; nom. pl. bęlgas. bellows. Selected plural instrument use.\par
% glossary-block:B00168:end
% glossary-block:B00169:begin
\needspace{3\baselineskip}
\textbf{belimpan} (v.). belong; relate. Supplied participial construction.\par
% glossary-block:B00169:end
% glossary-block:B00170:begin
\needspace{3\baselineskip}
\textbf{bēn} (n. f.). prayer; petition. Source dative plural after tō. petition; appeal.\par
% glossary-block:B00170:end
% glossary-block:B00171:begin
\needspace{3\baselineskip}
\textbf{Beniamin} (proper name); gen. sg. Beniamines. Benjamin.\par
% glossary-block:B00171:end
% glossary-block:B00172:begin
\needspace{3\baselineskip}
\textbf{bēodan} (v. strong); pret. 1 sg. bēad; pret. pl. budon; past ptcp. boden; pres. 3 sg. bīett. offer. command; bid.\par
% glossary-block:B00172:end
% glossary-block:B00173:begin
\needspace{3\baselineskip}
\textbf{bēon} — alternative infinitive: see \textbf{wesan}.\par
% glossary-block:B00173:end
% glossary-block:B00174:begin
\needspace{3\baselineskip}
\textbf{beorhtnes} (n. f.); gen. sg. beorhtnesse. brightness.\par
% glossary-block:B00174:end
% glossary-block:B00175:begin
\needspace{3\baselineskip}
\textbf{beorn} (n. m.); gen. sg. beornes; nom. pl. beornas. warrior.\par
% glossary-block:B00175:end
% glossary-block:B00176:begin
\needspace{3\baselineskip}
\textbf{bēorscipe} (n. m.); gen. sg. bēorscipes. feast; banquet.\par
% glossary-block:B00176:end
% glossary-block:B00177:begin
\needspace{3\baselineskip}
\textbf{beran} (v. strong); pret. 1 sg. bær; pret. pl. bǣron; past ptcp. boren; pres. 3 sg. birð. carry. Accusative person or thing carried. bear; produce. Accusative produce; distinct from physical carrying. rich in; productive of. Supplied berende on + dative ore expression; participial adjective.\par
% glossary-block:B00177:end
% glossary-block:B00178:begin
\needspace{3\baselineskip}
\textbf{berēafian} (v.); past ptcp. berēafod. despoil; deprive of. Accusative person deprived, genitive thing lost; passive retains that genitive.\par
% glossary-block:B00178:end
% glossary-block:B00179:begin
\needspace{3\baselineskip}
\textbf{berstan} (v. strong); pret. 1 sg. bærst; pret. pl. burston; past ptcp. borsten; pres. 3 sg. bierst. burst.\par
% glossary-block:B00179:end
% glossary-block:B00180:begin
\needspace{3\baselineskip}
\textbf{besēon} (v.). look; turn one's gaze. Supplied source form and frame. look toward. Target of the look follows the supplied phrase.\par
% glossary-block:B00180:end
% glossary-block:B00181:begin
\needspace{3\baselineskip}
\textbf{bęt} — comparative adverb: see \textbf{wel}.\par
% glossary-block:B00181:end
% glossary-block:B00182:begin
\needspace{3\baselineskip}
\textbf{bętera} (adj.). better. Comparative of gōd.\par
% glossary-block:B00182:end
% glossary-block:B00183:begin
\needspace{3\baselineskip}
\textbf{bętst} (adj.). best. Superlative of gōd. Also superlative adverb: see \textbf{wel}.\par
% glossary-block:B00183:end
% glossary-block:B00184:begin
\needspace{3\baselineskip}
\textbf{betwēonan} (adverb or particle). between; to one another. Supplied reciprocal betwēonan him in 45.6.\par
% glossary-block:B00184:end
% glossary-block:B00185:begin
\needspace{3\baselineskip}
\textbf{betweox} (prep.). among. Accusative in the exact supplied phrases; not a universal case rule. among; between.\par
% glossary-block:B00185:end
% glossary-block:B00186:begin
\needspace{3\baselineskip}
\textbf{beðearf} — pres. ind. 1/3 sg.: see \textbf{beðurfan}.\par
% glossary-block:B00186:end
% glossary-block:B00187:begin
\needspace{3\baselineskip}
\textbf{beðorfte} — pret. sg.: see \textbf{beðurfan}.\par
% glossary-block:B00187:end
% glossary-block:B00188:begin
\needspace{3\baselineskip}
\textbf{beðurfan} (v. preterite-present); pres. 1/3 sg. beðearf; pres. pl. beðurfon; pret. 1 sg. beðorfte. need; require. Genitive thing needed. No past-participle citation is given in the checked source paradigm.\par
% glossary-block:B00188:end
% glossary-block:B00189:begin
\needspace{3\baselineskip}
\textbf{bewǣfan} (v.); past ptcp. bewǣfed. wrap; clothe. Supplied participial clothing expression.\par
% glossary-block:B00189:end
% glossary-block:B00190:begin
\needspace{3\baselineskip}
\textbf{bewęndan} (v. weak); pret. 1 sg. bewęnde; past ptcp. bewęnd. turn. The assigned reflexive frame uses dative him and tō + dative goal.\par
% glossary-block:B00190:end
% glossary-block:B00191:begin
\needspace{3\baselineskip}
\textbf{biddan} (v. strong); pret. 1 sg. bæd; pret. pl. bǣdon; past ptcp. beden. ask; request. Genitive thing requested; an accusative person may be added. beseech; entreat. Accusative person; the requested action can be a ðæt clause. wish well; bid farewell. Supplied farewell formula.\par
% glossary-block:B00191:end
% glossary-block:B00192:begin
\needspace{3\baselineskip}
\textbf{bīegan} (v. weak); pret. 1 sg. bīegde; past ptcp. bīeged. bend. Supplied past bent; retain the source's unmarked ie in its actual form. bend.\par
% glossary-block:B00192:end
% glossary-block:B00193:begin
\needspace{3\baselineskip}
\textbf{biema} (n. m.); gen. sg. bieman; nom. pl. bieman. trumpet; clarion.\par
% glossary-block:B00193:end
% glossary-block:B00194:begin
\needspace{3\baselineskip}
\textbf{bierst} — pres. ind. 2 sg., variant: see \textbf{beran}; pres. ind. 3 sg.: see \textbf{berstan}.\par
% glossary-block:B00194:end
% glossary-block:B00195:begin
\needspace{3\baselineskip}
\textbf{bierð} — pres. ind. 3 sg., variant: see \textbf{beran}.\par
% glossary-block:B00195:end
% glossary-block:B00196:begin
\needspace{3\baselineskip}
\textbf{bīet} — pres. ind. 3 sg., variant: see \textbf{bēodan}.\par
% glossary-block:B00196:end
% glossary-block:B00197:begin
\needspace{3\baselineskip}
\textbf{bīett} — pres. ind. 3 sg.: see \textbf{bēatan}; pres. ind. 3 sg.: see \textbf{bēodan}.\par
% glossary-block:B00197:end
% glossary-block:B00198:begin
\needspace{3\baselineskip}
\textbf{bindan} (v. strong); pret. 1 sg. band; pret. pl. bundon; past ptcp. bunden. bind; tie. Accusative person or thing bound.\par
% glossary-block:B00198:end
% glossary-block:B00199:begin
\needspace{3\baselineskip}
\textbf{binnan} (prep.). within.\par
% glossary-block:B00199:end
% glossary-block:B00200:begin
\needspace{3\baselineskip}
\textbf{birð} — pres. ind. 3 sg.: see \textbf{beran}.\par
% glossary-block:B00200:end
% glossary-block:B00201:begin
\needspace{3\baselineskip}
\textbf{bisceop} (n. m.); gen. sg. bisceopes; nom. pl. bisceopas. bishop.\par
% glossary-block:B00201:end
% glossary-block:B00202:begin
\needspace{3\baselineskip}
\textbf{bisgian} (v.). occupy; trouble.\par
% glossary-block:B00202:end
% glossary-block:B00203:begin
\needspace{3\baselineskip}
\textbf{bisgu} (n. f.). occupations; cares.\par
% glossary-block:B00203:end
% glossary-block:B00204:begin
\needspace{3\baselineskip}
\textbf{bist} — pres. ind. 2 sg.: see \textbf{wesan}.\par
% glossary-block:B00204:end
% glossary-block:B00205:begin
\needspace{3\baselineskip}
\textbf{bið} — pres. ind. 3 sg.: see \textbf{wesan}.\par
% glossary-block:B00205:end
% glossary-block:B00206:begin
\needspace{3\baselineskip}
\textbf{blāwan} (v. strong); pret. 1 sg. blēow; pret. pl. blēowon; past ptcp. blāwen; pres. 3 sg. blǣwð. blow. Intransitive in the assigned bellows sentence.\par
% glossary-block:B00206:end
% glossary-block:B00207:begin
\needspace{3\baselineskip}
\textbf{blæc} (adj.). black.\par
% glossary-block:B00207:end
% glossary-block:B00208:begin
\needspace{3\baselineskip}
\textbf{blǣwð} — pres. ind. 3 sg.: see \textbf{blāwan}.\par
% glossary-block:B00208:end
% glossary-block:B00209:begin
\needspace{3\baselineskip}
\textbf{blētsian} (v.); past ptcp. blētsod. bless. Supplied third-singular preterite; God is the understood subject. bless.\par
% glossary-block:B00209:end
% glossary-block:B00210:begin
\needspace{3\baselineskip}
\textbf{blētsung} (n. f.); gen. sg. blētsunge. blessing.\par
% glossary-block:B00210:end
% glossary-block:B00211:begin
\needspace{3\baselineskip}
\textbf{blinnan} (v.). cease. Supplied wæs blinnende periphrasis.\par
% glossary-block:B00211:end
% glossary-block:B00212:begin
\needspace{3\baselineskip}
\textbf{blissesang} (n. m.). song of joy.\par
% glossary-block:B00212:end
% glossary-block:B00213:begin
\needspace{3\baselineskip}
\textbf{blissian} (v. weak); pret. 1 sg. blissode; past ptcp. blissod. rejoice. Source genitive plural present-participle use, of people rejoicing. rejoice. Intransitive; mid + dative companion or ofer + accusative cause of joy in the assigned frames.\par
% glossary-block:B00213:end
% glossary-block:B00214:begin
\needspace{3\baselineskip}
\textbf{blīðe} (adj.). merry; glad.\par
% glossary-block:B00214:end
% glossary-block:B00215:begin
\needspace{3\baselineskip}
\textbf{blīðe} (adv.). joyfully; gladly.\par
% glossary-block:B00215:end
% glossary-block:B00216:begin
\needspace{3\baselineskip}
\textbf{blōstma} (n. m.); gen. sg. blōstman; nom. pl. blōstman. blossom.\par
% glossary-block:B00216:end
% glossary-block:B00217:begin
\needspace{3\baselineskip}
\textbf{blōstmian} (v. weak); pret. 1 sg. blōstmode; past ptcp. blōstmod. blossom; bloom.\par
% glossary-block:B00217:end
% glossary-block:B00218:begin
\needspace{3\baselineskip}
\textbf{bōc} (n. f.); gen. sg. bēc; nom. pl. bēc. book. Genitive singular variant bōce.\par
% glossary-block:B00218:end
% glossary-block:B00219:begin
\needspace{3\baselineskip}
\textbf{Bōclæden} (n. n.). book-Latin; Latin.\par
% glossary-block:B00219:end
% glossary-block:B00220:begin
\needspace{3\baselineskip}
\textbf{bohte} — pret. sg.: see \textbf{bycgan}.\par
% glossary-block:B00220:end
% glossary-block:B00221:begin
\needspace{3\baselineskip}
\textbf{brād} (adj.). broad.\par
% glossary-block:B00221:end
% glossary-block:B00222:begin
\needspace{3\baselineskip}
\textbf{brādnes} (n. f.); gen. sg. brādnesse; nom. pl. brādnessa. surface; face. Of the deep in the assigned expression.\par
% glossary-block:B00222:end
% glossary-block:B00223:begin
\needspace{3\baselineskip}
\textbf{brǣd} — pret. sg., contracted variant: see \textbf{bregdan}.\par
% glossary-block:B00223:end
% glossary-block:B00224:begin
\needspace{3\baselineskip}
\textbf{brǣdan} (v.); pres. 3 sg. brǣt. spread; extend. Supplied source form spreads.\par
% glossary-block:B00224:end
% glossary-block:B00225:begin
\needspace{3\baselineskip}
\textbf{bregdan} (v. strong); pret. 1 sg. brægd; pret. pl. brugdon; past ptcp. brogden; pres. 3 sg. brītt. brandish; draw.\par
% glossary-block:B00225:end
% glossary-block:B00226:begin
\needspace{3\baselineskip}
\textbf{Breoton} (proper name). Britain.\par
% glossary-block:B00226:end
% glossary-block:B00227:begin
\needspace{3\baselineskip}
\textbf{brēðer} — dat. sg.: see \textbf{brōðor}.\par
% glossary-block:B00227:end
% glossary-block:B00228:begin
\needspace{3\baselineskip}
\textbf{bringan} (v. weak); pret. 1 sg. brōhte; past ptcp. (ge)brōht. bring. Accusative thing or person brought; tō + dative destination/person, or the assigned bare dative recipient.\par
% glossary-block:B00228:end
% glossary-block:B00229:begin
\needspace{3\baselineskip}
\textbf{brīt} — pres. ind. 3 sg., variant: see \textbf{bregdan}.\par
% glossary-block:B00229:end
% glossary-block:B00230:begin
\needspace{3\baselineskip}
\textbf{brītt} — pres. ind. 3 sg.: see \textbf{bregdan}.\par
% glossary-block:B00230:end
% glossary-block:B00231:begin
\needspace{3\baselineskip}
\textbf{brōden} — past ptcp., contracted variant: see \textbf{bregdan}.\par
% glossary-block:B00231:end
% glossary-block:B00232:begin
\needspace{3\baselineskip}
\textbf{brōhte} — pret. sg.: see \textbf{bringan}.\par
% glossary-block:B00232:end
% glossary-block:B00233:begin
\needspace{3\baselineskip}
\textbf{brōðor} (n. m.); gen. sg. brōðor; nom. pl. brōðor. brother. Dative singular brēðer.\par
% glossary-block:B00233:end
% glossary-block:B00234:begin
\needspace{3\baselineskip}
\textbf{brycg} (n. f.); gen. sg. brycge; nom. pl. brycga. bridge.\par
% glossary-block:B00234:end
% glossary-block:B00235:begin
\needspace{3\baselineskip}
\textbf{brȳd} (n. f.); gen. sg. brȳde; nom. pl. brȳda. bride.\par
% glossary-block:B00235:end
% glossary-block:B00236:begin
\needspace{3\baselineskip}
\textbf{Bryttas} (ethnonym). Britons.\par
% glossary-block:B00236:end
% glossary-block:B00237:begin
\needspace{3\baselineskip}
\textbf{bū} — nom./acc. neut. pl.: see \textbf{bēgen}.\par
% glossary-block:B00237:end
% glossary-block:B00238:begin
\needspace{3\baselineskip}
\textbf{bufan} (prep.). above. With dative landmark in the taught spatial frame.\par
% glossary-block:B00238:end
% glossary-block:B00239:begin
\needspace{3\baselineskip}
\textbf{būgan} (v. strong); pret. 1 sg. bēag; pret. pl. bugon; past ptcp. bogen; pres. 3 sg. bȳgð. bend.\par
% glossary-block:B00239:end
% glossary-block:B00240:begin
\needspace{3\baselineskip}
\textbf{būr} (n. m.). chamber; room.\par
% glossary-block:B00240:end
% glossary-block:B00241:begin
\needspace{3\baselineskip}
\textbf{burg} (n. f.); gen. sg. byrig; nom. pl. byrig. town; city. Genitive singular variant burge.\par
% glossary-block:B00241:end
% glossary-block:B00242:begin
\needspace{3\baselineskip}
\textbf{būtan} (prep.). besides; apart from. Dative excluded person or thing. without. With dative. outside.\par
% glossary-block:B00242:end
% glossary-block:B00243:begin
\needspace{3\baselineskip}
\textbf{būtan} (conj.). except.\par
% glossary-block:B00243:end
% glossary-block:B00244:begin
\needspace{3\baselineskip}
\textbf{būtan} (conj.). unless; except if.\par
% glossary-block:B00244:end
% glossary-block:B00245:begin
\needspace{3\baselineskip}
\textbf{bycgan} (v. weak); pret. 1 sg. bohte; past ptcp. (ge)boht. buy. Accusative goods bought.\par
% glossary-block:B00245:end
% glossary-block:B00246:begin
\needspace{3\baselineskip}
\textbf{bȳgð} — pres. ind. 3 sg.: see \textbf{būgan}.\par
% glossary-block:B00246:end
% glossary-block:B00247:begin
\needspace{3\baselineskip}
\textbf{bȳhð} — pres. ind. 3 sg., variant: see \textbf{būgan}.\par
% glossary-block:B00247:end
% glossary-block:B00248:begin
\needspace{3\baselineskip}
\textbf{byrig} — gen./dat. sg.; nom./acc. pl.: see \textbf{burg}.\par
% glossary-block:B00248:end
% glossary-block:B00249:begin
\needspace{3\baselineskip}
\textbf{bȳsen} (n. f.). example; illustration.\par
% glossary-block:B00249:end
\pdfbookmark[1]{C}{outline-B00250}
% glossary-block:B00250:begin
\needspace{5\baselineskip}
\subsection*{C}
% glossary-block:B00250:end
% glossary-block:B00251:begin
\textbf{campian} (v.). fight.\par
% glossary-block:B00251:end
% glossary-block:B00252:begin
\needspace{3\baselineskip}
\textbf{cann} — pres. ind. 1/3 sg.: see \textbf{cunnan}.\par
% glossary-block:B00252:end
% glossary-block:B00253:begin
\needspace{3\baselineskip}
\textbf{ceaf} (n. n.). chaff.\par
% glossary-block:B00253:end
% glossary-block:B00254:begin
\needspace{3\baselineskip}
\textbf{cēapian} (v. weak); pret. 1 sg. cēapode; past ptcp. cēapod. trade; traffic. Traffic means trade; retain the supplied trading frame.\par
% glossary-block:B00254:end
% glossary-block:B00255:begin
\needspace{3\baselineskip}
\textbf{cēapmann} — citation/spelling variant: see \textbf{cīepmann}.\par
% glossary-block:B00255:end
% glossary-block:B00256:begin
\needspace{3\baselineskip}
\textbf{ceaster} (n. f.); gen. sg. ceastre; nom. pl. ceastra. town; city.\par
% glossary-block:B00256:end
% glossary-block:B00257:begin
\needspace{3\baselineskip}
\textbf{ceasterwaran} (n. m.); nom. pl. ceasterwaran. townspeople; citizens.\par
% glossary-block:B00257:end
% glossary-block:B00258:begin
\needspace{3\baselineskip}
\textbf{cēderbēam} (n. m.); gen. sg. cēderbēames; nom. pl. cēderbēamas. cedar; cedar tree.\par
% glossary-block:B00258:end
% glossary-block:B00259:begin
\needspace{3\baselineskip}
\textbf{cēlan} (v.). cool. Pronunciation example.\par
% glossary-block:B00259:end
% glossary-block:B00260:begin
\needspace{3\baselineskip}
\textbf{cęmpa} (n. m.); gen. sg. cęmpan; nom. pl. cęmpan. soldier.\par
% glossary-block:B00260:end
% glossary-block:B00261:begin
\needspace{3\baselineskip}
\textbf{ceorl} (n. m.); gen. sg. ceorles; nom. pl. ceorlas. layman.\par
% glossary-block:B00261:end
% glossary-block:B00262:begin
\needspace{3\baselineskip}
\textbf{cēosan} (v. strong); pret. 1 sg. cēas; pret. pl. curon; past ptcp. coren; pres. 3 sg. cīest. choose. Accusative person or thing chosen.\par
% glossary-block:B00262:end
% glossary-block:B00263:begin
\needspace{3\baselineskip}
\textbf{cīepa} (n.). merchant. Reference in a source or explanatory note.\par
% glossary-block:B00263:end
% glossary-block:B00264:begin
\needspace{3\baselineskip}
\textbf{cīepan} (v. weak); pret. 1 sg. cīepte; past ptcp. cīeped. sell. Accusative goods sold.\par
% glossary-block:B00264:end
% glossary-block:B00265:begin
\needspace{3\baselineskip}
\textbf{cīepmann} (n.). merchant. Reference in a source or explanatory note.\par
% glossary-block:B00265:end
% glossary-block:B00266:begin
\needspace{3\baselineskip}
\textbf{cīest} — pres. ind. 3 sg.: see \textbf{cēosan}.\par
% glossary-block:B00266:end
% glossary-block:B00267:begin
\needspace{3\baselineskip}
\textbf{cirice} (n. f.); gen. sg. cirican; nom. pl. cirican. church. oratory. Chad's separate place of prayer; distinct from the ordinary church sense.\par
% glossary-block:B00267:end
% glossary-block:B00268:begin
\needspace{3\baselineskip}
\textbf{clǣne} (adj.). pure; clean. clean; pure; blameless.\par
% glossary-block:B00268:end
% glossary-block:B00269:begin
\needspace{3\baselineskip}
\textbf{cleopian} (v. weak); pret. 1 sg. cleopode; past ptcp. cleopod. call; summon. Accusative person summoned; distinguish intransitive cry aloud. cry aloud. Intransitive, with the words cried given as direct speech.\par
% glossary-block:B00269:end
% glossary-block:B00270:begin
\needspace{3\baselineskip}
\textbf{clif} (n. n.). cliff.\par
% glossary-block:B00270:end
% glossary-block:B00271:begin
\needspace{3\baselineskip}
\textbf{cnapa} (n. m.); gen. sg. cnapan; nom. pl. cnapan. boy.\par
% glossary-block:B00271:end
% glossary-block:B00272:begin
\needspace{3\baselineskip}
\textbf{cnāwan} (v. strong); pret. 1 sg. cnēow; pret. pl. cnēowon; past ptcp. cnāwen; pres. 3 sg. cnǣwð. know.\par
% glossary-block:B00272:end
% glossary-block:B00273:begin
\needspace{3\baselineskip}
\textbf{cnǣwð} — pres. ind. 3 sg.: see \textbf{cnāwan}.\par
% glossary-block:B00273:end
% glossary-block:B00274:begin
\needspace{3\baselineskip}
\textbf{cnēo} (n. n.); gen. sg. cnēowes; nom. pl. cnēowu. knee. Supplied accusative plural knees. knee.\par
% glossary-block:B00274:end
% glossary-block:B00275:begin
\needspace{3\baselineskip}
\textbf{cnēoris} (n. f.); gen. sg. cnēorisse; nom. pl. cnēorissa. generation.\par
% glossary-block:B00275:end
% glossary-block:B00276:begin
\needspace{3\baselineskip}
\textbf{cnīf} (n.). knife. Pronunciation example.\par
% glossary-block:B00276:end
% glossary-block:B00277:begin
\needspace{3\baselineskip}
\textbf{cniht} (n. m.); gen. sg. cnihtes; nom. pl. cnihtas. young man; youth. servant.\par
% glossary-block:B00277:end
% glossary-block:B00278:begin
\needspace{3\baselineskip}
\textbf{cnucian} (v.). knock.\par
% glossary-block:B00278:end
% glossary-block:B00279:begin
\needspace{3\baselineskip}
\textbf{coss} (n. m.); gen. sg. cosses; nom. pl. cossas. kiss.\par
% glossary-block:B00279:end
% glossary-block:B00280:begin
\needspace{3\baselineskip}
\textbf{costnung} (n. f.); gen. sg. costnunge; nom. pl. costnunga. temptation.\par
% glossary-block:B00280:end
% glossary-block:B00281:begin
\needspace{3\baselineskip}
\textbf{cræft} (n. m.); gen. sg. cræftes; nom. pl. cræftas. trade; occupation; craft. skill; accomplishment.\par
% glossary-block:B00281:end
% glossary-block:B00282:begin
\needspace{3\baselineskip}
\textbf{Crīst} (n. m.); gen. sg. Crīstes. Christ.\par
% glossary-block:B00282:end
% glossary-block:B00283:begin
\needspace{3\baselineskip}
\textbf{cucu} (adj.). alive; living. Supplied accusative neuter plural adjective.\par
% glossary-block:B00283:end
% glossary-block:B00284:begin
\needspace{3\baselineskip}
\textbf{cuma} (n. m.); gen. sg. cuman; nom. pl. cuman. stranger; guest.\par
% glossary-block:B00284:end
% glossary-block:B00285:begin
\needspace{3\baselineskip}
\textbf{cuman} (v. strong); pret. 1 sg. cōm; pret. pl. cōmon; past ptcp. cumen; pres. 3 sg. cymð. come.\par
% glossary-block:B00285:end
% glossary-block:B00286:begin
\needspace{3\baselineskip}
\textbf{cunnan} (v. preterite-present); pres. 1/3 sg. cann; pres. pl. cunnon; pret. 1 sg. cūðe; past ptcp. (ge)cunnen. know how; be able. Learned skill; plain infinitive complement. know; be acquainted with. Accusative person known; distinct from learned ability.\par
% glossary-block:B00286:end
% glossary-block:B00287:begin
\needspace{3\baselineskip}
\textbf{cūðe} — pret. sg.: see \textbf{cunnan}.\par
% glossary-block:B00287:end
% glossary-block:B00288:begin
\needspace{3\baselineskip}
\textbf{cweartern} (n. n.); gen. sg. cwearternes; nom. pl. cweartern. prison.\par
% glossary-block:B00288:end
% glossary-block:B00289:begin
\needspace{3\baselineskip}
\textbf{cwellan} (v.). kill.\par
% glossary-block:B00289:end
% glossary-block:B00290:begin
\needspace{3\baselineskip}
\textbf{cwēn} (n. f.); gen. sg. cwēne; nom. pl. cwēna. queen. princess.\par
% glossary-block:B00290:end
% glossary-block:B00291:begin
\needspace{3\baselineskip}
\textbf{cweðan} (v. strong); pret. 1 sg. cwæð; pret. pl. cwǣdon; past ptcp. cweden; pres. 3 sg. cwið. say. say. Source genitive plural present participle, of people saying. say. Supplied preterite said.\par
% glossary-block:B00291:end
% glossary-block:B00292:begin
\needspace{3\baselineskip}
\textbf{cwic} (adj.). alive; living. living people. Substantival adjective.\par
% glossary-block:B00292:end
% glossary-block:B00293:begin
\needspace{3\baselineskip}
\textbf{cwið} — pres. ind. 3 sg.: see \textbf{cweðan}.\par
% glossary-block:B00293:end
% glossary-block:B00294:begin
\needspace{3\baselineskip}
\textbf{cwōm} — pret. sg., variant: see \textbf{cuman}.\par
% glossary-block:B00294:end
% glossary-block:B00295:begin
\needspace{3\baselineskip}
\textbf{cymð} — pres. ind. 3 sg.: see \textbf{cuman}.\par
% glossary-block:B00295:end
% glossary-block:B00296:begin
\needspace{3\baselineskip}
\textbf{cynehelm} (n. m.); gen. sg. cynehelmes; nom. pl. cynehelmas. crown.\par
% glossary-block:B00296:end
% glossary-block:B00297:begin
\needspace{3\baselineskip}
\textbf{cynesętl} (n. n.). royal seat; throne.\par
% glossary-block:B00297:end
% glossary-block:B00298:begin
\needspace{3\baselineskip}
\textbf{cyning} (n. m.); gen. sg. cyninges; nom. pl. cyningas. king.\par
% glossary-block:B00298:end
% glossary-block:B00299:begin
\needspace{3\baselineskip}
\textbf{cynn} (n. n.); gen. sg. cynnes; nom. pl. cynn. kind.\par
% glossary-block:B00299:end
% glossary-block:B00300:begin
\needspace{3\baselineskip}
\textbf{Cyrenisc} (adj.). Cyrenian; of Cyrene. Geographical adjective.\par
% glossary-block:B00300:end
% glossary-block:B00301:begin
\needspace{3\baselineskip}
\textbf{cyrice} — citation/spelling variant: see \textbf{cirice}.\par
% glossary-block:B00301:end
% glossary-block:B00302:begin
\needspace{3\baselineskip}
\textbf{cyssan} (v. weak); pret. 1 sg. cyste; past ptcp. (ge)cyssed. kiss.\par
% glossary-block:B00302:end
% glossary-block:B00303:begin
\needspace{3\baselineskip}
\textbf{cȳðan} (v. weak); pret. 1 sg. cȳðde; past ptcp. cȳðed. announce; make known. Dative person informed; accusative thing announced or a ðæt content clause.\par
% glossary-block:B00303:end
\pdfbookmark[1]{D}{outline-B00304}
% glossary-block:B00304:begin
\needspace{5\baselineskip}
\subsection*{D}
% glossary-block:B00304:end
% glossary-block:B00305:begin
\textbf{dǣd} (n. f.). deed.\par
% glossary-block:B00305:end
% glossary-block:B00306:begin
\needspace{3\baselineskip}
\textbf{dæg} (n. m.); gen. sg. dæges; nom. pl. dagas. day.\par
% glossary-block:B00306:end
% glossary-block:B00307:begin
\needspace{3\baselineskip}
\textbf{dæghwāmlīc} (adj.). daily.\par
% glossary-block:B00307:end
% glossary-block:B00308:begin
\needspace{3\baselineskip}
\textbf{dǣl} (n. m.); gen. sg. dǣles; nom. pl. dǣlas. part.\par
% glossary-block:B00308:end
% glossary-block:B00309:begin
\needspace{3\baselineskip}
\textbf{dēad} (adj.). dead. dead people. Substantival adjective.\par
% glossary-block:B00309:end
% glossary-block:B00310:begin
\needspace{3\baselineskip}
\textbf{dearr} — pres. ind. 1/3 sg.: see \textbf{durran}.\par
% glossary-block:B00310:end
% glossary-block:B00311:begin
\needspace{3\baselineskip}
\textbf{dēað} (n. m.); gen. sg. dēaðes; nom. pl. dēaðas. death.\par
% glossary-block:B00311:end
% glossary-block:B00312:begin
\needspace{3\baselineskip}
\textbf{dehter} — dat. sg.: see \textbf{dohtor}.\par
% glossary-block:B00312:end
% glossary-block:B00313:begin
\needspace{3\baselineskip}
\textbf{dēma} (n. m.); gen. sg. dēman; nom. pl. dēman. judge.\par
% glossary-block:B00313:end
% glossary-block:B00314:begin
\needspace{3\baselineskip}
\textbf{dēman} (v. weak); pret. 1 sg. dēmde; past ptcp. dēmed. judge. condemn. Accusative person condemned.\par
% glossary-block:B00314:end
% glossary-block:B00315:begin
\needspace{3\baselineskip}
\textbf{Dęne} (n. m.); nom. pl. Dęne. Danes. Plural people-name.\par
% glossary-block:B00315:end
% glossary-block:B00316:begin
\needspace{3\baselineskip}
\textbf{dēofol} (n. m.); gen. sg. dēofles; nom. pl. dēoflas. devil; demon.\par
% glossary-block:B00316:end
% glossary-block:B00317:begin
\needspace{3\baselineskip}
\textbf{dēor} (n. n.); gen. sg. dēores; nom. pl. dēor. animal; beast.\par
% glossary-block:B00317:end
% glossary-block:B00318:begin
\needspace{3\baselineskip}
\textbf{dēorwierðe} (adj.); comp. dēorwierðra. precious; valuable.\par
% glossary-block:B00318:end
% glossary-block:B00319:begin
\needspace{3\baselineskip}
\textbf{dēorwierðra} (adj.). more precious. Comparative adjective; distinct from an adverb modifying an action. Comparative of dēorwierðe.\par
% glossary-block:B00319:end
% glossary-block:B00320:begin
\needspace{3\baselineskip}
\textbf{dēorwurð} (adj.). precious; costly.\par
% glossary-block:B00320:end
% glossary-block:B00321:begin
\needspace{3\baselineskip}
\textbf{dēð} — pres. ind. 3 sg.: see \textbf{dōn}.\par
% glossary-block:B00321:end
% glossary-block:B00322:begin
\needspace{3\baselineskip}
\textbf{dīegelnes} (n. f.). retreat; seclusion. Supplied plural after of.\par
% glossary-block:B00322:end
% glossary-block:B00323:begin
\needspace{3\baselineskip}
\textbf{diernan} (v.). conceal; hide.\par
% glossary-block:B00323:end
% glossary-block:B00324:begin
\needspace{3\baselineskip}
\textbf{disc} (n. m.); gen. sg. disces; nom. pl. discas. dish.\par
% glossary-block:B00324:end
% glossary-block:B00325:begin
\needspace{3\baselineskip}
\textbf{docga} (n.). dog. Pronunciation example.\par
% glossary-block:B00325:end
% glossary-block:B00326:begin
\needspace{3\baselineskip}
\textbf{dohtor} (n. f.); gen. sg. dohtor; nom. pl. dohtor. daughter. Dative singular dehter, also a late genitive variant.\par
% glossary-block:B00326:end
% glossary-block:B00327:begin
\needspace{3\baselineskip}
\textbf{dōm} (n. m.); gen. sg. dōmes; nom. pl. dōmas. judgment. dōmes dæg means Day of Judgment. doom; fate. judgment. A judgment made in judging, distinct from doom/fate.\par
% glossary-block:B00327:end
% glossary-block:B00328:begin
\needspace{3\baselineskip}
\textbf{dōn} (v. irregular); pres. 1 sg. dō; pres. 3 sg. dēð; pres. pl. dōð; pret. 1 sg. dyde; pret. pl. dydon; past ptcp. (ge)dōn. do; perform. do for; do to. Dative beneficiary/affected person; thing done is accusative when expressed. put; place. Accusative object; on + accusative goal in the assigned frame. make; render. Accusative object with agreeing object predicate. do. Replaces an action already mentioned; distinct from English dummy do. cause; bring about. With a finite ðæt clause; its actor is nominative within that clause. fulfil; carry out. A word or command is fulfilled. deal with; behave toward. Assigned wið person frame; manner is separately expressed.\par
% glossary-block:B00328:end
% glossary-block:B00329:begin
\needspace{3\baselineskip}
\textbf{dorste} — pret. sg.: see \textbf{durran}.\par
% glossary-block:B00329:end
% glossary-block:B00330:begin
\needspace{3\baselineskip}
\textbf{draca} (n. m.); gen. sg. dracan; nom. pl. dracan. dragon.\par
% glossary-block:B00330:end
% glossary-block:B00331:begin
\needspace{3\baselineskip}
\textbf{drēam} (n. m.). joy; bliss.\par
% glossary-block:B00331:end
% glossary-block:B00332:begin
\needspace{3\baselineskip}
\textbf{dręnc} (n. m.); gen. sg. dręnces; nom. pl. dręncas. drink.\par
% glossary-block:B00332:end
% glossary-block:B00333:begin
\needspace{3\baselineskip}
\textbf{drēogan} (v. strong); pret. 1 sg. drēag; pret. pl. drugon; past ptcp. drogen; pres. 3 sg. drīegð. endure.\par
% glossary-block:B00333:end
% glossary-block:B00334:begin
\needspace{3\baselineskip}
\textbf{drīegð} — pres. ind. 3 sg.: see \textbf{drēogan}.\par
% glossary-block:B00334:end
% glossary-block:B00335:begin
\needspace{3\baselineskip}
\textbf{drīehð} — pres. ind. 3 sg., variant: see \textbf{drēogan}.\par
% glossary-block:B00335:end
% glossary-block:B00336:begin
\needspace{3\baselineskip}
\textbf{drīfan} (v. strong); pret. 1 sg. drāf; pret. pl. drifon; past ptcp. drifen. drive. Accusative person or thing driven.\par
% glossary-block:B00336:end
% glossary-block:B00337:begin
\needspace{3\baselineskip}
\textbf{drincan} (v. strong); pret. 1 sg. dranc; pret. pl. druncon; past ptcp. druncen. drink. Accusative liquid drunk.\par
% glossary-block:B00337:end
% glossary-block:B00338:begin
\needspace{3\baselineskip}
\textbf{drȳgnes} (n. f.); gen. sg. drȳgnesse; nom. pl. drȳgnessa. dry land.\par
% glossary-block:B00338:end
% glossary-block:B00339:begin
\needspace{3\baselineskip}
\textbf{dryhten} (n. m.); gen. sg. dryhtnes; nom. pl. dryhtnas. lord; the Lord.\par
% glossary-block:B00339:end
% glossary-block:B00340:begin
\needspace{3\baselineskip}
\textbf{dūn} (n. f.). mountain; hill.\par
% glossary-block:B00340:end
% glossary-block:B00341:begin
\needspace{3\baselineskip}
\textbf{durran} (v. preterite-present); pres. 1/3 sg. dearr; pres. pl. durron; pret. 1 sg. dorste. dare to. Conventional citation label; no usable infinitive is supplied. See the finite forms in Unit 6. No past-participle citation is given in the checked source paradigm.\par
% glossary-block:B00341:end
% glossary-block:B00342:begin
\needspace{3\baselineskip}
\textbf{dūst} (n. n.); gen. sg. dūstes; nom. pl. dūst. dust.\par
% glossary-block:B00342:end
% glossary-block:B00343:begin
\needspace{3\baselineskip}
\textbf{dyde} — pret. sg.: see \textbf{dōn}.\par
% glossary-block:B00343:end
\pdfbookmark[1]{E}{outline-B00344}
% glossary-block:B00344:begin
\needspace{5\baselineskip}
\subsection*{E}
% glossary-block:B00344:end
% glossary-block:B00345:begin
\textbf{ēac} (adv.). also. Also in the combination ēac swilce.\par
% glossary-block:B00345:end
% glossary-block:B00346:begin
\needspace{3\baselineskip}
\textbf{ēadig} (adj.). happy; blessed.\par
% glossary-block:B00346:end
% glossary-block:B00347:begin
\needspace{3\baselineskip}
\textbf{ēage} (n. n.); gen. sg. ēagan; nom. pl. ēagan. eye.\par
% glossary-block:B00347:end
% glossary-block:B00348:begin
\needspace{3\baselineskip}
\textbf{ēagðȳrel} (n. n.); gen. sg. ēagðȳrles; nom. pl. ēagðȳrlu. window.\par
% glossary-block:B00348:end
% glossary-block:B00349:begin
\needspace{3\baselineskip}
\textbf{eahta} (cardinal numeral). eight. Cardinal 8.\par
% glossary-block:B00349:end
% glossary-block:B00350:begin
\needspace{3\baselineskip}
\textbf{ēalā} (interj.). O!. Interjection in address, exclamation or wish.\par
% glossary-block:B00350:end
% glossary-block:B00351:begin
\needspace{3\baselineskip}
\textbf{eald} (adj.); comp. ieldra; superl. ieldest. old. old people; the elderly. Substantival adjective.\par
% glossary-block:B00351:end
% glossary-block:B00352:begin
\needspace{3\baselineskip}
\textbf{ealdor} (n. m.); gen. sg. ealdres; nom. pl. ealdras. chief; king. This noun is distinct from the homonymous noun meaning life.\par
% glossary-block:B00352:end
% glossary-block:B00353:begin
\needspace{3\baselineskip}
\textbf{ealdordōm} (n. m.). chief place; primacy.\par
% glossary-block:B00353:end
% glossary-block:B00354:begin
\needspace{3\baselineskip}
\textbf{ealdormann} (n. m.); gen. sg. ealdormannes; nom. pl. ealdormęnn. prince; leader; noble.\par
% glossary-block:B00354:end
% glossary-block:B00355:begin
\needspace{3\baselineskip}
\textbf{ealdorscipe} (n. m.); gen. sg. ealdorscipes; nom. pl. ealdorscipas. chief place; primacy.\par
% glossary-block:B00355:end
% glossary-block:B00356:begin
\needspace{3\baselineskip}
\textbf{eall} (adj.; substantive; adv.). all; whole. all people; everything. Substantival: plural for people, neuter singular for everything. completely. Only in the fixed phrase mid ealle. just; exactly. Degree emphasis in eall swā.\par
% glossary-block:B00356:end
% glossary-block:B00357:begin
\needspace{3\baselineskip}
\textbf{ēare} (n. n.); gen. sg. ēaran; nom. pl. ēaran. ear.\par
% glossary-block:B00357:end
% glossary-block:B00358:begin
\needspace{3\baselineskip}
\textbf{earfoðrīme} (adj.). innumerable; hard to count. Assigned source adjective.\par
% glossary-block:B00358:end
% glossary-block:B00359:begin
\needspace{3\baselineskip}
\textbf{earg} (adj.). cowardly.\par
% glossary-block:B00359:end
% glossary-block:B00360:begin
\needspace{3\baselineskip}
\textbf{earm} (n.). arm. Pronunciation example.\par
% glossary-block:B00360:end
% glossary-block:B00361:begin
\needspace{3\baselineskip}
\textbf{earm} (adj.). the poor; the unfortunate. Substantival adjective.\par
% glossary-block:B00361:end
% glossary-block:B00362:begin
\needspace{3\baselineskip}
\textbf{earmlice} (adverb or particle). cruelly; miserably.\par
% glossary-block:B00362:end
% glossary-block:B00363:begin
\needspace{3\baselineskip}
\textbf{earnung} (n. f.). merit; good service. Source plural merits/services.\par
% glossary-block:B00363:end
% glossary-block:B00364:begin
\needspace{3\baselineskip}
\textbf{eart} — pres. ind. 2 sg.: see \textbf{wesan}.\par
% glossary-block:B00364:end
% glossary-block:B00365:begin
\needspace{3\baselineskip}
\textbf{Eastęngle} (proper name). East Anglians.\par
% glossary-block:B00365:end
% glossary-block:B00366:begin
\needspace{3\baselineskip}
\textbf{ēastnorðerne} (adj.). northeasterly.\par
% glossary-block:B00366:end
% glossary-block:B00367:begin
\needspace{3\baselineskip}
\textbf{ēastsūðdǣl} (n. m.). southeastern part.\par
% glossary-block:B00367:end
% glossary-block:B00368:begin
\needspace{3\baselineskip}
\textbf{ēaðmōd} (adj.). lowly; humble. humble people; the lowly. Substantival adjective.\par
% glossary-block:B00368:end
% glossary-block:B00369:begin
\needspace{3\baselineskip}
\textbf{ēaðmōdnes} (n. f.); gen. sg. ēaðmōdnesse. humility.\par
% glossary-block:B00369:end
% glossary-block:B00370:begin
\needspace{3\baselineskip}
\textbf{ēce} (adj.). eternal.\par
% glossary-block:B00370:end
% glossary-block:B00371:begin
\needspace{3\baselineskip}
\textbf{ecg} (n.). edge. Pronunciation example.\par
% glossary-block:B00371:end
% glossary-block:B00372:begin
\needspace{3\baselineskip}
\textbf{edlēan} (n. n.); gen. sg. edlēanes; nom. pl. edlēan. reward; recompense.\par
% glossary-block:B00372:end
% glossary-block:B00373:begin
\needspace{3\baselineskip}
\textbf{ęfnan} (v. weak); pret. 1 sg. ęfnede; past ptcp. (ge)ęfned. perform.\par
% glossary-block:B00373:end
% glossary-block:B00374:begin
\needspace{3\baselineskip}
\textbf{efne} (interj.). behold!. Interjection, distinct from the verb behealdan look at. just; precisely; even.\par
% glossary-block:B00374:end
% glossary-block:B00375:begin
\needspace{3\baselineskip}
\textbf{eft, ęft} (adv.). again; also. The source note also permits the additive sense. again; back. again; afterwards.\par
% glossary-block:B00375:end
% glossary-block:B00376:begin
\needspace{3\baselineskip}
\textbf{ęfthweorfan} (v. strong); pret. 1 sg. ęfthwearf; pret. pl. ęfthwurfon; past ptcp. ęfthworfen; pres. 3 sg. ęfthwierfð. return; turn back.\par
% glossary-block:B00376:end
% glossary-block:B00377:begin
\needspace{3\baselineskip}
\textbf{ęfthwierfð} — pres. ind. 3 sg.: see \textbf{ęfthweorfan}.\par
% glossary-block:B00377:end
% glossary-block:B00378:begin
\needspace{3\baselineskip}
\textbf{ęgesa} (n. m.). fear; terror. Reference in a source or explanatory note.\par
% glossary-block:B00378:end
% glossary-block:B00379:begin
\needspace{3\baselineskip}
\textbf{ęgeslic} (adj.). dreadful; terrible.\par
% glossary-block:B00379:end
% glossary-block:B00380:begin
\needspace{3\baselineskip}
\textbf{Egypta} (proper name). Egyptians. Egypta land denotes Egypt, literally land of the Egyptians; supplied name forms. Supplied name forms: Egypta / Ægypta, genitive plural; Egyptum, dative plural. Egypta land is the land of the Egyptians.\par
% glossary-block:B00380:end
% glossary-block:B00381:begin
\needspace{3\baselineskip}
\textbf{Egypte} (ethnonym). Egyptians.\par
% glossary-block:B00381:end
% glossary-block:B00382:begin
\needspace{3\baselineskip}
\textbf{ēhtan} (v. weak); pret. 1 sg. ēhte; past ptcp. (ge)ēht. pursue; persecute. Accusative person in the selected production frame; a separately supported genitive alternative is not the selected target.\par
% glossary-block:B00382:end
% glossary-block:B00383:begin
\needspace{3\baselineskip}
\textbf{ęlcung} (n. f.); gen. sg. ęlcunge. delay.\par
% glossary-block:B00383:end
% glossary-block:B00384:begin
\needspace{3\baselineskip}
\textbf{ęle} (n. m.); gen. sg. ęles; nom. pl. ęlas. oil.\par
% glossary-block:B00384:end
% glossary-block:B00385:begin
\needspace{3\baselineskip}
\textbf{elp} (n. m.); gen. sg. elpes; nom. pl. elpas. elephant.\par
% glossary-block:B00385:end
% glossary-block:B00386:begin
\needspace{3\baselineskip}
\textbf{elpenbǣnen} (adj.). ivory; made of ivory. Adjective of material.\par
% glossary-block:B00386:end
% glossary-block:B00387:begin
\needspace{3\baselineskip}
\textbf{ęlðēodig} (adj.). foreigners; strangers. Substantival adjective. foreign.\par
% glossary-block:B00387:end
% glossary-block:B00388:begin
\needspace{3\baselineskip}
\textbf{ęnde} (n. m.). end.\par
% glossary-block:B00388:end
% glossary-block:B00389:begin
\needspace{3\baselineskip}
\textbf{endlefan} (cardinal numeral). eleven. Cardinal 11.\par
% glossary-block:B00389:end
% glossary-block:B00390:begin
\needspace{3\baselineskip}
\textbf{endlefta} (ordinal). eleventh.\par
% glossary-block:B00390:end
% glossary-block:B00391:begin
\needspace{3\baselineskip}
\textbf{endleofan} (numeral). eleven.\par
% glossary-block:B00391:end
% glossary-block:B00392:begin
\needspace{3\baselineskip}
\textbf{ęngel} (n. m.); gen. sg. ęngles; nom. pl. ęnglas. angel.\par
% glossary-block:B00392:end
% glossary-block:B00393:begin
\needspace{3\baselineskip}
\textbf{Engle} (n. m.); nom. pl. Engle. the English. People, distinct from the language/adjective.\par
% glossary-block:B00393:end
% glossary-block:B00394:begin
\needspace{3\baselineskip}
\textbf{Englisc} (adj.). English.\par
% glossary-block:B00394:end
% glossary-block:B00395:begin
\needspace{3\baselineskip}
\textbf{Englisc} (adjective used substantivally). English, the language. Language noun.\par
% glossary-block:B00395:end
% glossary-block:B00396:begin
\needspace{3\baselineskip}
\textbf{ēode} — pret. sg.: see \textbf{gān}.\par
% glossary-block:B00396:end
% glossary-block:B00397:begin
\needspace{3\baselineskip}
\textbf{eom} — pres. ind. 1 sg.: see \textbf{wesan}.\par
% glossary-block:B00397:end
% glossary-block:B00398:begin
\needspace{3\baselineskip}
\textbf{eorl} (n.). man. Pronunciation example.\par
% glossary-block:B00398:end
% glossary-block:B00399:begin
\needspace{3\baselineskip}
\textbf{eornostlice} (adv.). accordingly; thus.\par
% glossary-block:B00399:end
% glossary-block:B00400:begin
\needspace{3\baselineskip}
\textbf{eorðe} (n. f.); gen. sg. eorðan; nom. pl. eorðan. earth.\par
% glossary-block:B00400:end
% glossary-block:B00401:begin
\needspace{3\baselineskip}
\textbf{eorðlic} (adj.). earthly.\par
% glossary-block:B00401:end
% glossary-block:B00402:begin
\needspace{3\baselineskip}
\textbf{eorðtilð} (n. f.); gen. sg. eorðtilðe; nom. pl. eorðtilða. agriculture.\par
% glossary-block:B00402:end
% glossary-block:B00403:begin
\needspace{3\baselineskip}
\textbf{ēow} — dat. pl.; acc. pl.: see \textbf{ðū}.\par
% glossary-block:B00403:end
% glossary-block:B00404:begin
\needspace{3\baselineskip}
\textbf{ēower} (possessive adjective). your; yours. Plural possessors; inflection agrees with the possessed noun. Also gen. pl.: see \textbf{ðū}.\par
% glossary-block:B00404:end
% glossary-block:B00405:begin
\needspace{3\baselineskip}
\textbf{ēowic} — acc. pl., variant: see \textbf{ðū}.\par
% glossary-block:B00405:end
% glossary-block:B00406:begin
\needspace{3\baselineskip}
\textbf{ęrian} (v.). plow.\par
% glossary-block:B00406:end
% glossary-block:B00407:begin
\needspace{3\baselineskip}
\textbf{ēst} (n. m.). provisions; food.\par
% glossary-block:B00407:end
% glossary-block:B00408:begin
\needspace{3\baselineskip}
\textbf{etan} (v. strong); pret. 1 sg. ǣt; pret. pl. ǣton; past ptcp. eten; pres. 3 sg. itt. eat.\par
% glossary-block:B00408:end
% glossary-block:B00409:begin
\needspace{3\baselineskip}
\textbf{ēðel} (n. m.); gen. sg. ēðles; nom. pl. ēðlas. native land; homeland.\par
% glossary-block:B00409:end
\pdfbookmark[1]{F}{outline-B00410}
% glossary-block:B00410:begin
\needspace{5\baselineskip}
\subsection*{F}
% glossary-block:B00410:end
% glossary-block:B00411:begin
\textbf{faran} (v. strong); pret. 1 sg. fōr; pret. pl. fōron; past ptcp. faren; pres. 3 sg. færð. go; travel.\par
% glossary-block:B00411:end
% glossary-block:B00412:begin
\needspace{3\baselineskip}
\textbf{Farne} (proper name). Farne. Island; supplied proper-name form after in.\par
% glossary-block:B00412:end
% glossary-block:B00413:begin
\needspace{3\baselineskip}
\textbf{fæc} (n. n.). interval; period of time; while. Noun in the supplied time expression.\par
% glossary-block:B00413:end
% glossary-block:B00414:begin
\needspace{3\baselineskip}
\textbf{fæder} (n. m.); gen. sg. fæder; nom. pl. fæderas. father. Includes the spiritual father relationship in the assigned prompts.\par
% glossary-block:B00414:end
% glossary-block:B00415:begin
\needspace{3\baselineskip}
\textbf{fæger} (adj.); comp. fægerra; supplied superl. acc. f. sg. fægerestan. beautiful; fair.\par
% glossary-block:B00415:end
% glossary-block:B00416:begin
\needspace{3\baselineskip}
\textbf{fægerest} (adj.). most beautiful. Superlative of fæger; source form fægerestan.\par
% glossary-block:B00416:end
% glossary-block:B00417:begin
\needspace{3\baselineskip}
\textbf{fægerra} (adj.). more beautiful; fairer. Comparative of fæger.\par
% glossary-block:B00417:end
% glossary-block:B00418:begin
\needspace{3\baselineskip}
\textbf{fǣmne} (n. f.); gen. sg. fǣmnan; nom. pl. fǣmnan. maiden; young woman.\par
% glossary-block:B00418:end
% glossary-block:B00419:begin
\needspace{3\baselineskip}
\textbf{fǣringa} (adv.). suddenly.\par
% glossary-block:B00419:end
% glossary-block:B00420:begin
\needspace{3\baselineskip}
\textbf{færð} — pres. ind. 3 sg.: see \textbf{faran}.\par
% glossary-block:B00420:end
% glossary-block:B00421:begin
\needspace{3\baselineskip}
\textbf{fæsten} (n. n.). fortification; stronghold.\par
% glossary-block:B00421:end
% glossary-block:B00422:begin
\needspace{3\baselineskip}
\textbf{fæsthafol} (adj.). tenacious; withholding. Genitive of what is held back in the source.\par
% glossary-block:B00422:end
% glossary-block:B00423:begin
\needspace{3\baselineskip}
\textbf{fæstnes} (n. f.); gen. sg. fæstnesse; nom. pl. fæstnessa. firmament.\par
% glossary-block:B00423:end
% glossary-block:B00424:begin
\needspace{3\baselineskip}
\textbf{fæt} (n. n.); gen. sg. fætes; nom. pl. fatu. vessel; container.\par
% glossary-block:B00424:end
% glossary-block:B00425:begin
\needspace{3\baselineskip}
\textbf{feallan} (v. strong); pret. 1 sg. fēoll; pret. pl. fēollon; past ptcp. feallen; pres. 3 sg. fielð. fall.\par
% glossary-block:B00425:end
% glossary-block:B00426:begin
\needspace{3\baselineskip}
\textbf{fēawe} (adj.). few people. Substantival use in the assigned expression. few.\par
% glossary-block:B00426:end
% glossary-block:B00427:begin
\needspace{3\baselineskip}
\textbf{feccean} (v.). fetch.\par
% glossary-block:B00427:end
% glossary-block:B00428:begin
\needspace{3\baselineskip}
\textbf{fēdan} (v. weak); pret. 1 sg. fēdde; past ptcp. fēded. feed; nourish. Accusative being fed; food may be mid + dative or a bare instrumental in the assigned expressions.\par
% glossary-block:B00428:end
% glossary-block:B00429:begin
\needspace{3\baselineskip}
\textbf{fēhð} — pres. ind. 3 sg.: see \textbf{fōn}.\par
% glossary-block:B00429:end
% glossary-block:B00430:begin
\needspace{3\baselineskip}
\textbf{fela} (adj.). many; a large number of. Genitive of the things counted in the taught frame.\par
% glossary-block:B00430:end
% glossary-block:B00431:begin
\needspace{3\baselineskip}
\textbf{felgan} — citation/spelling variant: see \textbf{fēolan}.\par
% glossary-block:B00431:end
% glossary-block:B00432:begin
\needspace{3\baselineskip}
\textbf{feoh} (n. n.); gen. sg. fēos. money.\par
% glossary-block:B00432:end
% glossary-block:B00433:begin
\needspace{3\baselineskip}
\textbf{feohtan} (v. strong); pret. 1 sg. feaht; pret. pl. fuhton; past ptcp. fohten; pres. 3 sg. fieht. fight; contend. Intransitive; wið + accusative opponent.\par
% glossary-block:B00433:end
% glossary-block:B00434:begin
\needspace{3\baselineskip}
\textbf{fēolan} (v. strong); pret. 1 sg. fealh; pret. pl. fulgon; past ptcp. folen; pres. 3 sg. filhð. penetrate. Preserve the supplied/taught lexical frame; not a synonym for every kind of entering.\par
% glossary-block:B00434:end
% glossary-block:B00435:begin
\needspace{3\baselineskip}
\textbf{fēond} (n. m.); gen. sg. fēondes; nom. pl. fīend. enemy; foe.\par
% glossary-block:B00435:end
% glossary-block:B00436:begin
\needspace{3\baselineskip}
\textbf{feorr} (adv.). far; far off.\par
% glossary-block:B00436:end
% glossary-block:B00437:begin
\needspace{3\baselineskip}
\textbf{feorran} (adv.). from afar.\par
% glossary-block:B00437:end
% glossary-block:B00438:begin
\needspace{3\baselineskip}
\textbf{fēorða} (ordinal numeral). fourth.\par
% glossary-block:B00438:end
% glossary-block:B00439:begin
\needspace{3\baselineskip}
\textbf{fēower} (cardinal numeral). four. Cardinal 4.\par
% glossary-block:B00439:end
% glossary-block:B00440:begin
\needspace{3\baselineskip}
\textbf{fēowertig} (numeral). forty.\par
% glossary-block:B00440:end
% glossary-block:B00441:begin
\needspace{3\baselineskip}
\textbf{fēran} (v.). go; journey. Reference in a source or explanatory note.\par
% glossary-block:B00441:end
% glossary-block:B00442:begin
\needspace{3\baselineskip}
\textbf{fēt} — pres. ind. 3 sg.: see \textbf{fēdan}; dat. sg.; nom./acc. pl.: see \textbf{fōt}.\par
% glossary-block:B00442:end
% glossary-block:B00443:begin
\needspace{3\baselineskip}
\textbf{fieht} — pres. ind. 3 sg.: see \textbf{feohtan}.\par
% glossary-block:B00443:end
% glossary-block:B00444:begin
\needspace{3\baselineskip}
\textbf{fielð} — pres. ind. 3 sg.: see \textbf{feallan}.\par
% glossary-block:B00444:end
% glossary-block:B00445:begin
\needspace{3\baselineskip}
\textbf{fielþ} — pres. ind. 3 sg., variant: see \textbf{feallan}.\par
% glossary-block:B00445:end
% glossary-block:B00446:begin
\needspace{3\baselineskip}
\textbf{fīend} — nom./acc. pl., variant: see \textbf{fēond}.\par
% glossary-block:B00446:end
% glossary-block:B00447:begin
\needspace{3\baselineskip}
\textbf{fierst} (n. m.). period; interval. Masculine noun in the source time expression.\par
% glossary-block:B00447:end
% glossary-block:B00448:begin
\needspace{3\baselineskip}
\textbf{fīf} (cardinal numeral). five. Cardinal 5.\par
% glossary-block:B00448:end
% glossary-block:B00449:begin
\needspace{3\baselineskip}
\textbf{fīfta} (ordinal numeral). fifth.\par
% glossary-block:B00449:end
% glossary-block:B00450:begin
\needspace{3\baselineskip}
\textbf{fīftig} (numeral). fifty.\par
% glossary-block:B00450:end
% glossary-block:B00451:begin
\needspace{3\baselineskip}
\textbf{filhð} — pres. ind. 3 sg.: see \textbf{fēolan}.\par
% glossary-block:B00451:end
% glossary-block:B00452:begin
\needspace{3\baselineskip}
\textbf{filhþ} — pres. ind. 3 sg., variant: see \textbf{fēolan}.\par
% glossary-block:B00452:end
% glossary-block:B00453:begin
\needspace{3\baselineskip}
\textbf{findan} (v. strong); pret. 1 sg. fand; pret. pl. fundon; past ptcp. funden. find. Accusative person or thing found.\par
% glossary-block:B00453:end
% glossary-block:B00454:begin
\needspace{3\baselineskip}
\textbf{finger} (n. m.); gen. sg. fingres; nom. pl. fingras. finger.\par
% glossary-block:B00454:end
% glossary-block:B00455:begin
\needspace{3\baselineskip}
\textbf{fisc} (n. m.); gen. sg. fisces; nom. pl. fiscas. fish.\par
% glossary-block:B00455:end
% glossary-block:B00456:begin
\needspace{3\baselineskip}
\textbf{fisccynn} (n. n.). kinds of fish. Supplied accusative neuter plural.\par
% glossary-block:B00456:end
% glossary-block:B00457:begin
\needspace{3\baselineskip}
\textbf{fiscere} (n. m.); gen. sg. fisceres; nom. pl. fisceras. fisherman.\par
% glossary-block:B00457:end
% glossary-block:B00458:begin
\needspace{3\baselineskip}
\textbf{fiscnoð} (n. m.). fishing.\par
% glossary-block:B00458:end
% glossary-block:B00459:begin
\needspace{3\baselineskip}
\textbf{flǣsc} (n. n.); gen. sg. flǣsces; nom. pl. flǣsc. flesh.\par
% glossary-block:B00459:end
% glossary-block:B00460:begin
\needspace{3\baselineskip}
\textbf{flēogan} (v. strong); pret. 1 sg. flēag; pret. pl. flugon; past ptcp. flogen; pres. 3 sg. flīehð. fly.\par
% glossary-block:B00460:end
% glossary-block:B00461:begin
\needspace{3\baselineskip}
\textbf{flēon} (v. strong); pret. 1 sg. flēah; pret. pl. flugon; past ptcp. flogen; pres. 3 sg. flīehð. flee. Supplied third-plural present; distinct from flēogan fly. flee; flee from; avoid. Intransitive or with accusative person/thing avoided.\par
% glossary-block:B00461:end
% glossary-block:B00462:begin
\needspace{3\baselineskip}
\textbf{flēwð} — pres. ind. 3 sg.: see \textbf{flōwan}.\par
% glossary-block:B00462:end
% glossary-block:B00463:begin
\needspace{3\baselineskip}
\textbf{flīegð} — pres. ind. 3 sg., variant: see \textbf{flēogan}.\par
% glossary-block:B00463:end
% glossary-block:B00464:begin
\needspace{3\baselineskip}
\textbf{flīehð} — pres. ind. 3 sg.: see \textbf{flēogan}; pres. ind. 3 sg.: see \textbf{flēon}.\par
% glossary-block:B00464:end
% glossary-block:B00465:begin
\needspace{3\baselineskip}
\textbf{flocc} (n. m.). group; company.\par
% glossary-block:B00465:end
% glossary-block:B00466:begin
\needspace{3\baselineskip}
\textbf{flōd} (n. m.); gen. sg. flōdes; nom. pl. flōdas. flood; stream of water.\par
% glossary-block:B00466:end
% glossary-block:B00467:begin
\needspace{3\baselineskip}
\textbf{flōwan} (v. strong); pret. 1 sg. flēow; pret. pl. flēowon; past ptcp. flōwen; pres. 3 sg. flēwð. flow.\par
% glossary-block:B00467:end
% glossary-block:B00468:begin
\needspace{3\baselineskip}
\textbf{fōdor} (n. n.); gen. sg. fōdres; nom. pl. fōdru. fodder.\par
% glossary-block:B00468:end
% glossary-block:B00469:begin
\needspace{3\baselineskip}
\textbf{folc} (n. n.); gen. sg. folces; nom. pl. folc. people. Collective singular. peoples. Distinct peoples or national groups; plural sense.\par
% glossary-block:B00469:end
% glossary-block:B00470:begin
\needspace{3\baselineskip}
\textbf{fōn} (v. strong); pret. 1 sg. fēng; pret. pl. fēngon; past ptcp. fangen; pres. 3 sg. fēhð. seize; catch. Accusative patient in the selected frame.\par
% glossary-block:B00470:end
% glossary-block:B00471:begin
\needspace{3\baselineskip}
\textbf{for} (prep.). for; for the sake of. Case follows each complete supplied accusative/dative expression. for; on behalf of. Dative beneficiary. why; for what reason. In for hwȳ or for hwon. because of. Dative hindering circumstance. by; for the sake of. Grounds of entreaty in the prayer; keep the supplied phrase. in exchange for. With dative. instead of; in place of. Accusative in the supplied substitution frame. because of; on account of. for; to the benefit of.\par
% glossary-block:B00471:end
% glossary-block:B00472:begin
\needspace{3\baselineskip}
\textbf{forbærnan} (v. weak); pret. 1 sg. forbærnde; past ptcp. forbærned. parch; scorch.\par
% glossary-block:B00472:end
% glossary-block:B00473:begin
\needspace{3\baselineskip}
\textbf{forceorfan} (v. strong); pret. 1 sg. forcearf; pret. pl. forcurfon; past ptcp. forcorfen; pres. 3 sg. forcierfð. cut down; cut off. Accusative tree or thing cut.\par
% glossary-block:B00473:end
% glossary-block:B00474:begin
\needspace{3\baselineskip}
\textbf{forcierfð} — pres. ind. 3 sg.: see \textbf{forceorfan}.\par
% glossary-block:B00474:end
% glossary-block:B00475:begin
\needspace{3\baselineskip}
\textbf{fordēð} — pres. ind. 3 sg.: see \textbf{fordōn}.\par
% glossary-block:B00475:end
% glossary-block:B00476:begin
\needspace{3\baselineskip}
\textbf{fordilgian} (v. weak); pret. 1 sg. fordilgode; past ptcp. fordilgod. destroy; annihilate.\par
% glossary-block:B00476:end
% glossary-block:B00477:begin
\needspace{3\baselineskip}
\textbf{fordōn} (v. irregular); pres. 1 sg. fordō; pres. 3 sg. fordēð; pres. pl. fordōð; pret. 1 sg. fordyde; pret. pl. fordydon; past ptcp. fordōn. destroy.\par
% glossary-block:B00477:end
% glossary-block:B00478:begin
\needspace{3\baselineskip}
\textbf{fordyde} — pret. sg.: see \textbf{fordōn}.\par
% glossary-block:B00478:end
% glossary-block:B00479:begin
\needspace{3\baselineskip}
\textbf{foresægde} — pret. sg.: see \textbf{foresęcgean}.\par
% glossary-block:B00479:end
% glossary-block:B00480:begin
\needspace{3\baselineskip}
\textbf{foresægð} — pres. ind. 3 sg.: see \textbf{foresęcgean}.\par
% glossary-block:B00480:end
% glossary-block:B00481:begin
\needspace{3\baselineskip}
\textbf{foresęcgean} (v. weak); pret. 1 sg. foresægde; past ptcp. foresægd; pres. 3 sg. foresægð. mention before; tell of beforehand.\par
% glossary-block:B00481:end
% glossary-block:B00482:begin
\needspace{3\baselineskip}
\textbf{forgiefan} (v. strong); pret. 1 sg. forgeaf; pret. pl. forgēafon; past ptcp. forgiefen. give; grant. Supplied granted sense at Cook 143.18; not forgive here. give; grant. Accusative thing given and dative recipient; the thing may be understood from context. forgive. Accusative sin, dative person forgiven; in a passive the sin becomes nominative and the person remains dative.\par
% glossary-block:B00482:end
% glossary-block:B00483:begin
\needspace{3\baselineskip}
\textbf{forgieldan} (v. strong); pret. 1 sg. forgeald; pret. pl. forguldon; past ptcp. forgolden. pay; repay; requite. Dative recipient, accusative thing repaid in the selected frame.\par
% glossary-block:B00483:end
% glossary-block:B00484:begin
\needspace{3\baselineskip}
\textbf{forierman} (v.). ruin; reduce to misery.\par
% glossary-block:B00484:end
% glossary-block:B00485:begin
\needspace{3\baselineskip}
\textbf{forlǣtan} (v. strong); pret. 1 sg. forlēt; pret. pl. forlēton; past ptcp. forlǣten. abandon; forsake. Accusative person or thing left. leave to. Thing left and recipient retain their taught roles. let go; dismiss. set aside. Supplied source sense.\par
% glossary-block:B00485:end
% glossary-block:B00486:begin
\needspace{3\baselineskip}
\textbf{forlēosan} (v. strong); pret. 1 sg. forlēas; pret. pl. forluron; past ptcp. forloren; pres. 3 sg. forlīest. lose. Accusative thing lost.\par
% glossary-block:B00486:end
% glossary-block:B00487:begin
\needspace{3\baselineskip}
\textbf{forliden} (adj.). shipwrecked.\par
% glossary-block:B00487:end
% glossary-block:B00488:begin
\needspace{3\baselineskip}
\textbf{forlidennes} (n. f.); gen. sg. forlidennesse; nom. pl. forlidennessa. shipwreck.\par
% glossary-block:B00488:end
% glossary-block:B00489:begin
\needspace{3\baselineskip}
\textbf{forlīest} — pres. ind. 3 sg.: see \textbf{forlēosan}.\par
% glossary-block:B00489:end
% glossary-block:B00490:begin
\needspace{3\baselineskip}
\textbf{forniman} (v.). destroy; lay waste.\par
% glossary-block:B00490:end
% glossary-block:B00491:begin
\needspace{3\baselineskip}
\textbf{forstelan} (v. strong); pret. 1 sg. forstæl; pret. pl. forstǣlon; past ptcp. forstolen; pres. 3 sg. forstilð. steal away; steal.\par
% glossary-block:B00491:end
% glossary-block:B00492:begin
\needspace{3\baselineskip}
\textbf{forstilð} — pres. ind. 3 sg.: see \textbf{forstelan}.\par
% glossary-block:B00492:end
% glossary-block:B00493:begin
\needspace{3\baselineskip}
\textbf{forswīgian} (v.). keep silent about; conceal.\par
% glossary-block:B00493:end
% glossary-block:B00494:begin
\needspace{3\baselineskip}
\textbf{fortrēdan} (v.). trample down.\par
% glossary-block:B00494:end
% glossary-block:B00495:begin
\needspace{3\baselineskip}
\textbf{forð} (adv.). forth. Adverbial particle with the bringing/producing verb.\par
% glossary-block:B00495:end
% glossary-block:B00496:begin
\needspace{3\baselineskip}
\textbf{for-ðām} (conj.). because.\par
% glossary-block:B00496:end
% glossary-block:B00497:begin
\needspace{3\baselineskip}
\textbf{forðām} (adv.). therefore. Pronunciation example.\par
% glossary-block:B00497:end
% glossary-block:B00498:begin
\needspace{3\baselineskip}
\textbf{for-ðām-ðe} (conj.). because.\par
% glossary-block:B00498:end
% glossary-block:B00499:begin
\needspace{3\baselineskip}
\textbf{forðfōr} (n. f.); gen. sg. forðfōre; nom. pl. forðfōra. departure.\par
% glossary-block:B00499:end
% glossary-block:B00500:begin
\needspace{3\baselineskip}
\textbf{for-ðon-ðe} (conj.; adv. in for-ðon). because. Causal conjunction. therefore; for that reason. Adverbial consequence, distinct from because.\par
% glossary-block:B00500:end
% glossary-block:B00501:begin
\needspace{3\baselineskip}
\textbf{forðtēon} (v.). present; perform. Source performance sense; retain the actual form's spelling.\par
% glossary-block:B00501:end
% glossary-block:B00502:begin
\needspace{3\baselineskip}
\textbf{forweorðan} (v. strong); pret. 1 sg. forwearð; pret. pl. forwurdon; past ptcp. forworden; pres. 3 sg. forwierð. perish.\par
% glossary-block:B00502:end
% glossary-block:B00503:begin
\needspace{3\baselineskip}
\textbf{forwiernan} (v. weak); pret. 1 sg. forwiernde; past ptcp. forwierned. deny; refuse. Dative person denied, genitive thing denied.\par
% glossary-block:B00503:end
% glossary-block:B00504:begin
\needspace{3\baselineskip}
\textbf{forwierð} — pres. ind. 3 sg.: see \textbf{forweorðan}.\par
% glossary-block:B00504:end
% glossary-block:B00505:begin
\needspace{3\baselineskip}
\textbf{fōt} (n. m.); gen. sg. fōtes; nom. pl. fēt. foot. Supplied body-part form feet; distinct from the length-measure sense. foot. Unit of length, distinct from the body-part sense.\par
% glossary-block:B00505:end
% glossary-block:B00506:begin
\needspace{3\baselineskip}
\textbf{fōtsceamul} (n. m.); gen. sg. fōtsceamules; nom. pl. fōtsceamulas. footstool.\par
% glossary-block:B00506:end
% glossary-block:B00507:begin
\needspace{3\baselineskip}
\textbf{fox} (n. m.); gen. sg. foxes; nom. pl. foxas. fox.\par
% glossary-block:B00507:end
% glossary-block:B00508:begin
\needspace{3\baselineskip}
\textbf{fram} (prep.). from. Separation, with dative in the taught frame. from. Only the supplied rest-from-work construction.\par
% glossary-block:B00508:end
% glossary-block:B00509:begin
\needspace{3\baselineskip}
\textbf{frætwung} (n. f.); gen. sg. frætwunge; nom. pl. frætwunga. array.\par
% glossary-block:B00509:end
% glossary-block:B00510:begin
\needspace{3\baselineskip}
\textbf{frēa} (n. m.). lord. Reference in a source or explanatory note.\par
% glossary-block:B00510:end
% glossary-block:B00511:begin
\needspace{3\baselineskip}
\textbf{frēcnes} (n. f.). danger; peril.\par
% glossary-block:B00511:end
% glossary-block:B00512:begin
\needspace{3\baselineskip}
\textbf{fręmman} (v. weak); pret. 1 sg. fręmede; past ptcp. (ge)fręmed. benefit. Dative beneficiary.\par
% glossary-block:B00512:end
% glossary-block:B00513:begin
\needspace{3\baselineskip}
\textbf{frēo} (adj.). free.\par
% glossary-block:B00513:end
% glossary-block:B00514:begin
\needspace{3\baselineskip}
\textbf{frēond} (n. m.); gen. sg. frēondes; nom. pl. frīend. friend.\par
% glossary-block:B00514:end
% glossary-block:B00515:begin
\needspace{3\baselineskip}
\textbf{frīend} — nom./acc. pl., variant: see \textbf{frēond}.\par
% glossary-block:B00515:end
% glossary-block:B00516:begin
\needspace{3\baselineskip}
\textbf{frōfor} (n. f.); gen. sg. frōfre; nom. pl. frōfra. comfort; consolation.\par
% glossary-block:B00516:end
% glossary-block:B00517:begin
\needspace{3\baselineskip}
\textbf{frosc} (n.). frog. Pronunciation example.\par
% glossary-block:B00517:end
% glossary-block:B00518:begin
\needspace{3\baselineskip}
\textbf{fugol} (n. m.); gen. sg. fugles; nom. pl. fuglas. bird.\par
% glossary-block:B00518:end
% glossary-block:B00519:begin
\needspace{3\baselineskip}
\textbf{fūl} (adj.). foul; vile.\par
% glossary-block:B00519:end
% glossary-block:B00520:begin
\needspace{3\baselineskip}
\textbf{full} (adj.). full. Pronunciation example.\par
% glossary-block:B00520:end
% glossary-block:B00521:begin
\needspace{3\baselineskip}
\textbf{full} (adv.). very.\par
% glossary-block:B00521:end
% glossary-block:B00522:begin
\needspace{3\baselineskip}
\textbf{fullfręmman} (v.). finish; complete. Supplied passive use.\par
% glossary-block:B00522:end
% glossary-block:B00523:begin
\needspace{3\baselineskip}
\textbf{fultum} (n. m.); gen. sg. fultumes. help; aid. help; support.\par
% glossary-block:B00523:end
% glossary-block:B00524:begin
\needspace{3\baselineskip}
\textbf{fuluhtere} (n. m.); gen. sg. fuluhteres. baptizer; Baptist. Baptist as John's title in the assigned passage.\par
% glossary-block:B00524:end
% glossary-block:B00525:begin
\needspace{3\baselineskip}
\textbf{furlang} (n. n.); gen. sg. furlanges; nom. pl. furlang. furlong.\par
% glossary-block:B00525:end
% glossary-block:B00526:begin
\needspace{3\baselineskip}
\textbf{furðor} (adv.). further; any more. Adverb in the supplied no-spirit-more expression.\par
% glossary-block:B00526:end
% glossary-block:B00527:begin
\needspace{3\baselineskip}
\textbf{furðra} (adj.). foremost; prior.\par
% glossary-block:B00527:end
% glossary-block:B00528:begin
\needspace{3\baselineskip}
\textbf{fȳr} (n. n.); gen. sg. fȳres; nom. pl. fȳr. fire.\par
% glossary-block:B00528:end
% glossary-block:B00529:begin
\needspace{3\baselineskip}
\textbf{fyrmest} (adv.). first. first of all; at the outset.\par
% glossary-block:B00529:end
% glossary-block:B00530:begin
\needspace{3\baselineskip}
\textbf{fȳrspearca} (n. m.); gen. sg. fȳrspearcan; nom. pl. fȳrspearcan. spark.\par
% glossary-block:B00530:end
\pdfbookmark[1]{G}{outline-B00531}
% glossary-block:B00531:begin
\needspace{5\baselineskip}
\subsection*{G}
% glossary-block:B00531:end
% glossary-block:B00532:begin
\textbf{gagates} (n.). jet (Latin loan-name). The stone; Latin loan-name in the source.\par
% glossary-block:B00532:end
% glossary-block:B00533:begin
\needspace{3\baselineskip}
\textbf{gālnes} (n. f.). lust. Assigned source sense.\par
% glossary-block:B00533:end
% glossary-block:B00534:begin
\needspace{3\baselineskip}
\textbf{gān} (v. irregular); pres. 1 sg. gā; pres. 3 sg. gǣð; pres. pl. gāð; pret. 1 sg. ēode; pret. pl. ēodon; past ptcp. (ge)gān. go; walk. go. Supplied imperative gang.\par
% glossary-block:B00534:end
% glossary-block:B00535:begin
\needspace{3\baselineskip}
\textbf{gangan} (v.). go; enter. Supplied circumstantial participial phrase.\par
% glossary-block:B00535:end
% glossary-block:B00536:begin
\needspace{3\baselineskip}
\textbf{gāst} (n. m.); gen. sg. gāstes; nom. pl. gāstas. spirit.\par
% glossary-block:B00536:end
% glossary-block:B00537:begin
\needspace{3\baselineskip}
\textbf{gāt} (n. f.); gen. sg. gāte; nom. pl. gǣt. goat.\par
% glossary-block:B00537:end
% glossary-block:B00538:begin
\needspace{3\baselineskip}
\textbf{gærs} (n. n.); gen. sg. gærses; nom. pl. gærs. grass.\par
% glossary-block:B00538:end
% glossary-block:B00539:begin
\needspace{3\baselineskip}
\textbf{gǣt} — dat. sg.; nom./acc. pl.; also gen. sg. variant: see \textbf{gāt}.\par
% glossary-block:B00539:end
% glossary-block:B00540:begin
\needspace{3\baselineskip}
\textbf{gǣð} — pres. ind. 3 sg.: see \textbf{gān}.\par
% glossary-block:B00540:end
% glossary-block:B00541:begin
\needspace{3\baselineskip}
\textbf{ge} (conj.). both ... and. both … and. Paired ge … ge.\par
% glossary-block:B00541:end
% glossary-block:B00542:begin
\needspace{3\baselineskip}
\textbf{gē} — nom. pl.: see \textbf{ðū}.\par
% glossary-block:B00542:end
% glossary-block:B00543:begin
\needspace{3\baselineskip}
\textbf{gēar} (n. n.); gen. sg. gēares; nom. pl. gēar. year.\par
% glossary-block:B00543:end
% glossary-block:B00544:begin
\needspace{3\baselineskip}
\textbf{gearcian} (v. weak); pret. 1 sg. gearcode; past ptcp. gearcod. prepare.\par
% glossary-block:B00544:end
% glossary-block:B00545:begin
\needspace{3\baselineskip}
\textbf{gearu} (adj.). ready.\par
% glossary-block:B00545:end
% glossary-block:B00546:begin
\needspace{3\baselineskip}
\textbf{geat} (n. n.); gen. sg. geates; nom. pl. gatu. gate.\par
% glossary-block:B00546:end
% glossary-block:B00547:begin
\needspace{3\baselineskip}
\textbf{gebed} (n. n.). prayer.\par
% glossary-block:B00547:end
% glossary-block:B00548:begin
\needspace{3\baselineskip}
\textbf{gebeorg} (n. n.). protection; defence.\par
% glossary-block:B00548:end
% glossary-block:B00549:begin
\needspace{3\baselineskip}
\textbf{gebēorscipe} (n. m.); gen. sg. gebēorscipes; nom. pl. gebēorscipas. feast; banquet.\par
% glossary-block:B00549:end
% glossary-block:B00550:begin
\needspace{3\baselineskip}
\textbf{gebiddan} (v. strong); pret. 1 sg. gebæd; pret. pl. gebǣdon; past ptcp. gebeden. pray. Supplied wæs gebiddende, rendered prayed in this source. pray. With an accusative reflexive pronoun in the assigned prayer frame.\par
% glossary-block:B00550:end
% glossary-block:B00551:begin
\needspace{3\baselineskip}
\textbf{geblēdsian} — citation/spelling variant: see \textbf{geblētsian}.\par
% glossary-block:B00551:end
% glossary-block:B00552:begin
\needspace{3\baselineskip}
\textbf{geblētsian} (v. weak); pret. 1 sg. geblētsode; past ptcp. geblētsod. bless.\par
% glossary-block:B00552:end
% glossary-block:B00553:begin
\needspace{3\baselineskip}
\textbf{gebrec} (n. n.). crash; noise; din. Supplied accusative singular.\par
% glossary-block:B00553:end
% glossary-block:B00554:begin
\needspace{3\baselineskip}
\textbf{gebringan} (v.). bring. Supplied preterite with plural winds as subject; retain the prefix.\par
% glossary-block:B00554:end
% glossary-block:B00555:begin
\needspace{3\baselineskip}
\textbf{gebrōhton} — pret. ind. pl.; only supplied form: see \textbf{gebringan}.\par
% glossary-block:B00555:end
% glossary-block:B00556:begin
\needspace{3\baselineskip}
\textbf{gebyrd} (n. f.); gen. sg. gebyrde; nom. pl. gebyrda. birth.\par
% glossary-block:B00556:end
% glossary-block:B00557:begin
\needspace{3\baselineskip}
\textbf{gecwēme} (adj.). pleasing; acceptable.\par
% glossary-block:B00557:end
% glossary-block:B00558:begin
\needspace{3\baselineskip}
\textbf{gecwēmlīce} (adv.). pleasantly; agreeably.\par
% glossary-block:B00558:end
% glossary-block:B00559:begin
\needspace{3\baselineskip}
\textbf{gecweðan} (v. strong); pret. 1 sg. gecwæð; pret. pl. gecwǣdon; past ptcp. gecweden; pres. 3 sg. gecwið. speak; say.\par
% glossary-block:B00559:end
% glossary-block:B00560:begin
\needspace{3\baselineskip}
\textbf{gecwið} — pres. ind. 3 sg.: see \textbf{gecweðan}.\par
% glossary-block:B00560:end
% glossary-block:B00561:begin
\needspace{3\baselineskip}
\textbf{ge-cȳðan} (v.). show.\par
% glossary-block:B00561:end
% glossary-block:B00562:begin
\needspace{3\baselineskip}
\textbf{gedǣlan} (v.). divide; separate. Preserve the assigned separation frame.\par
% glossary-block:B00562:end
% glossary-block:B00563:begin
\needspace{3\baselineskip}
\textbf{gedeorf} (n. n.); gen. sg. gedeorfes; nom. pl. gedeorf. labor; toil.\par
% glossary-block:B00563:end
% glossary-block:B00564:begin
\needspace{3\baselineskip}
\textbf{gedēð} — pres. ind. 3 sg.: see \textbf{gedōn}.\par
% glossary-block:B00564:end
% glossary-block:B00565:begin
\needspace{3\baselineskip}
\textbf{gedōn} (v.); past ptcp. gedōn. do; make. Supplied passive fulfilment formula. do; perform. make; arrange. The present form/arrangement of the translation in the source. make; cause to be.\par
% glossary-block:B00565:end
% glossary-block:B00566:begin
\needspace{3\baselineskip}
\textbf{geęndian} (v.). end; come to an end.\par
% glossary-block:B00566:end
% glossary-block:B00567:begin
\needspace{3\baselineskip}
\textbf{geęndung} (n. f.). end; close.\par
% glossary-block:B00567:end
% glossary-block:B00568:begin
\needspace{3\baselineskip}
\textbf{gefeallan} (v.). chance upon. fall in love. Figurative falling of the heart; supplied idiom differs in prefix from befeallan in 42.1.\par
% glossary-block:B00568:end
% glossary-block:B00569:begin
\needspace{3\baselineskip}
\textbf{gefēra} (n. m.); gen. sg. gefēran; nom. pl. gefēran. companion; fellow.\par
% glossary-block:B00569:end
% glossary-block:B00570:begin
\needspace{3\baselineskip}
\textbf{gefęrian} (v.); past ptcp. gefęred. carry; bear. Supplied passive movement; was carried/borne.\par
% glossary-block:B00570:end
% glossary-block:B00571:begin
\needspace{3\baselineskip}
\textbf{gefērscipe} (n. m.); gen. sg. gefērscipes; nom. pl. gefērscipas. companionship.\par
% glossary-block:B00571:end
% glossary-block:B00572:begin
\needspace{3\baselineskip}
\textbf{gefiðred} (adj.). feathered.\par
% glossary-block:B00572:end
% glossary-block:B00573:begin
\needspace{3\baselineskip}
\textbf{geflīt} (n. n.); gen. sg. geflītes; nom. pl. geflīt. dispute; strife.\par
% glossary-block:B00573:end
% glossary-block:B00574:begin
\needspace{3\baselineskip}
\textbf{gefręmman} (v. weak); pret. 1 sg. gefręmede; past ptcp. gefręmed. perform; accomplish. A work is the thing accomplished.\par
% glossary-block:B00574:end
% glossary-block:B00575:begin
\needspace{3\baselineskip}
\textbf{gefultumian} (v.). help; assist.\par
% glossary-block:B00575:end
% glossary-block:B00576:begin
\needspace{3\baselineskip}
\textbf{gefyllan} (v. weak); pret. 1 sg. gefylde; past ptcp. gefylled. finish; accomplish. A work is the thing completed; distinct from fill. fill. fulfil. Accusative prediction fulfilled; passive subjunctive in the assigned purpose construction.\par
% glossary-block:B00576:end
% glossary-block:B00577:begin
\needspace{3\baselineskip}
\textbf{gegaderian} (v. weak); pret. 1 sg. gegaderode; past ptcp. gegaderod. gather. Accusative people or things gathered.\par
% glossary-block:B00577:end
% glossary-block:B00578:begin
\needspace{3\baselineskip}
\textbf{gegaderung} (n. f.); gen. sg. gegaderunge; nom. pl. gegaderunga. gathering; collection; congregation. Congregations of waters in the source.\par
% glossary-block:B00578:end
% glossary-block:B00579:begin
\needspace{3\baselineskip}
\textbf{gegierela} (n. m.); gen. sg. gegierelan; nom. pl. gegierelan. garment; clothing.\par
% glossary-block:B00579:end
% glossary-block:B00580:begin
\needspace{3\baselineskip}
\textbf{geglęngan} (v.); past ptcp. geglęnged. adorn. The supplied passive adornment expression uses mid + dative material.\par
% glossary-block:B00580:end
% glossary-block:B00581:begin
\needspace{3\baselineskip}
\textbf{gehālgian} (v. weak); pret. 1 sg. gehālgode; past ptcp. gehālgod. hallow; consecrate.\par
% glossary-block:B00581:end
% glossary-block:B00582:begin
\needspace{3\baselineskip}
\textbf{gehātan} (v. strong); pret. 1 sg. gehēt; pret. pl. gehēton; past ptcp. gehāten; pres. 3 sg. gehǣtt. promise. Accusative thing promised; a dative recipient may be added.\par
% glossary-block:B00582:end
% glossary-block:B00583:begin
\needspace{3\baselineskip}
\textbf{gehǣtt} — pres. ind. 3 sg.: see \textbf{gehātan}.\par
% glossary-block:B00583:end
% glossary-block:B00584:begin
\needspace{3\baselineskip}
\textbf{gehīeran} (v.). hear. Supplied prefixed third-singular preterite. hear; listen.\par
% glossary-block:B00584:end
% glossary-block:B00585:begin
\needspace{3\baselineskip}
\textbf{gehīersumian} (v.). serve; obey. Supplied source form and government.\par
% glossary-block:B00585:end
% glossary-block:B00586:begin
\needspace{3\baselineskip}
\textbf{gehīersumnes} (n. f.); gen. sg. gehīersumnesse; nom. pl. gehīersumnessa. obedience.\par
% glossary-block:B00586:end
% glossary-block:B00587:begin
\needspace{3\baselineskip}
\textbf{gehwā / gehwæt} (pron.). each one; everyone; everything. Personal gehwā and neuter gehwæt; the early reference status does not by itself establish production.\par
% glossary-block:B00587:end
% glossary-block:B00588:begin
\needspace{3\baselineskip}
\textbf{gehwǣm} — dat. m./f. sg.; dat. n. sg.: see \textbf{gehwā / gehwæt}.\par
% glossary-block:B00588:end
% glossary-block:B00589:begin
\needspace{3\baselineskip}
\textbf{gehwæs} — gen. m./f. sg.; gen. n. sg.: see \textbf{gehwā / gehwæt}.\par
% glossary-block:B00589:end
% glossary-block:B00590:begin
\needspace{3\baselineskip}
\textbf{gehwæt} — nom. n. sg.; acc. n. sg.: see \textbf{gehwā / gehwæt}.\par
% glossary-block:B00590:end
% glossary-block:B00591:begin
\needspace{3\baselineskip}
\textbf{gehwilc} (indefinite adjective). each; every.\par
% glossary-block:B00591:end
% glossary-block:B00592:begin
\needspace{3\baselineskip}
\textbf{gehwone} — acc. m./f. sg.: see \textbf{gehwā / gehwæt}.\par
% glossary-block:B00592:end
% glossary-block:B00593:begin
\needspace{3\baselineskip}
\textbf{gelǣccan} (v.). catch. Supplied caught form.\par
% glossary-block:B00593:end
% glossary-block:B00594:begin
\needspace{3\baselineskip}
\textbf{gelǣdan} (v.). bring; carry. Reference in a source or explanatory note.\par
% glossary-block:B00594:end
% glossary-block:B00595:begin
\needspace{3\baselineskip}
\textbf{gelæhte} — pret. sg.; supplied source form: see \textbf{gelǣccan}.\par
% glossary-block:B00595:end
% glossary-block:B00596:begin
\needspace{3\baselineskip}
\textbf{gelǣran} (v.). teach; instruct. Reference in a source or explanatory note.\par
% glossary-block:B00596:end
% glossary-block:B00597:begin
\needspace{3\baselineskip}
\textbf{gelǣred} (adj.). educated; trained; skilled.\par
% glossary-block:B00597:end
% glossary-block:B00598:begin
\needspace{3\baselineskip}
\textbf{gelēafa} (n. m.); gen. sg. gelēafan; nom. pl. gelēafan. faith.\par
% glossary-block:B00598:end
% glossary-block:B00599:begin
\needspace{3\baselineskip}
\textbf{geleornian} (v.); past ptcp. geleornod. learn; study. learn; master. Learn a skill or performance in the source.\par
% glossary-block:B00599:end
% glossary-block:B00600:begin
\needspace{3\baselineskip}
\textbf{gelīc} (adj.). equal; like. Dative standard of comparison.\par
% glossary-block:B00600:end
% glossary-block:B00601:begin
\needspace{3\baselineskip}
\textbf{gelīca} (n. m.). equal; peer.\par
% glossary-block:B00601:end
% glossary-block:B00602:begin
\needspace{3\baselineskip}
\textbf{gelīce} (adv.). likewise; similarly.\par
% glossary-block:B00602:end
% glossary-block:B00603:begin
\needspace{3\baselineskip}
\textbf{gelīcnes} (n. f.); gen. sg. gelīcnesse; nom. pl. gelīcnessa. likeness.\par
% glossary-block:B00603:end
% glossary-block:B00604:begin
\needspace{3\baselineskip}
\textbf{gelīefan} (v. weak); pret. 1 sg. gelīefde; past ptcp. gelīefed. believe.\par
% glossary-block:B00604:end
% glossary-block:B00605:begin
\needspace{3\baselineskip}
\textbf{gelīefst} (v.; supplied pres. ind. 2 sg.). permit; allow. You permit; you allow. In gif ðū mē gelīefst, mē is the dative person granted permission; speaking is understood.\par
% glossary-block:B00605:end
% glossary-block:B00606:begin
\needspace{3\baselineskip}
\textbf{gelimpan} (v.); pres. 3 sg. gelimpð. happen; befall. Supplied preterite subjunctive in might befall us. happen; come about. Supplied preterite source form. happen; come about.\par
% glossary-block:B00606:end
% glossary-block:B00607:begin
\needspace{3\baselineskip}
\textbf{gelōgian} (v.); pres. 3 sg. gelōgað. place; set.\par
% glossary-block:B00607:end
% glossary-block:B00608:begin
\needspace{3\baselineskip}
\textbf{gelōmlice} (adv.). frequently; regularly.\par
% glossary-block:B00608:end
% glossary-block:B00609:begin
\needspace{3\baselineskip}
\textbf{geman} — pres. ind. 1/3 sg.: see \textbf{gemunan}.\par
% glossary-block:B00609:end
% glossary-block:B00610:begin
\needspace{3\baselineskip}
\textbf{gemang} (prep.). among. Accusative group in the selected sending-among frame.\par
% glossary-block:B00610:end
% glossary-block:B00611:begin
\needspace{3\baselineskip}
\textbf{gemanigfieldan} (v. weak); pret. 1 sg. gemanigfielde; past ptcp. gemanigfield. be multiplied. Assigned passive with agreeing gemanigfield/gemanigfielde; no independent active causative frame is assigned.\par
% glossary-block:B00611:end
% glossary-block:B00612:begin
\needspace{3\baselineskip}
\textbf{gemǣne} (adj.). common; shared.\par
% glossary-block:B00612:end
% glossary-block:B00613:begin
\needspace{3\baselineskip}
\textbf{gemǣre} (n. n.); gen. sg. gemǣres; nom. pl. gemǣru. boundary; border.\par
% glossary-block:B00613:end
% glossary-block:B00614:begin
\needspace{3\baselineskip}
\textbf{gemęngan} (v.). join; mingle. Supplied source form. mix. Mix sound with song in the supplied performance expression.\par
% glossary-block:B00614:end
% glossary-block:B00615:begin
\needspace{3\baselineskip}
\textbf{gemētan} (v. weak); pret. 1 sg. gemētte; past ptcp. gemētt. find. Supplied past participle gemētt in a passive. find; meet. Accusative person or thing found/met.\par
% glossary-block:B00615:end
% glossary-block:B00616:begin
\needspace{3\baselineskip}
\textbf{gemiltsian} (v.). pity; have compassion on. Dative person.\par
% glossary-block:B00616:end
% glossary-block:B00617:begin
\needspace{3\baselineskip}
\textbf{gemiltsiend} (n. m.). one who shows mercy; pitier.\par
% glossary-block:B00617:end
% glossary-block:B00618:begin
\needspace{3\baselineskip}
\textbf{gemunan} (v. preterite-present); pres. 1/3 sg. geman; pres. pl. gemunon; pret. 1 sg. gemunde; past ptcp. gemunen. remember.\par
% glossary-block:B00618:end
% glossary-block:B00619:begin
\needspace{3\baselineskip}
\textbf{genēahhe} (adv.). often; frequently.\par
% glossary-block:B00619:end
% glossary-block:B00620:begin
\needspace{3\baselineskip}
\textbf{genēalǣcan} (v. weak); pret. 1 sg. genēalǣcte; past ptcp. genēalǣced. approach; draw near. Dative person, place or event approached. approach. Supplied source form approached.\par
% glossary-block:B00620:end
% glossary-block:B00621:begin
\needspace{3\baselineskip}
\textbf{genēosian} (v. weak); pret. 1 sg. genēosode; past ptcp. genēosod. visit. Accusative person in the selected prefixed-verb frame.\par
% glossary-block:B00621:end
% glossary-block:B00622:begin
\needspace{3\baselineskip}
\textbf{geniman} (v.). take.\par
% glossary-block:B00622:end
% glossary-block:B00623:begin
\needspace{3\baselineskip}
\textbf{genyhtsumian} (v.). suffice; be sufficient for. Complete supplied construction with dative beneficiary.\par
% glossary-block:B00623:end
% glossary-block:B00624:begin
\needspace{3\baselineskip}
\textbf{geoguð} (n. f.). youth.\par
% glossary-block:B00624:end
% glossary-block:B00625:begin
\needspace{3\baselineskip}
\textbf{geond} (prep.). along; through. With accusative.\par
% glossary-block:B00625:end
% glossary-block:B00626:begin
\needspace{3\baselineskip}
\textbf{geong} (adj.); supplied superl. acc./gen. m. sg. gingstan. young. In supplied superlative gingstan: youngest.\par
% glossary-block:B00626:end
% glossary-block:B00627:begin
\needspace{3\baselineskip}
\textbf{geopenian} (v. weak); pret. 1 sg. geopenode; past ptcp. geopenod. open. Accusative thing opened, including mouth or eyes.\par
% glossary-block:B00627:end
% glossary-block:B00628:begin
\needspace{3\baselineskip}
\textbf{geornfull} (adj.). occupied; intent. Source mental-attention context.\par
% glossary-block:B00628:end
% glossary-block:B00629:begin
\needspace{3\baselineskip}
\textbf{geornlice} (adv.). earnestly; attentively.\par
% glossary-block:B00629:end
% glossary-block:B00630:begin
\needspace{3\baselineskip}
\textbf{geręccan} (v.). explain; render.\par
% glossary-block:B00630:end
% glossary-block:B00631:begin
\needspace{3\baselineskip}
\textbf{gerēfa} (n. m.); gen. sg. gerēfan; nom. pl. gerēfan. steward; reeve.\par
% glossary-block:B00631:end
% glossary-block:B00632:begin
\needspace{3\baselineskip}
\textbf{gereord} (n. n.). meal; banquet.\par
% glossary-block:B00632:end
% glossary-block:B00633:begin
\needspace{3\baselineskip}
\textbf{geręstan} (v. weak); pret. 1 sg. geręste; past ptcp. geręst. rest; repose.\par
% glossary-block:B00633:end
% glossary-block:B00634:begin
\needspace{3\baselineskip}
\textbf{gerīman} (v.). number; count. Exact supplied infinitive gerīman; not an independently assigned full paradigm.\par
% glossary-block:B00634:end
% glossary-block:B00635:begin
\needspace{3\baselineskip}
\textbf{gerīmcræft} (n. m.); gen. sg. gerīmcræftes. arithmetic; computation.\par
% glossary-block:B00635:end
% glossary-block:B00636:begin
\needspace{3\baselineskip}
\textbf{gesǣd} — past ptcp., variant: see \textbf{sęcgan}.\par
% glossary-block:B00636:end
% glossary-block:B00637:begin
\needspace{3\baselineskip}
\textbf{gesceaft} (n. f.); gen. sg. gesceafte; nom. pl. gesceafta. creature; created being or thing.\par
% glossary-block:B00637:end
% glossary-block:B00638:begin
\needspace{3\baselineskip}
\textbf{gescieldan} (v.). shield; protect.\par
% glossary-block:B00638:end
% glossary-block:B00639:begin
\needspace{3\baselineskip}
\textbf{gescieppan} (v. strong); pret. 1 sg. gescēop; pret. pl. gescēopon; past ptcp. gesceapen. create.\par
% glossary-block:B00639:end
% glossary-block:B00640:begin
\needspace{3\baselineskip}
\textbf{gescrȳdan} (v. weak); pret. 1 sg. gescrȳdde; past ptcp. gescrȳd. clothe.\par
% glossary-block:B00640:end
% glossary-block:B00641:begin
\needspace{3\baselineskip}
\textbf{gesęllan} (v.). give.\par
% glossary-block:B00641:end
% glossary-block:B00642:begin
\needspace{3\baselineskip}
\textbf{gesęndan} (v.). send back.\par
% glossary-block:B00642:end
% glossary-block:B00643:begin
\needspace{3\baselineskip}
\textbf{gesēon} (v. strong); pret. 1 sg. geseah; pret. pl. gesāwon; past ptcp. gesewen; pres. 3 sg. gesiehð. see. see. Accusative perceived actor and plain infinitive. see. Absolute use, with no expressed object. see. With a content clause stating the circumstance seen. regard; take notice of. see. With the assigned clause describing how a state of affairs stands.\par
% glossary-block:B00643:end
% glossary-block:B00644:begin
\needspace{3\baselineskip}
\textbf{gesęttan} (v. weak); pret. 1 sg. gesętte; past ptcp. gesętt. set; place. Supplied gesętte. placed; situated. Supplied participial complement of understood us; not the identical-looking finite preterite. set; place; establish. appoint as lord. Supplied frame with person appointed, tō + dative role, and dative people governed.\par
% glossary-block:B00644:end
% glossary-block:B00645:begin
\needspace{3\baselineskip}
\textbf{gesewenlic} (adj.). visible.\par
% glossary-block:B00645:end
% glossary-block:B00646:begin
\needspace{3\baselineskip}
\textbf{gesiehð} — pres. ind. 3 sg.: see \textbf{gesēon}.\par
% glossary-block:B00646:end
% glossary-block:B00647:begin
\needspace{3\baselineskip}
\textbf{gesihð} (n. f.). sight; gaze. Supplied accusative singular.\par
% glossary-block:B00647:end
% glossary-block:B00648:begin
\needspace{3\baselineskip}
\textbf{geslēan} (v.). strike; smite. Supplied past/participial source form. strike. Supplied source form struck.\par
% glossary-block:B00648:end
% glossary-block:B00649:begin
\needspace{3\baselineskip}
\textbf{gesmierwan} (v.); past ptcp. gesmierwed. anoint. Supplied source form anointed.\par
% glossary-block:B00649:end
% glossary-block:B00650:begin
\needspace{3\baselineskip}
\textbf{gesmyltan} (v.). become calm. Supplied preterite of the sea becoming calm.\par
% glossary-block:B00650:end
% glossary-block:B00651:begin
\needspace{3\baselineskip}
\textbf{gestandan} (v.); pres. 3 sg. gestęnt. assail; come upon. Supplied source sense.\par
% glossary-block:B00651:end
% glossary-block:B00652:begin
\needspace{3\baselineskip}
\textbf{gestaðelian} (v.). establish; make steadfast.\par
% glossary-block:B00652:end
% glossary-block:B00653:begin
\needspace{3\baselineskip}
\textbf{gestęnt} — pres. ind. 3 sg.: see \textbf{gestandan}.\par
% glossary-block:B00653:end
% glossary-block:B00654:begin
\needspace{3\baselineskip}
\textbf{gestillan} (v. weak); pret. 1 sg. gestilde; past ptcp. gestilled. quell; settle. Accusative disturbance settled; distinguish intransitive become still. cease; become still. Intransitive in the source storm relative.\par
% glossary-block:B00654:end
% glossary-block:B00655:begin
\needspace{3\baselineskip}
\textbf{gestrangian} (v.). strengthen.\par
% glossary-block:B00655:end
% glossary-block:B00656:begin
\needspace{3\baselineskip}
\textbf{gestrēon} (n. n.); gen. sg. gestrēones; nom. pl. gestrēon. gain; profit.\par
% glossary-block:B00656:end
% glossary-block:B00657:begin
\needspace{3\baselineskip}
\textbf{gesund} (adj.). well; in good health.\par
% glossary-block:B00657:end
% glossary-block:B00658:begin
\needspace{3\baselineskip}
\textbf{geswęngan} (v.); past ptcp. geswęnged. toss; shake. Supplied participial source use.\par
% glossary-block:B00658:end
% glossary-block:B00659:begin
\needspace{3\baselineskip}
\textbf{gesweotolian} (v.); pres. 3 sg. gesweotolað; past ptcp. gesweotolod. reveal; make clear. Supplied source form. make manifest; reveal.\par
% glossary-block:B00659:end
% glossary-block:B00660:begin
\needspace{3\baselineskip}
\textbf{geswīcan} (v. strong); pret. 1 sg. geswāc; pret. pl. geswicon; past ptcp. geswicen. cease. Genitive thing ceased from. fail; cease. Intransitive without a genitive complement in the selected frame.\par
% glossary-block:B00660:end
% glossary-block:B00661:begin
\needspace{3\baselineskip}
\textbf{gesyndig} (adj.). favourable; fair. Of winds.\par
% glossary-block:B00661:end
% glossary-block:B00662:begin
\needspace{3\baselineskip}
\textbf{getǣcan} (v.); past ptcp. getǣht. show; direct. Supplied source form.\par
% glossary-block:B00662:end
% glossary-block:B00663:begin
\needspace{3\baselineskip}
\textbf{getǣht} — past ptcp.; supplied source form: see \textbf{getǣcan}.\par
% glossary-block:B00663:end
% glossary-block:B00664:begin
\needspace{3\baselineskip}
\textbf{getæl} (n. n.). number; company.\par
% glossary-block:B00664:end
% glossary-block:B00665:begin
\needspace{3\baselineskip}
\textbf{getēon} (v.); past ptcp. getogen. train; educate.\par
% glossary-block:B00665:end
% glossary-block:B00666:begin
\needspace{3\baselineskip}
\textbf{getimbrung} (n. f.); gen. sg. getimbrunge; nom. pl. getimbrunga. building; structure built.\par
% glossary-block:B00666:end
% glossary-block:B00667:begin
\needspace{3\baselineskip}
\textbf{geðanc} (n. n.). thought; mind.\par
% glossary-block:B00667:end
% glossary-block:B00668:begin
\needspace{3\baselineskip}
\textbf{geðeahtend} (n. m.); gen. sg. geðeahtendes; nom. pl. geðeahtend. counsellor.\par
% glossary-block:B00668:end
% glossary-block:B00669:begin
\needspace{3\baselineskip}
\textbf{geðęncean} (v.). remember; bear in mind.\par
% glossary-block:B00669:end
% glossary-block:B00670:begin
\needspace{3\baselineskip}
\textbf{geðrǣstan} (v.). afflict; beat. Supplied participial source use.\par
% glossary-block:B00670:end
% glossary-block:B00671:begin
\needspace{3\baselineskip}
\textbf{geðūht} — past ptcp.: see \textbf{ðyncan}.\par
% glossary-block:B00671:end
% glossary-block:B00672:begin
\needspace{3\baselineskip}
\textbf{gewǣgan} (v.); past ptcp. gewǣged. buffet; shake. Supplied participial source use.\par
% glossary-block:B00672:end
% glossary-block:B00673:begin
\needspace{3\baselineskip}
\textbf{geweald} (n. n.); gen. sg. gewealdes; nom. pl. geweald. control; rule; dominion. Ordinary singular; on gewealde is the taught control phrase.\par
% glossary-block:B00673:end
% glossary-block:B00674:begin
\needspace{3\baselineskip}
\textbf{gewęndan} (v. weak); pret. 1 sg. gewęnde; past ptcp. gewęnd. return; turn. translate; turn from one language into another. turn; change back. Supplied source form.\par
% glossary-block:B00674:end
% glossary-block:B00675:begin
\needspace{3\baselineskip}
\textbf{geweorðan} (v.); past ptcp. geworden. come to be; become. come to pass; befall. become; come about.\par
% glossary-block:B00675:end
% glossary-block:B00676:begin
\needspace{3\baselineskip}
\textbf{geweorðian} (v. weak); pret. 1 sg. geweorðode; past ptcp. geweorðod. distinguish; adorn. Accusative person adorned; mid + dative equipment or adornment.\par
% glossary-block:B00676:end
% glossary-block:B00677:begin
\needspace{3\baselineskip}
\textbf{gewiht} (n. n.). weight.\par
% glossary-block:B00677:end
% glossary-block:B00678:begin
\needspace{3\baselineskip}
\textbf{gewilnian} (v. weak); pret. 1 sg. gewilnode; past ptcp. gewilnod. desire; wish.\par
% glossary-block:B00678:end
% glossary-block:B00679:begin
\needspace{3\baselineskip}
\textbf{gewinn} (n. n.). distress; struggle. labour; work. Noun, distinct from labour as a verb.\par
% glossary-block:B00679:end
% glossary-block:B00680:begin
\needspace{3\baselineskip}
\textbf{gewinna} (n. m.). enemy; opponent. Supplied plural source forms.\par
% glossary-block:B00680:end
% glossary-block:B00681:begin
\needspace{3\baselineskip}
\textbf{gewissian} (v.). guide; direct.\par
% glossary-block:B00681:end
% glossary-block:B00682:begin
\needspace{3\baselineskip}
\textbf{gewītan} (v. strong); pret. 1 sg. gewāt; pret. pl. gewiton; past ptcp. gewiten. depart; go away. Intransitive; fram + dative departure point is optional. The assigned frame requires no reflexive. pass away; cease to hold. Of a decree.\par
% glossary-block:B00682:end
% glossary-block:B00683:begin
\needspace{3\baselineskip}
\textbf{gewitan} (v.). find out; learn. Knowledge verb, distinct from gewītan depart.\par
% glossary-block:B00683:end
% glossary-block:B00684:begin
\needspace{3\baselineskip}
\textbf{geworhte} — pret. sg.: see \textbf{gewyrcean}.\par
% glossary-block:B00684:end
% glossary-block:B00685:begin
\needspace{3\baselineskip}
\textbf{gewrit} (n. n.); gen. sg. gewrites; nom. pl. gewritu. letter; written document.\par
% glossary-block:B00685:end
% glossary-block:B00686:begin
\needspace{3\baselineskip}
\textbf{gewuna} (n. m.). custom; habit.\par
% glossary-block:B00686:end
% glossary-block:B00687:begin
\needspace{3\baselineskip}
\textbf{gewunian} (v. weak); pret. 1 sg. gewunode; past ptcp. gewunod. be accustomed; be wont.\par
% glossary-block:B00687:end
% glossary-block:B00688:begin
\needspace{3\baselineskip}
\textbf{gewyrcean} (v. weak); pret. 1 sg. geworhte; past ptcp. geworht. make; build. make; render. Accusative object with an agreeing predicate noun/adjective.\par
% glossary-block:B00688:end
% glossary-block:B00689:begin
\needspace{3\baselineskip}
\textbf{giefan} (v. strong); pret. 1 sg. geaf; pret. pl. gēafon; past ptcp. giefen. give.\par
% glossary-block:B00689:end
% glossary-block:B00690:begin
\needspace{3\baselineskip}
\textbf{giefu} (n. f.). gift. grace; divine favour. gift. Supplied plural gifts.\par
% glossary-block:B00690:end
% glossary-block:B00691:begin
\needspace{3\baselineskip}
\textbf{gieldan} (v. strong); pret. 1 sg. geald; pret. pl. guldon; past ptcp. golden. yield.\par
% glossary-block:B00691:end
% glossary-block:B00692:begin
\needspace{3\baselineskip}
\textbf{gīeman} (v. weak); pret. 1 sg. gīemde; past ptcp. gīemed. rule over. Genitive of what is ruled; no additional over preposition.\par
% glossary-block:B00692:end
% glossary-block:B00693:begin
\needspace{3\baselineskip}
\textbf{giernan} (v. weak); pret. 1 sg. giernde; past ptcp. gierned. desire; request; woo. Genitive desired thing/person in the selected frames, including the maiden wooed.\par
% glossary-block:B00693:end
% glossary-block:B00694:begin
\needspace{3\baselineskip}
\textbf{gierwan} (v. weak); pret. 1 sg. gierede; past ptcp. (ge)gier(w)ed. prepare. Accusative thing prepared.\par
% glossary-block:B00694:end
% glossary-block:B00695:begin
\needspace{3\baselineskip}
\textbf{giesthūs} (n. n.); gen. sg. giesthūses; nom. pl. giesthūs. inn; guest-house.\par
% glossary-block:B00695:end
% glossary-block:B00696:begin
\needspace{3\baselineskip}
\textbf{gīet} (adv.). still; further.\par
% glossary-block:B00696:end
% glossary-block:B00697:begin
\needspace{3\baselineskip}
\textbf{gif} (conj.). if. Mood is determined by the open, tentative or counterfactual condition, not by gif alone.\par
% glossary-block:B00697:end
% glossary-block:B00698:begin
\needspace{3\baselineskip}
\textbf{gimm} (n. m.); gen. sg. gimmes; nom. pl. gimmas. gem; precious stone.\par
% glossary-block:B00698:end
% glossary-block:B00699:begin
\needspace{3\baselineskip}
\textbf{gingstan} — inflected superlative; supplied source form: see \textbf{geong}.\par
% glossary-block:B00699:end
% glossary-block:B00700:begin
\needspace{3\baselineskip}
\textbf{git} — nom. dual: see \textbf{ðū}.\par
% glossary-block:B00700:end
% glossary-block:B00701:begin
\needspace{3\baselineskip}
\textbf{glæd} (adj.). glad. gracious; favourably disposed. With wið + person in the supplied gracious-toward frame.\par
% glossary-block:B00701:end
% glossary-block:B00702:begin
\needspace{3\baselineskip}
\textbf{God} (n. m.); gen. sg. Godes. God. Proper divine name; no productive plural sense is assigned. god; deity. Common deity sense, distinct from the proper divine name.\par
% glossary-block:B00702:end
% glossary-block:B00703:begin
\needspace{3\baselineskip}
\textbf{gōd} (adj.); comp. bętera; superl. bętst. good. kind; benevolent.\par
% glossary-block:B00703:end
% glossary-block:B00704:begin
\needspace{3\baselineskip}
\textbf{gōd} (n. n.); gen. sg. gōdes; nom. pl. gōd. goods; property. Plural lexical noun. good; benefit; kindness. Nominal good in the repayment contrast.\par
% glossary-block:B00704:end
% glossary-block:B00705:begin
\needspace{3\baselineskip}
\textbf{godspell} (n. n.); gen. sg. godspelles; nom. pl. godspell. Gospel.\par
% glossary-block:B00705:end
% glossary-block:B00706:begin
\needspace{3\baselineskip}
\textbf{gōdwębb} (n. n.); gen. sg. gōdwębbes; nom. pl. gōdwębb. purple cloth.\par
% glossary-block:B00706:end
% glossary-block:B00707:begin
\needspace{3\baselineskip}
\textbf{gold} (n. n.); gen. sg. goldes; nom. pl. gold. gold.\par
% glossary-block:B00707:end
% glossary-block:B00708:begin
\needspace{3\baselineskip}
\textbf{goldhord} (n. n.); gen. sg. goldhordes; nom. pl. goldhord. treasure; hoard of gold.\par
% glossary-block:B00708:end
% glossary-block:B00709:begin
\needspace{3\baselineskip}
\textbf{grēne} (adj.). green. Reference in a source or explanatory note.\par
% glossary-block:B00709:end
% glossary-block:B00710:begin
\needspace{3\baselineskip}
\textbf{grētan} (v. weak); pret. 1 sg. grētte; past ptcp. grētt. greet; take leave of. Accusative person/place; greeting and farewell depend on context.\par
% glossary-block:B00710:end
% glossary-block:B00711:begin
\needspace{3\baselineskip}
\textbf{grōwan} (v.). grow. Reference in a source or explanatory note.\par
% glossary-block:B00711:end
% glossary-block:B00712:begin
\needspace{3\baselineskip}
\textbf{grundweall} (n. m.); gen. sg. grundwealles; nom. pl. grundweallas. foundation.\par
% glossary-block:B00712:end
% glossary-block:B00713:begin
\needspace{3\baselineskip}
\textbf{guma} (n. m.). man. poetic/source quotation.\par
% glossary-block:B00713:end
% glossary-block:B00714:begin
\needspace{3\baselineskip}
\textbf{gylt} (n. m.); gen. sg. gyltes; nom. pl. gyltas. debt. Includes the prayer's moral debt.\par
% glossary-block:B00714:end
% glossary-block:B00715:begin
\needspace{3\baselineskip}
\textbf{gyltend} (n. m.); gen. sg. gyltendes; nom. pl. gyltende. debtor.\par
% glossary-block:B00715:end
\pdfbookmark[1]{H}{outline-B00716}
% glossary-block:B00716:begin
\needspace{5\baselineskip}
\subsection*{H}
% glossary-block:B00716:end
% glossary-block:B00717:begin
\textbf{habban} (v. weak); pret. 1 sg. hæfde; past ptcp. (ge)hæfd; pres. 3 sg. hæfð. have; possess. Accusative possession; distinguish the perfect auxiliary. have. Auxiliary with transitive past participle for the taught completed-action/result construction.\par
% glossary-block:B00717:end
% glossary-block:B00718:begin
\needspace{3\baselineskip}
\textbf{hāl} (adj.). hale; well. Supplied greeting. hale; healthy.\par
% glossary-block:B00718:end
% glossary-block:B00719:begin
\needspace{3\baselineskip}
\textbf{hālga} (n. m.); gen. sg. hālgan; nom. pl. hālgan. saint; holy one.\par
% glossary-block:B00719:end
% glossary-block:B00720:begin
\needspace{3\baselineskip}
\textbf{hālig} (adj.). holy. holy ones; saints. Substantival adjective in the source. Before vowel endings hālg-; retain hālig- before -ne, -re and -ra.\par
% glossary-block:B00720:end
% glossary-block:B00721:begin
\needspace{3\baselineskip}
\textbf{hālsian} (v. weak); pret. 1 sg. hālsode; past ptcp. hālsod. entreat; implore. conjure; solemnly charge. Accusative person; ðurh + accusative oath standard; a requested ðæt clause may follow.\par
% glossary-block:B00721:end
% glossary-block:B00722:begin
\needspace{3\baselineskip}
\textbf{hām} (adv.). home; homeward. Directional adverb; distinct from the noun home.\par
% glossary-block:B00722:end
% glossary-block:B00723:begin
\needspace{3\baselineskip}
\textbf{hām} (n. m.); gen. sg. hāmes; nom. pl. hāmas. home; dwelling.\par
% glossary-block:B00723:end
% glossary-block:B00724:begin
\needspace{3\baselineskip}
\textbf{hāmweard} (adverb or particle). homeward.\par
% glossary-block:B00724:end
% glossary-block:B00725:begin
\needspace{3\baselineskip}
\textbf{hand} (n. f.); gen. sg. hande; nom. pl. handa. hand; side. Side in right-/left-hand expressions.\par
% glossary-block:B00725:end
% glossary-block:B00726:begin
\needspace{3\baselineskip}
\textbf{hāt} (adj.). hot.\par
% glossary-block:B00726:end
% glossary-block:B00727:begin
\needspace{3\baselineskip}
\textbf{hātan} (v. strong); pret. 1 sg. hēt; pret. pl. hēton; past ptcp. hāten; pres. 3 sg. hǣtt. call; name. command; order. Accusative actor + plain infinitive for the commanded action; distinct from call/name. be called; be named. Special hātte form in the assigned naming expression.\par
% glossary-block:B00727:end
% glossary-block:B00728:begin
\needspace{3\baselineskip}
\textbf{hatian} (v. weak); pret. 1 sg. hatode; past ptcp. hatod. hate.\par
% glossary-block:B00728:end
% glossary-block:B00729:begin
\needspace{3\baselineskip}
\textbf{hæfde} — pret. sg.: see \textbf{habban}.\par
% glossary-block:B00729:end
% glossary-block:B00730:begin
\needspace{3\baselineskip}
\textbf{hæfð} — pres. ind. 3 sg.: see \textbf{habban}.\par
% glossary-block:B00730:end
% glossary-block:B00731:begin
\needspace{3\baselineskip}
\textbf{hægel} (n. m.). hail.\par
% glossary-block:B00731:end
% glossary-block:B00732:begin
\needspace{3\baselineskip}
\textbf{hǣl} (n. n.). hail; greeting (in hǣl ābēodan, bid someone hail). greeting formula; him is the dative recipient.\par
% glossary-block:B00732:end
% glossary-block:B00733:begin
\needspace{3\baselineskip}
\textbf{hǣlan} (v. weak); pret. 1 sg. hǣlde; past ptcp. hǣled. save.\par
% glossary-block:B00733:end
% glossary-block:B00734:begin
\needspace{3\baselineskip}
\textbf{Hǣlend} (n. m.); gen. sg. Hǣlendes. Saviour; Jesus. Proper title.\par
% glossary-block:B00734:end
% glossary-block:B00735:begin
\needspace{3\baselineskip}
\textbf{hǣlðu} (n.). health. Pronunciation example.\par
% glossary-block:B00735:end
% glossary-block:B00736:begin
\needspace{3\baselineskip}
\textbf{hǣlu} (n. f.); gen. sg. hǣlu. safety; rescue. salvation.\par
% glossary-block:B00736:end
% glossary-block:B00737:begin
\needspace{3\baselineskip}
\textbf{hǣr} (n. n.); gen. sg. hǣres; nom. pl. hǣr. hair; a hair.\par
% glossary-block:B00737:end
% glossary-block:B00738:begin
\needspace{3\baselineskip}
\textbf{hærfest} (n. m.); gen. sg. hærfestes; nom. pl. hærfestas. harvest.\par
% glossary-block:B00738:end
% glossary-block:B00739:begin
\needspace{3\baselineskip}
\textbf{hǣs} (n. f.). command; order.\par
% glossary-block:B00739:end
% glossary-block:B00740:begin
\needspace{3\baselineskip}
\textbf{hǣt} — pres. ind. 3 sg., variant: see \textbf{hātan}.\par
% glossary-block:B00740:end
% glossary-block:B00741:begin
\needspace{3\baselineskip}
\textbf{hǣtt} — pres. ind. 3 sg.: see \textbf{hātan}.\par
% glossary-block:B00741:end
% glossary-block:B00742:begin
\needspace{3\baselineskip}
\textbf{hǣtu} (n. f.); gen. sg. hǣte. heat.\par
% glossary-block:B00742:end
% glossary-block:B00743:begin
\needspace{3\baselineskip}
\textbf{hǣðen} (adj.). heathen people.\par
% glossary-block:B00743:end
% glossary-block:B00744:begin
\needspace{3\baselineskip}
\textbf{hē} (personal pron.). he; she; it; they. Personal pronoun; gender/number and governed case determine the selected form and reference. his; her; its; their. Invariant genitive possessors, unlike agreeing possessive adjectives. their. Source genitive possessor.\par
% glossary-block:B00744:end
% glossary-block:B00745:begin
\needspace{3\baselineskip}
\textbf{hēaclif} (n. n.). high cliff.\par
% glossary-block:B00745:end
% glossary-block:B00746:begin
\needspace{3\baselineskip}
\textbf{hēafod} (n. n.); gen. sg. hēafdes; nom. pl. hēafdu. head.\par
% glossary-block:B00746:end
% glossary-block:B00747:begin
\needspace{3\baselineskip}
\textbf{hēah} (adj.); supplied superl. gen. m. sg. Hēahstan. high. high in rank; exalted. Rank/status sense in the supplied source. the Highest. Divine title. Contracted forms include hēan and hēam; see the Unit 3 pattern.\par
% glossary-block:B00747:end
% glossary-block:B00748:begin
\needspace{3\baselineskip}
\textbf{hēah} (adv.). high. Adverbial height, distinct from an agreeing adjective.\par
% glossary-block:B00748:end
% glossary-block:B00749:begin
\needspace{3\baselineskip}
\textbf{hēahstan} — inflected superlative; supplied source form: see \textbf{hēah}.\par
% glossary-block:B00749:end
% glossary-block:B00750:begin
\needspace{3\baselineskip}
\textbf{healdan} (v. strong); pret. 1 sg. hēold; pret. pl. hēoldon; past ptcp. healden; pres. 3 sg. hielt. hold. keep; observe. Accusative saying or command observed. maintain; sustain.\par
% glossary-block:B00750:end
% glossary-block:B00751:begin
\needspace{3\baselineskip}
\textbf{healf} (adj.). half.\par
% glossary-block:B00751:end
% glossary-block:B00752:begin
\needspace{3\baselineskip}
\textbf{heall} (n. f.). hall.\par
% glossary-block:B00752:end
% glossary-block:B00753:begin
\needspace{3\baselineskip}
\textbf{healt} (adj.). lame; halt. Adjective describing impaired walking; not the command halt/stop.\par
% glossary-block:B00753:end
% glossary-block:B00754:begin
\needspace{3\baselineskip}
\textbf{hēanes} (n. f.); gen. sg. hēanesse; nom. pl. hēanessa. height; highest point. heights; highest places.\par
% glossary-block:B00754:end
% glossary-block:B00755:begin
\needspace{3\baselineskip}
\textbf{hearpe} (n. f.); gen. sg. hearpan; nom. pl. hearpan. harp.\par
% glossary-block:B00755:end
% glossary-block:B00756:begin
\needspace{3\baselineskip}
\textbf{hearpenægl} (n. m.). plectrum; harp pick.\par
% glossary-block:B00756:end
% glossary-block:B00757:begin
\needspace{3\baselineskip}
\textbf{hearpere} (n. m.); gen. sg. hearperes; nom. pl. hearperas. harper.\par
% glossary-block:B00757:end
% glossary-block:B00758:begin
\needspace{3\baselineskip}
\textbf{hearpestręng} (n. m.). harp strings.\par
% glossary-block:B00758:end
% glossary-block:B00759:begin
\needspace{3\baselineskip}
\textbf{hearpian} (v.). harp. Supplied genitive harper/harping construction, with instrument; do not reconstruct it from the infinitive alone. play harp.\par
% glossary-block:B00759:end
% glossary-block:B00760:begin
\needspace{3\baselineskip}
\textbf{hēawan} (v. strong); pret. 1 sg. hēow; pret. pl. hēowon; past ptcp. hēawen; pres. 3 sg. hīewð. hew.\par
% glossary-block:B00760:end
% glossary-block:B00761:begin
\needspace{3\baselineskip}
\textbf{hębban} (v. strong); pret. 1 sg. hōf; pret. pl. hōfon; past ptcp. hafen; pres. 3 sg. hęfð. raise; lift.\par
% glossary-block:B00761:end
% glossary-block:B00762:begin
\needspace{3\baselineskip}
\textbf{hęfig} (adj.). grievous.\par
% glossary-block:B00762:end
% glossary-block:B00763:begin
\needspace{3\baselineskip}
\textbf{hęfignes} (n. f.); gen. sg. hęfignesse; nom. pl. hęfignessa. burden.\par
% glossary-block:B00763:end
% glossary-block:B00764:begin
\needspace{3\baselineskip}
\textbf{hęfð} — pres. ind. 3 sg.: see \textbf{hębban}.\par
% glossary-block:B00764:end
% glossary-block:B00765:begin
\needspace{3\baselineskip}
\textbf{hēht} — pret. sg., older variant: see \textbf{hātan}.\par
% glossary-block:B00765:end
% glossary-block:B00766:begin
\needspace{3\baselineskip}
\textbf{hęll} (n. f.); gen. sg. hęlle; nom. pl. hęlla. hell.\par
% glossary-block:B00766:end
% glossary-block:B00767:begin
\needspace{3\baselineskip}
\textbf{helpan} (v. strong); pret. 1 sg. healp; pret. pl. hulpon; past ptcp. holpen; pres. 3 sg. hilpð. help. Dative person helped.\par
% glossary-block:B00767:end
% glossary-block:B00768:begin
\needspace{3\baselineskip}
\textbf{hēo} — nom. f. sg.; also acc. f. sg. and nom./acc. pl., variant: see \textbf{hē}.\par
% glossary-block:B00768:end
% glossary-block:B00769:begin
\needspace{3\baselineskip}
\textbf{heofon} (n. m.); gen. sg. heofones; nom. pl. heofonas. heaven; sky.\par
% glossary-block:B00769:end
% glossary-block:B00770:begin
\needspace{3\baselineskip}
\textbf{heofone} (n. f.). heaven; sky. Weak feminine lexical alternative in the source.\par
% glossary-block:B00770:end
% glossary-block:B00771:begin
\needspace{3\baselineskip}
\textbf{heofonfȳr} (n. n.). fire from heaven; heavenly fire.\par
% glossary-block:B00771:end
% glossary-block:B00772:begin
\needspace{3\baselineskip}
\textbf{heofonlic} (adj.). heavenly.\par
% glossary-block:B00772:end
% glossary-block:B00773:begin
\needspace{3\baselineskip}
\textbf{heofonlice} (adv.). from heaven; heavenly. Source expression at Cook 143.17.\par
% glossary-block:B00773:end
% glossary-block:B00774:begin
\needspace{3\baselineskip}
\textbf{heonan} (adv.). hence; from here.\par
% glossary-block:B00774:end
% glossary-block:B00775:begin
\needspace{3\baselineskip}
\textbf{heonanforð} (adv.). henceforth.\par
% glossary-block:B00775:end
% glossary-block:B00776:begin
\needspace{3\baselineskip}
\textbf{heora} — gen. pl., variant: see \textbf{hē}.\par
% glossary-block:B00776:end
% glossary-block:B00777:begin
\needspace{3\baselineskip}
\textbf{heorte} (n. f.); gen. sg. heortan; nom. pl. heortan. heart.\par
% glossary-block:B00777:end
% glossary-block:B00778:begin
\needspace{3\baselineskip}
\textbf{hēr} (adv.). here.\par
% glossary-block:B00778:end
% glossary-block:B00779:begin
\needspace{3\baselineskip}
\textbf{hęre} (n. m.); gen. sg. hęres; nom. pl. hęras. army.\par
% glossary-block:B00779:end
% glossary-block:B00780:begin
\needspace{3\baselineskip}
\textbf{hęrgian} (v.). plunder; ravage.\par
% glossary-block:B00780:end
% glossary-block:B00781:begin
\needspace{3\baselineskip}
\textbf{hęrian} (v. weak); pret. 1 sg. hęrede; past ptcp. hęred. praise. Genitive plural present-participle use, of people praising. praise. Supplied source form praised. praise.\par
% glossary-block:B00781:end
% glossary-block:B00782:begin
\needspace{3\baselineskip}
\textbf{hī} — nom./acc. f. sg.; nom./acc. pl., variant: see \textbf{hē}.\par
% glossary-block:B00782:end
% glossary-block:B00783:begin
\needspace{3\baselineskip}
\textbf{hider} (adv.). hither; here. Motion toward the speaker, distinct from location hēr.\par
% glossary-block:B00783:end
% glossary-block:B00784:begin
\needspace{3\baselineskip}
\textbf{hīe} — acc. f sg.; nom. pl.; acc. pl.: see \textbf{hē}.\par
% glossary-block:B00784:end
% glossary-block:B00785:begin
\needspace{3\baselineskip}
\textbf{hielt} — pres. ind. 3 sg.: see \textbf{healdan}.\par
% glossary-block:B00785:end
% glossary-block:B00786:begin
\needspace{3\baselineskip}
\textbf{hīenan} (v.). oppress; humble.\par
% glossary-block:B00786:end
% glossary-block:B00787:begin
\needspace{3\baselineskip}
\textbf{hiene} — acc. m. sg., variant: see \textbf{hē}.\par
% glossary-block:B00787:end
% glossary-block:B00788:begin
\needspace{3\baselineskip}
\textbf{hīera} — gen. pl., source spelling: see \textbf{hē}.\par
% glossary-block:B00788:end
% glossary-block:B00789:begin
\needspace{3\baselineskip}
\textbf{hiera} — gen. pl.: see \textbf{hē}.\par
% glossary-block:B00789:end
% glossary-block:B00790:begin
\needspace{3\baselineskip}
\textbf{hīeran} (v. weak); pret. 1 sg. hīerde; past ptcp. (ge)hīered. hear. hear an action. Accusative perceived actor plus plain infinitive in the taught perception construction. obey. Dative person obeyed.\par
% glossary-block:B00790:end
% glossary-block:B00791:begin
\needspace{3\baselineskip}
\textbf{hiere} — gen. f sg.; dat. f sg.: see \textbf{hē}.\par
% glossary-block:B00791:end
% glossary-block:B00792:begin
\needspace{3\baselineskip}
\textbf{hīernes} (n. f.). obedience.\par
% glossary-block:B00792:end
% glossary-block:B00793:begin
\needspace{3\baselineskip}
\textbf{Hierusalem} (proper name). Hierusalem.\par
% glossary-block:B00793:end
% glossary-block:B00794:begin
\needspace{3\baselineskip}
\textbf{hīewð} — pres. ind. 3 sg.: see \textbf{hēawan}.\par
% glossary-block:B00794:end
% glossary-block:B00795:begin
\needspace{3\baselineskip}
\textbf{hilpð} — pres. ind. 3 sg.: see \textbf{helpan}.\par
% glossary-block:B00795:end
% glossary-block:B00796:begin
\needspace{3\baselineskip}
\textbf{him} — dat. m sg.; dat. n sg.; dat. pl.: see \textbf{hē}.\par
% glossary-block:B00796:end
% glossary-block:B00797:begin
\needspace{3\baselineskip}
\textbf{hine} — acc. m sg.: see \textbf{hē}.\par
% glossary-block:B00797:end
% glossary-block:B00798:begin
\needspace{3\baselineskip}
\textbf{hīo} — nom. f. sg., variant: see \textbf{hē}.\par
% glossary-block:B00798:end
% glossary-block:B00799:begin
\needspace{3\baselineskip}
\textbf{hira} — gen. pl., variant: see \textbf{hē}.\par
% glossary-block:B00799:end
% glossary-block:B00800:begin
\needspace{3\baselineskip}
\textbf{hire} — gen./dat. f. sg., variant: see \textbf{hē}.\par
% glossary-block:B00800:end
% glossary-block:B00801:begin
\needspace{3\baselineskip}
\textbf{his} — gen. m sg.; gen. n sg.: see \textbf{hē}.\par
% glossary-block:B00801:end
% glossary-block:B00802:begin
\needspace{3\baselineskip}
\textbf{hit} — nom. n sg.; acc. n sg.: see \textbf{hē}.\par
% glossary-block:B00802:end
% glossary-block:B00803:begin
\needspace{3\baselineskip}
\textbf{hīw} (n. n.). form; kind. Supplied dative neuter plural.\par
% glossary-block:B00803:end
% glossary-block:B00804:begin
\needspace{3\baselineskip}
\textbf{hlāf} (n. m.); gen. sg. hlāfes. bread; loaf.\par
% glossary-block:B00804:end
% glossary-block:B00805:begin
\needspace{3\baselineskip}
\textbf{hlāford} (n. m.); gen. sg. hlāfordes; nom. pl. hlāfordas. lord; master.\par
% glossary-block:B00805:end
% glossary-block:B00806:begin
\needspace{3\baselineskip}
\textbf{hlǣfdige} (n. f.); gen. sg. hlǣfdigan; nom. pl. hlǣfdigan. lady.\par
% glossary-block:B00806:end
% glossary-block:B00807:begin
\needspace{3\baselineskip}
\textbf{hlæst} (n. n.); gen. sg. hlæstes; nom. pl. hlæst. wares; merchandise; cargo. Selected plural source use.\par
% glossary-block:B00807:end
% glossary-block:B00808:begin
\needspace{3\baselineskip}
\textbf{hlīfian} (v.); pres. 3 sg. hlīfað. tower. Verb meaning stand/tower high in the source; distinct from noun tower.\par
% glossary-block:B00808:end
% glossary-block:B00809:begin
\needspace{3\baselineskip}
\textbf{hlīsa} (n. m.); gen. sg. hlīsan; nom. pl. hlīsan. fame. report; account heard.\par
% glossary-block:B00809:end
% glossary-block:B00810:begin
\needspace{3\baselineskip}
\textbf{hlūd} (adj.). loud.\par
% glossary-block:B00810:end
% glossary-block:B00811:begin
\needspace{3\baselineskip}
\textbf{hlystan} (v.). listen to; heed. Genitive person or thing listened to; hear is the contextual prompt wording.\par
% glossary-block:B00811:end
% glossary-block:B00812:begin
\needspace{3\baselineskip}
\textbf{hof} (n. n.); gen. sg. hofes; nom. pl. hofu. dwelling.\par
% glossary-block:B00812:end
% glossary-block:B00813:begin
\needspace{3\baselineskip}
\textbf{hogað} — pres. ind. 3 sg., variant: see \textbf{hycgan}.\par
% glossary-block:B00813:end
% glossary-block:B00814:begin
\needspace{3\baselineskip}
\textbf{hogde} — pret. sg., variant: see \textbf{hycgan}.\par
% glossary-block:B00814:end
% glossary-block:B00815:begin
\needspace{3\baselineskip}
\textbf{hogode} — pret. sg.: see \textbf{hycgan}.\par
% glossary-block:B00815:end
% glossary-block:B00816:begin
\needspace{3\baselineskip}
\textbf{hol} (n. n.); gen. sg. holes; nom. pl. holu. hole.\par
% glossary-block:B00816:end
% glossary-block:B00817:begin
\needspace{3\baselineskip}
\textbf{hōlunga} (adv.). in vain.\par
% glossary-block:B00817:end
% glossary-block:B00818:begin
\needspace{3\baselineskip}
\textbf{hord} (n. n/m.). hoard; treasure.\par
% glossary-block:B00818:end
% glossary-block:B00819:begin
\needspace{3\baselineskip}
\textbf{horg} (adj.). dirty; filthy.\par
% glossary-block:B00819:end
% glossary-block:B00820:begin
\needspace{3\baselineskip}
\textbf{horn} (n. m.); gen. sg. hornes; nom. pl. hornas. horn; strength; dignity. Strength/dignity is the figurative force in 37.14, not a replacement for the literal source horn.\par
% glossary-block:B00820:end
% glossary-block:B00821:begin
\needspace{3\baselineskip}
\textbf{hors} (n. n.); gen. sg. horses; nom. pl. hors. horse.\par
% glossary-block:B00821:end
% glossary-block:B00822:begin
\needspace{3\baselineskip}
\textbf{hraðe} (adv.). quickly.\par
% glossary-block:B00822:end
% glossary-block:B00823:begin
\needspace{3\baselineskip}
\textbf{hrædlīce} (adv.). quickly.\par
% glossary-block:B00823:end
% glossary-block:B00824:begin
\needspace{3\baselineskip}
\textbf{hrædnes} (n. f.). speed; quickness.\par
% glossary-block:B00824:end
% glossary-block:B00825:begin
\needspace{3\baselineskip}
\textbf{hrēod} (n. n.); gen. sg. hrēodes; nom. pl. hrēod. reed.\par
% glossary-block:B00825:end
% glossary-block:B00826:begin
\needspace{3\baselineskip}
\textbf{hrēohnes} (n.). tempest; storm.\par
% glossary-block:B00826:end
% glossary-block:B00827:begin
\needspace{3\baselineskip}
\textbf{hrēosan} (v.); pret. pl. hruron. fall. Supplied present participle falling.\par
% glossary-block:B00827:end
% glossary-block:B00828:begin
\needspace{3\baselineskip}
\textbf{hrōf} (n. m.); gen. sg. hrōfes; nom. pl. hrōfas. roof.\par
% glossary-block:B00828:end
% glossary-block:B00829:begin
\needspace{3\baselineskip}
\textbf{hū} (adv.). how?.\par
% glossary-block:B00829:end
% glossary-block:B00830:begin
\needspace{3\baselineskip}
\textbf{hūlic} (interrogative adjective). of what sort?.\par
% glossary-block:B00830:end
% glossary-block:B00831:begin
\needspace{3\baselineskip}
\textbf{hund} (cardinal). hundred. Only the complete supplied hundred-sheep expression at its first appearance; later numerical frames remain locally specified.\par
% glossary-block:B00831:end
% glossary-block:B00832:begin
\needspace{3\baselineskip}
\textbf{hund} (n. m.); gen. sg. hundes; nom. pl. hundas. dog.\par
% glossary-block:B00832:end
% glossary-block:B00833:begin
\needspace{3\baselineskip}
\textbf{hundnigontig} (numeral). ninety.\par
% glossary-block:B00833:end
% glossary-block:B00834:begin
\needspace{3\baselineskip}
\textbf{hundtēontig} (numeral). a hundredfold. A hundred in return, in the supplied recompense expression. hundred. Within the complete supplied 144,000 expression; do not infer an independent high-number system.\par
% glossary-block:B00834:end
% glossary-block:B00835:begin
\needspace{3\baselineskip}
\textbf{hunger} (n. m.); gen. sg. hungres. famine; hunger.\par
% glossary-block:B00835:end
% glossary-block:B00836:begin
\needspace{3\baselineskip}
\textbf{hungrig} (adj.). hungry. hungry people. Substantival adjective.\par
% glossary-block:B00836:end
% glossary-block:B00837:begin
\needspace{3\baselineskip}
\textbf{hunig} (n. n.); gen. sg. huniges. honey.\par
% glossary-block:B00837:end
% glossary-block:B00838:begin
\needspace{3\baselineskip}
\textbf{hūs} (n. n.); gen. sg. hūses; nom. pl. hūs. house.\par
% glossary-block:B00838:end
% glossary-block:B00839:begin
\needspace{3\baselineskip}
\textbf{hwā, hwæt} (interrogative/indefinite pron.). who? what?. Case follows role/government; distinguish personal and neuter reference. anyone. what reason. Instrumental hwon in the supplied why expressions.\par
% glossary-block:B00839:end
% glossary-block:B00840:begin
\needspace{3\baselineskip}
\textbf{hwām} — dat. m./f. sg.; dat. n. sg.: see \textbf{hwā, hwæt}.\par
% glossary-block:B00840:end
% glossary-block:B00841:begin
\needspace{3\baselineskip}
\textbf{hwanon} (adv.). whence; from where.\par
% glossary-block:B00841:end
% glossary-block:B00842:begin
\needspace{3\baselineskip}
\textbf{hwæl} (n. m.); gen. sg. hwæles; nom. pl. hwalas. whale.\par
% glossary-block:B00842:end
% glossary-block:B00843:begin
\needspace{3\baselineskip}
\textbf{hwǣm} — dat. m./f. sg.; dat. n. sg.: see \textbf{hwā, hwæt}.\par
% glossary-block:B00843:end
% glossary-block:B00844:begin
\needspace{3\baselineskip}
\textbf{hwǣr} (adv.). where.\par
% glossary-block:B00844:end
% glossary-block:B00845:begin
\needspace{3\baselineskip}
\textbf{hwæs} — gen. m./f. sg.; gen. n. sg.: see \textbf{hwā, hwæt}.\par
% glossary-block:B00845:end
% glossary-block:B00846:begin
\needspace{3\baselineskip}
\textbf{hwæt} — nom. n. sg.; acc. n. sg.: see \textbf{hwā, hwæt}.\par
% glossary-block:B00846:end
% glossary-block:B00847:begin
\needspace{3\baselineskip}
\textbf{hwǣte} (n. m.); gen. sg. hwǣtes. wheat.\par
% glossary-block:B00847:end
% glossary-block:B00848:begin
\needspace{3\baselineskip}
\textbf{hwæthwega} (indefinite pronoun/substantive). something. Supplied local indefinite form; no unrestricted new paradigm.\par
% glossary-block:B00848:end
% glossary-block:B00849:begin
\needspace{3\baselineskip}
\textbf{hwæthwugu} (indefinite). something. Supplied hwæthwugu swilces.\par
% glossary-block:B00849:end
% glossary-block:B00850:begin
\needspace{3\baselineskip}
\textbf{hwæðer} (interrogative pronoun/adjective). which of two?. whichever of two. In swā hwæðer swā; the supplied attraction example has a special case arrangement.\par
% glossary-block:B00850:end
% glossary-block:B00851:begin
\needspace{3\baselineskip}
\textbf{hwæðer} (conj.). whether.\par
% glossary-block:B00851:end
% glossary-block:B00852:begin
\needspace{3\baselineskip}
\textbf{hwelc} — citation/spelling variant: see \textbf{hwilc}.\par
% glossary-block:B00852:end
% glossary-block:B00853:begin
\needspace{3\baselineskip}
\textbf{hweorfan} (v.). return; turn back. Supplied source preterite.\par
% glossary-block:B00853:end
% glossary-block:B00854:begin
\needspace{3\baselineskip}
\textbf{hwider} (adv.). whither; to where.\par
% glossary-block:B00854:end
% glossary-block:B00855:begin
\needspace{3\baselineskip}
\textbf{hwilc} (interrogative/indefinite adj. or pron.). which? of what kind?. which. Selection from a group expressed in genitive. any; some. Indefinite adjective, distinct from an interrogative.\par
% glossary-block:B00855:end
% glossary-block:B00856:begin
\needspace{3\baselineskip}
\textbf{hwīlum} (adv.). sometimes. a while ago. Only in supplied hwīlum ǣr; earlier hwīlum sometimes is not enough.\par
% glossary-block:B00856:end
% glossary-block:B00857:begin
\needspace{3\baselineskip}
\textbf{hwīt} (adj.). white.\par
% glossary-block:B00857:end
% glossary-block:B00858:begin
\needspace{3\baselineskip}
\textbf{hwon} — instr. n. sg.: see \textbf{hwā, hwæt}.\par
% glossary-block:B00858:end
% glossary-block:B00859:begin
\needspace{3\baselineskip}
\textbf{hwone} — acc. m./f. sg.: see \textbf{hwā, hwæt}.\par
% glossary-block:B00859:end
% glossary-block:B00860:begin
\needspace{3\baselineskip}
\textbf{hwonne} (adv.). when?. Direct interrogative.\par
% glossary-block:B00860:end
% glossary-block:B00861:begin
\needspace{3\baselineskip}
\textbf{hwȳ} (adv.). why?. Also instr. n. sg.: see \textbf{hwā, hwæt}.\par
% glossary-block:B00861:end
% glossary-block:B00862:begin
\needspace{3\baselineskip}
\textbf{hwylc} — citation/spelling variant: see \textbf{hwilc}.\par
% glossary-block:B00862:end
% glossary-block:B00863:begin
\needspace{3\baselineskip}
\textbf{hycgan} (v. weak); pret. 1 sg. hogode; past ptcp. (ge)hogod; pres. 3 sg. hygð. think.\par
% glossary-block:B00863:end
% glossary-block:B00864:begin
\needspace{3\baselineskip}
\textbf{hȳdan} (v. weak); pret. 1 sg. hȳdde; past ptcp. hȳded. hide; conceal. Accusative thing hidden; wesan forms the assigned passive.\par
% glossary-block:B00864:end
% glossary-block:B00865:begin
\needspace{3\baselineskip}
\textbf{hygð} — pres. ind. 3 sg.: see \textbf{hycgan}.\par
% glossary-block:B00865:end
% glossary-block:B00866:begin
\needspace{3\baselineskip}
\textbf{hyll} (n. m.); gen. sg. hylles; nom. pl. hyllas. hill.\par
% glossary-block:B00866:end
% glossary-block:B00867:begin
\needspace{3\baselineskip}
\textbf{hyngran} (v. impersonal); pres. 3 sg. hyngreð; pret. 3 sg. hyngrede; past ptcp. (ge)hyngred. be hungry. Impersonal construction with accusative experiencer. hungry; hungering. Present participle used substantivally in the assigned frame. Impersonal use: the course supplies third-singular finite forms.\par
% glossary-block:B00867:end
% glossary-block:B00868:begin
\needspace{3\baselineskip}
\textbf{hȳran} — citation/spelling variant: see \textbf{hīeran}.\par
% glossary-block:B00868:end
% glossary-block:B00869:begin
\needspace{3\baselineskip}
\textbf{hȳrian} (v. weak); pret. 1 sg. hȳrode; past ptcp. hȳrod. hire.\par
% glossary-block:B00869:end
\pdfbookmark[1]{I}{outline-B00870}
% glossary-block:B00870:begin
\needspace{5\baselineskip}
\subsection*{I}
% glossary-block:B00870:end
% glossary-block:B00871:begin
\textbf{Iacob} (proper name). Jacob.\par
% glossary-block:B00871:end
% glossary-block:B00872:begin
\needspace{3\baselineskip}
\textbf{ic} (personal pron.). I; me; we; us. Singular, dual and plural personal forms; select governed case. of me; of us. Personal genitive citation forms identifying person/number; this record alone does not teach an attributive possessive use. our. Personal genitive plural used attributively in the exact ūre hǣle expression. of us two. Personal genitive dual after a selection word.\par
% glossary-block:B00872:end
% glossary-block:B00873:begin
\needspace{3\baselineskip}
\textbf{īdel} (adj.). empty; void. idle. Distinct from empty/void.\par
% glossary-block:B00873:end
% glossary-block:B00874:begin
\needspace{3\baselineskip}
\textbf{īdelnes} (n. f.); gen. sg. īdelnesse; nom. pl. īdelnessa. idleness.\par
% glossary-block:B00874:end
% glossary-block:B00875:begin
\needspace{3\baselineskip}
\textbf{īecan} (v. weak); pret. 1 sg. īecte; past ptcp. (ge)īeced. increase.\par
% glossary-block:B00875:end
% glossary-block:B00876:begin
\needspace{3\baselineskip}
\textbf{ieldest} (adj.). oldest. eldest; oldest. Substantival person/eldest brother in this source context. Superlative of eald.\par
% glossary-block:B00876:end
% glossary-block:B00877:begin
\needspace{3\baselineskip}
\textbf{ieldra} (adj.). older. Comparative of eald.\par
% glossary-block:B00877:end
% glossary-block:B00878:begin
\needspace{3\baselineskip}
\textbf{ieldu} (n. f.). age; old age.\par
% glossary-block:B00878:end
% glossary-block:B00879:begin
\needspace{3\baselineskip}
\textbf{ierman} (v.). afflict.\par
% glossary-block:B00879:end
% glossary-block:B00880:begin
\needspace{3\baselineskip}
\textbf{iernan} (v. strong); pret. 1 sg. arn; pret. pl. urnon; past ptcp. urnen. run. Intransitive; geond + accusative path. A perception construction may make its actor accusative.\par
% glossary-block:B00880:end
% glossary-block:B00881:begin
\needspace{3\baselineskip}
\textbf{ierð} (n. f.). crops; produce of the earth.\par
% glossary-block:B00881:end
% glossary-block:B00882:begin
\needspace{3\baselineskip}
\textbf{ierðling} (n. m.); gen. sg. ierðlinges; nom. pl. ierðlingas. farmer; plowman.\par
% glossary-block:B00882:end
% glossary-block:B00883:begin
\needspace{3\baselineskip}
\textbf{īgland} (n. n.); gen. sg. īglandes; nom. pl. īgland. island.\par
% glossary-block:B00883:end
% glossary-block:B00884:begin
\needspace{3\baselineskip}
\textbf{ilca} (demonstrative adjective). same.\par
% glossary-block:B00884:end
% glossary-block:B00885:begin
\needspace{3\baselineskip}
\textbf{in} (prep.). in; on. With dative in the exact locations/circumstances. in; inward. Adverbial direction. toward; onto. Supplied phrase directing the mind toward these things.\par
% glossary-block:B00885:end
% glossary-block:B00886:begin
\needspace{3\baselineskip}
\textbf{inc} — dat. dual; acc. dual: see \textbf{ðū}.\par
% glossary-block:B00886:end
% glossary-block:B00887:begin
\needspace{3\baselineskip}
\textbf{incer} (possessive adjective). your; yours. Two possessors; agreeing possessive. Also gen. dual: see \textbf{ðū}.\par
% glossary-block:B00887:end
% glossary-block:B00888:begin
\needspace{3\baselineskip}
\textbf{inēode} — pret. sg.: see \textbf{ingān}.\par
% glossary-block:B00888:end
% glossary-block:B00889:begin
\needspace{3\baselineskip}
\textbf{ingān} (v. irregular); pres. 1 sg. ingā; pres. 3 sg. ingǣð; pres. pl. ingāð; pret. 1 sg. inēode; pret. pl. inēodon; past ptcp. ingān. enter; go in. enter; go in. Supplied past source form.\par
% glossary-block:B00889:end
% glossary-block:B00890:begin
\needspace{3\baselineskip}
\textbf{ingangan} (v.). enter; go in. Supplied present-participle phrase.\par
% glossary-block:B00890:end
% glossary-block:B00891:begin
\needspace{3\baselineskip}
\textbf{ingǣð} — pres. ind. 3 sg.: see \textbf{ingān}.\par
% glossary-block:B00891:end
% glossary-block:B00892:begin
\needspace{3\baselineskip}
\textbf{innan} (adv.). within; inside.\par
% glossary-block:B00892:end
% glossary-block:B00893:begin
\needspace{3\baselineskip}
\textbf{intinga} (n. m.). cause; sake (in for … intingan). Within supplied for intingan expressions.\par
% glossary-block:B00893:end
% glossary-block:B00894:begin
\needspace{3\baselineskip}
\textbf{in-tō} (prep.). into; destination expressed with dative after intō.\par
% glossary-block:B00894:end
% glossary-block:B00895:begin
\needspace{3\baselineskip}
\textbf{inweardlīce} (adv.). fervently; wholeheartedly.\par
% glossary-block:B00895:end
% glossary-block:B00896:begin
\needspace{3\baselineskip}
\textbf{Iosep} (proper name). Joseph.\par
% glossary-block:B00896:end
% glossary-block:B00897:begin
\needspace{3\baselineskip}
\textbf{īs} (n. n.). ice.\par
% glossary-block:B00897:end
% glossary-block:B00898:begin
\needspace{3\baselineskip}
\textbf{is} — pres. ind. 3 sg.: see \textbf{wesan}.\par
% glossary-block:B00898:end
% glossary-block:B00899:begin
\needspace{3\baselineskip}
\textbf{īsern} (n. n.); gen. sg. īsernes; nom. pl. īsern. iron. Metal noun.\par
% glossary-block:B00899:end
% glossary-block:B00900:begin
\needspace{3\baselineskip}
\textbf{īsern} (adj.). iron; made of iron. Adjective of material.\par
% glossary-block:B00900:end
% glossary-block:B00901:begin
\needspace{3\baselineskip}
\textbf{Israel} (proper name); gen. sg. Israeles. Israel.\par
% glossary-block:B00901:end
% glossary-block:B00902:begin
\needspace{3\baselineskip}
\textbf{Israela} (proper name). Israelites. Supplied genitive plural name form. Israela is the supplied genitive plural name-form, of the Israelites.\par
% glossary-block:B00902:end
% glossary-block:B00903:begin
\needspace{3\baselineskip}
\textbf{it} — pres. ind. 3 sg., variant: see \textbf{etan}.\par
% glossary-block:B00903:end
% glossary-block:B00904:begin
\needspace{3\baselineskip}
\textbf{itt} — pres. ind. 3 sg.: see \textbf{etan}.\par
% glossary-block:B00904:end
% glossary-block:B00905:begin
\needspace{3\baselineskip}
\textbf{Iudas} (proper name). Judah.\par
% glossary-block:B00905:end
% glossary-block:B00906:begin
\needspace{3\baselineskip}
\textbf{Iudea} (proper name). Jews. Supplied genitive plural name form. Iudea is the supplied genitive plural name-form, of the Jews.\par
% glossary-block:B00906:end
\pdfbookmark[1]{J}{outline-B00907}
% glossary-block:B00907:begin
\needspace{5\baselineskip}
\subsection*{J}
% glossary-block:B00907:end
% glossary-block:B00908:begin
\textbf{Johannes} (proper name); gen. sg. Johannes. John.\par
% glossary-block:B00908:end
\pdfbookmark[1]{L}{outline-B00909}
% glossary-block:B00909:begin
\needspace{5\baselineskip}
\subsection*{L}
% glossary-block:B00909:end
% glossary-block:B00910:begin
\textbf{lāc} (n. n.); gen. sg. lāces; nom. pl. lāc. gift; present.\par
% glossary-block:B00910:end
% glossary-block:B00911:begin
\needspace{3\baselineskip}
\textbf{land} (n. n.); gen. sg. landes; nom. pl. land. land; country.\par
% glossary-block:B00911:end
% glossary-block:B00912:begin
\needspace{3\baselineskip}
\textbf{landgemǣre} (n. n.). territory; borders.\par
% glossary-block:B00912:end
% glossary-block:B00913:begin
\needspace{3\baselineskip}
\textbf{lang} (adj.). long. Adjective agreeing with the time interval fæc.\par
% glossary-block:B00913:end
% glossary-block:B00914:begin
\needspace{3\baselineskip}
\textbf{lange} (adv.). long. In supplied swā lange swā; its biblical inasmuch-as force is explained locally.\par
% glossary-block:B00914:end
% glossary-block:B00915:begin
\needspace{3\baselineskip}
\textbf{lār} (n. f.); gen. sg. lāre. learning; study.\par
% glossary-block:B00915:end
% glossary-block:B00916:begin
\needspace{3\baselineskip}
\textbf{lārēow} (n. m.); gen. sg. lārēowes; nom. pl. lārēowas. master; teacher.\par
% glossary-block:B00916:end
% glossary-block:B00917:begin
\needspace{3\baselineskip}
\textbf{Lazarus} (proper name). Lazarus. The parable's poor man; supplied plural denotes people like him.\par
% glossary-block:B00917:end
% glossary-block:B00918:begin
\needspace{3\baselineskip}
\textbf{lǣdan} (v. weak); pret. 1 sg. lǣdde; past ptcp. (ge)lǣd(ed). lead. lead; bring; take. A person is the object in the assigned frame. carry; take. A person is conveyed in this supplied frame. take; carry along.\par
% glossary-block:B00918:end
% glossary-block:B00919:begin
\needspace{3\baselineskip}
\textbf{Læden} (n. n.). Latin. The language.\par
% glossary-block:B00919:end
% glossary-block:B00920:begin
\needspace{3\baselineskip}
\textbf{læfel} (n. m.); gen. sg. læfeles; nom. pl. læfelas. cup.\par
% glossary-block:B00920:end
% glossary-block:B00921:begin
\needspace{3\baselineskip}
\textbf{lǣran} (v. weak); pret. 1 sg. lǣrde; past ptcp. lǣred. advise; instruct. teach; instruct.\par
% glossary-block:B00921:end
% glossary-block:B00922:begin
\needspace{3\baselineskip}
\textbf{lǣs} (degree word). less; a smaller amount.\par
% glossary-block:B00922:end
% glossary-block:B00923:begin
\needspace{3\baselineskip}
\textbf{lǣssa} (adj.). smaller. Comparative of lȳtel.\par
% glossary-block:B00923:end
% glossary-block:B00924:begin
\needspace{3\baselineskip}
\textbf{lǣst} (adj.). smallest. least. In rank or importance. Superlative of lȳtel.\par
% glossary-block:B00924:end
% glossary-block:B00925:begin
\needspace{3\baselineskip}
\textbf{lǣtan} (v. strong); pret. 1 sg. lēt; pret. pl. lēton; past ptcp. lǣten. let.\par
% glossary-block:B00925:end
% glossary-block:B00926:begin
\needspace{3\baselineskip}
\textbf{lēad} (n. n.); gen. sg. lēades; nom. pl. lēad. lead. Metal noun, not the verb lead.\par
% glossary-block:B00926:end
% glossary-block:B00927:begin
\needspace{3\baselineskip}
\textbf{lēaf} (n. f.). permission.\par
% glossary-block:B00927:end
% glossary-block:B00928:begin
\needspace{3\baselineskip}
\textbf{lęcgean} (v. weak); pret. 1 sg. lęgde; past ptcp. (ge)lęgd. lay; place. Supplied source form. lay; place.\par
% glossary-block:B00928:end
% glossary-block:B00929:begin
\needspace{3\baselineskip}
\textbf{lēde} — pret. sg., contracted variant: see \textbf{lęcgean}.\par
% glossary-block:B00929:end
% glossary-block:B00930:begin
\needspace{3\baselineskip}
\textbf{lęng} (adverb or particle). longer.\par
% glossary-block:B00930:end
% glossary-block:B00931:begin
\needspace{3\baselineskip}
\textbf{lēo} (n.). lion (Latin loan-form).\par
% glossary-block:B00931:end
% glossary-block:B00932:begin
\needspace{3\baselineskip}
\textbf{lēof} (adj.); comp. lēofra; superl. lēofost. dear. sir. Supplied respectful address lēof.\par
% glossary-block:B00932:end
% glossary-block:B00933:begin
\needspace{3\baselineskip}
\textbf{leofað} — pres. ind. 3 sg.: see \textbf{libban}.\par
% glossary-block:B00933:end
% glossary-block:B00934:begin
\needspace{3\baselineskip}
\textbf{lēofost} (adj.). dearest. Superlative of lēof.\par
% glossary-block:B00934:end
% glossary-block:B00935:begin
\needspace{3\baselineskip}
\textbf{lēofra} (adj.). dearer. preferable; dearer. Dative experiencer + is lēofre + infinitive in the taught preference frame. Comparative of lēof.\par
% glossary-block:B00935:end
% glossary-block:B00936:begin
\needspace{3\baselineskip}
\textbf{lēoht} (n. n.); gen. sg. lēohtes; nom. pl. lēoht. light. light; luminary. Individual countable light, distinct from mass light.\par
% glossary-block:B00936:end
% glossary-block:B00937:begin
\needspace{3\baselineskip}
\textbf{leornian} (v. weak); pret. 1 sg. leornode; past ptcp. leornod. learn. Accusative thing learned; æt + dative person/source.\par
% glossary-block:B00937:end
% glossary-block:B00938:begin
\needspace{3\baselineskip}
\textbf{leornung} (n. f.); gen. sg. leornunge; nom. pl. leornunga. study; learning.\par
% glossary-block:B00938:end
% glossary-block:B00939:begin
\needspace{3\baselineskip}
\textbf{leornungcniht} (n. m.); gen. sg. leornungcnihtes; nom. pl. leornungcnihtas. disciple.\par
% glossary-block:B00939:end
% glossary-block:B00940:begin
\needspace{3\baselineskip}
\textbf{lēoð} (n. n.); gen. sg. lēoðes; nom. pl. lēoð. poem; poetry; verse.\par
% glossary-block:B00940:end
% glossary-block:B00941:begin
\needspace{3\baselineskip}
\textbf{lēoðwyrhta} (n. m.); gen. sg. lēoðwyrhtan; nom. pl. lēoðwyrhtan. poet.\par
% glossary-block:B00941:end
% glossary-block:B00942:begin
\needspace{3\baselineskip}
\textbf{libban} (v. weak); pret. 1 sg. lifde; past ptcp. (ge)lifd; pres. 3 sg. leofað. live. living. Supplied weak participial adjective.\par
% glossary-block:B00942:end
% glossary-block:B00943:begin
\needspace{3\baselineskip}
\textbf{liccian} (v. weak); pret. 1 sg. liccode; past ptcp. liccod. lick.\par
% glossary-block:B00943:end
% glossary-block:B00944:begin
\needspace{3\baselineskip}
\textbf{licgan} (v. strong); pret. 1 sg. læg; pret. pl. lǣgon; past ptcp. legen. lie; be at rest. lie; be lying. Supplied participle describing the possessed persons, not a habban perfect.\par
% glossary-block:B00944:end
% glossary-block:B00945:begin
\needspace{3\baselineskip}
\textbf{līchama} (n. m.); gen. sg. līchaman; nom. pl. līchaman. body.\par
% glossary-block:B00945:end
% glossary-block:B00946:begin
\needspace{3\baselineskip}
\textbf{līchamlic} (adj.). bodily.\par
% glossary-block:B00946:end
% glossary-block:B00947:begin
\needspace{3\baselineskip}
\textbf{lician} (v.). please. The thing pleasing is the subject; the persons pleased are dative in the supplied reading.\par
% glossary-block:B00947:end
% glossary-block:B00948:begin
\needspace{3\baselineskip}
\textbf{līefan} (v.). allow; permit.\par
% glossary-block:B00948:end
% glossary-block:B00949:begin
\needspace{3\baselineskip}
\textbf{līeg} (n. m.). thunderbolt; lightning. Source levin sense.\par
% glossary-block:B00949:end
% glossary-block:B00950:begin
\needspace{3\baselineskip}
\textbf{līehting} (n. f.); gen. sg. līehtinge; nom. pl. līehtinga. illumination.\par
% glossary-block:B00950:end
% glossary-block:B00951:begin
\needspace{3\baselineskip}
\textbf{līf} (n. n.); gen. sg. līfes; nom. pl. līf. life. life. An individual's lived life, including cognate accusative with libban.\par
% glossary-block:B00951:end
% glossary-block:B00952:begin
\needspace{3\baselineskip}
\textbf{lifde} — pret. sg.: see \textbf{libban}.\par
% glossary-block:B00952:end
% glossary-block:B00953:begin
\needspace{3\baselineskip}
\textbf{lifian} — citation/spelling variant: see \textbf{libban}.\par
% glossary-block:B00953:end
% glossary-block:B00954:begin
\needspace{3\baselineskip}
\textbf{līð} — pres. ind. 3 sg., variant: see \textbf{licgan}.\par
% glossary-block:B00954:end
% glossary-block:B00955:begin
\needspace{3\baselineskip}
\textbf{lof} (n. m.); gen. sg. lofes. praise.\par
% glossary-block:B00955:end
% glossary-block:B00956:begin
\needspace{3\baselineskip}
\textbf{lor} (n. n.). destruction; ruin.\par
% glossary-block:B00956:end
% glossary-block:B00957:begin
\needspace{3\baselineskip}
\textbf{lufe} (n. f.). love.\par
% glossary-block:B00957:end
% glossary-block:B00958:begin
\needspace{3\baselineskip}
\textbf{lufian} (v. weak); pret. 1 sg. lufode; past ptcp. (ge)lufod. love.\par
% glossary-block:B00958:end
% glossary-block:B00959:begin
\needspace{3\baselineskip}
\textbf{lufiend} (n. m.). lover.\par
% glossary-block:B00959:end
% glossary-block:B00960:begin
\needspace{3\baselineskip}
\textbf{lufu} (n. f.). love.\par
% glossary-block:B00960:end
% glossary-block:B00961:begin
\needspace{3\baselineskip}
\textbf{lyft} (n. n.); gen. sg. lyftes; nom. pl. lyft. air.\par
% glossary-block:B00961:end
% glossary-block:B00962:begin
\needspace{3\baselineskip}
\textbf{lyre} (n. m.); gen. sg. lyres; nom. pl. lyras. loss.\par
% glossary-block:B00962:end
% glossary-block:B00963:begin
\needspace{3\baselineskip}
\textbf{lystan} (v.). please; be desired. Impersonal source expression; keep the complete supplied frame.\par
% glossary-block:B00963:end
% glossary-block:B00964:begin
\needspace{3\baselineskip}
\textbf{lȳtel} (adj.); comp. lǣssa; superl. lǣst. little; small. Stem lȳtl- before vowel endings; lȳtel- before -ne, -re and -ra.\par
% glossary-block:B00964:end
\pdfbookmark[1]{M}{outline-B00965}
% glossary-block:B00965:begin
\needspace{5\baselineskip}
\subsection*{M}
% glossary-block:B00965:end
% glossary-block:B00966:begin
\textbf{mā} (adv.). more. Adverb of amount. more; again. In the negative phrase never more/never again. more. More in number, distinct from degree or greater value.\par
% glossary-block:B00966:end
% glossary-block:B00967:begin
\needspace{3\baselineskip}
\textbf{magan} — citation/spelling variant: see \textbf{mugan}.\par
% glossary-block:B00967:end
% glossary-block:B00968:begin
\needspace{3\baselineskip}
\textbf{man} — citation/spelling variant: see \textbf{mann}.\par
% glossary-block:B00968:end
% glossary-block:B00969:begin
\needspace{3\baselineskip}
\textbf{mangere} (n. m.); gen. sg. mangeres; nom. pl. mangeras. merchant.\par
% glossary-block:B00969:end
% glossary-block:B00970:begin
\needspace{3\baselineskip}
\textbf{manig} (adj.). many.\par
% glossary-block:B00970:end
% glossary-block:B00971:begin
\needspace{3\baselineskip}
\textbf{manigfeald} (adj.). manifold; many.\par
% glossary-block:B00971:end
% glossary-block:B00972:begin
\needspace{3\baselineskip}
\textbf{mann} (n.; impersonal pronoun m.); gen. sg. mannes; nom. pl. męnn. man; human being. one; people in general. Impersonal man with singular finite verb.\par
% glossary-block:B00972:end
% glossary-block:B00973:begin
\needspace{3\baselineskip}
\textbf{māra} (adj.). greater; larger. Comparative adjective. greater; more important. Value/importance in the life/raiment comparison. more; a greater amount. Nominal/attributive quantity, distinct from more in number. Comparative of micel.\par
% glossary-block:B00973:end
% glossary-block:B00974:begin
\needspace{3\baselineskip}
\textbf{Maria} (proper name). Mary.\par
% glossary-block:B00974:end
% glossary-block:B00975:begin
\needspace{3\baselineskip}
\textbf{mǣden} (n. n.); gen. sg. mǣdenes; nom. pl. mǣden. girl; maiden.\par
% glossary-block:B00975:end
% glossary-block:B00976:begin
\needspace{3\baselineskip}
\textbf{mæg} — pres. ind. 1/3 sg.: see \textbf{mugan}.\par
% glossary-block:B00976:end
% glossary-block:B00977:begin
\needspace{3\baselineskip}
\textbf{mægen} (n. n.); gen. sg. mægenes; nom. pl. mægnu. power; strength. Heavenly powers in 37.2 is the contextual plural use.\par
% glossary-block:B00977:end
% glossary-block:B00978:begin
\needspace{3\baselineskip}
\textbf{mægenðrymm} (n. m.); gen. sg. mægenðrymmes. majesty; glory.\par
% glossary-block:B00978:end
% glossary-block:B00979:begin
\needspace{3\baselineskip}
\textbf{mægenðrymnes} (n. f.); gen. sg. mægenðrymnesse. glory; majesty.\par
% glossary-block:B00979:end
% glossary-block:B00980:begin
\needspace{3\baselineskip}
\textbf{mægð} (n. f.). maiden; girl. Reference in a source or explanatory note.\par
% glossary-block:B00980:end
% glossary-block:B00981:begin
\needspace{3\baselineskip}
\textbf{mǣgð} (n. f.); gen. sg. mǣgðe; nom. pl. mǣgða. tribe; nation.\par
% glossary-block:B00981:end
% glossary-block:B00982:begin
\needspace{3\baselineskip}
\textbf{mægðhād} (n. m.). maidenhood; virginity.\par
% glossary-block:B00982:end
% glossary-block:B00983:begin
\needspace{3\baselineskip}
\textbf{mǣre} (adj.). great; renowned. Distinct from simple physical size where the task requires the renowned sense.\par
% glossary-block:B00983:end
% glossary-block:B00984:begin
\needspace{3\baselineskip}
\textbf{mǣrsian} (v. weak); pret. 1 sg. mǣrsode; past ptcp. mǣrsod. make great; magnify; glorify.\par
% glossary-block:B00984:end
% glossary-block:B00985:begin
\needspace{3\baselineskip}
\textbf{mæsseprēost} (n. m.); gen. sg. mæsseprēostes; nom. pl. mæsseprēostas. priest.\par
% glossary-block:B00985:end
% glossary-block:B00986:begin
\needspace{3\baselineskip}
\textbf{mǣst} (adj.). greatest; largest. Superlative adjective. Superlative of micel.\par
% glossary-block:B00986:end
% glossary-block:B00987:begin
\needspace{3\baselineskip}
\textbf{mǣð} (n. f.). measure; capacity.\par
% glossary-block:B00987:end
% glossary-block:B00988:begin
\needspace{3\baselineskip}
\textbf{mē} — dat. sg.; acc. sg.: see \textbf{ic}.\par
% glossary-block:B00988:end
% glossary-block:B00989:begin
\needspace{3\baselineskip}
\textbf{meahte} — pret. sg.: see \textbf{mugan}.\par
% glossary-block:B00989:end
% glossary-block:B00990:begin
\needspace{3\baselineskip}
\textbf{mec} — acc. sg., variant: see \textbf{ic}.\par
% glossary-block:B00990:end
% glossary-block:B00991:begin
\needspace{3\baselineskip}
\textbf{mēd} (n. f.); gen. sg. mēde; nom. pl. mēda. reward. hire; wages paid for work.\par
% glossary-block:B00991:end
% glossary-block:B00992:begin
\needspace{3\baselineskip}
\textbf{mēder} — dat. sg.: see \textbf{mōdor}.\par
% glossary-block:B00992:end
% glossary-block:B00993:begin
\needspace{3\baselineskip}
\textbf{medmicel} (adj.). short; brief.\par
% glossary-block:B00993:end
% glossary-block:B00994:begin
\needspace{3\baselineskip}
\textbf{mehte} — pret. sg., variant: see \textbf{mugan}.\par
% glossary-block:B00994:end
% glossary-block:B00995:begin
\needspace{3\baselineskip}
\textbf{męnigu} (n. f.). company; crowd.\par
% glossary-block:B00995:end
% glossary-block:B00996:begin
\needspace{3\baselineskip}
\textbf{męnn} — dat. sg.; nom./acc. pl.: see \textbf{mann}.\par
% glossary-block:B00996:end
% glossary-block:B00997:begin
\needspace{3\baselineskip}
\textbf{męregrēote} (n. f.); gen. sg. męregrēotan; nom. pl. męregrēotan. pearl.\par
% glossary-block:B00997:end
% glossary-block:B00998:begin
\needspace{3\baselineskip}
\textbf{męreswīn} (n. n.); gen. sg. męreswīnes; nom. pl. męreswīn. dolphin.\par
% glossary-block:B00998:end
% glossary-block:B00999:begin
\needspace{3\baselineskip}
\textbf{męrgen} — citation/spelling variant: see \textbf{morgen}.\par
% glossary-block:B00999:end
% glossary-block:B01000:begin
\needspace{3\baselineskip}
\textbf{mētan} (v. weak); pret. 1 sg. mētte; past ptcp. (ge)mēt(t). meet; find.\par
% glossary-block:B01000:end
% glossary-block:B01001:begin
\needspace{3\baselineskip}
\textbf{męte} (n. m.); gen. sg. mętes; nom. pl. mętas. food.\par
% glossary-block:B01001:end
% glossary-block:B01002:begin
\needspace{3\baselineskip}
\textbf{micel} (adj.); comp. māra; superl. mǣst. great; large. much; an amount. Quantity sense. great; loud. Loud is the contextual force with a voice. Stem micl- before vowel endings; micel- before -ne, -re and -ra.\par
% glossary-block:B01002:end
% glossary-block:B01003:begin
\needspace{3\baselineskip}
\textbf{Michael} (proper name). Michael.\par
% glossary-block:B01003:end
% glossary-block:B01004:begin
\needspace{3\baselineskip}
\textbf{mid} (prep.; adv. in fixed expressions). with; by means of. Dative instrument or means. with. Manner in the supplied compassion phrase. with; accompanying. Dative person or cargo in the ordinary accompaniment frame. with. Circumstances of risk or loss in the supplied phrase. with. Manner/circumstance of victory. completely. Only in fixed mid ealle. with. Dative garment in the assigned clothing frame. with. Dative conversational partner. with; along with. Supplied mid wes, be with them. with. Accompanying quality or mental disposition in the source. with. Material or clothing in the supplied adornment phrase. with. Exact cloth/cloak complements of begyrd and bewǣfed. with; by means of.\par
% glossary-block:B01004:end
% glossary-block:B01005:begin
\needspace{3\baselineskip}
\textbf{middangeard} (n. m.); gen. sg. middangeardes; nom. pl. middangeardas. world; earth.\par
% glossary-block:B01005:end
% glossary-block:B01006:begin
\needspace{3\baselineskip}
\textbf{middes} (adverb or particle). at noon. Whole supplied tō middes dæges, with genitives.\par
% glossary-block:B01006:end
% glossary-block:B01007:begin
\needspace{3\baselineskip}
\textbf{mid-ðȳ} (conj.). when; while.\par
% glossary-block:B01007:end
% glossary-block:B01008:begin
\needspace{3\baselineskip}
\textbf{miht} (n. f.). ability; power. powers; powerful people. Those holding high station in the source's concluding human application.\par
% glossary-block:B01008:end
% glossary-block:B01009:begin
\needspace{3\baselineskip}
\textbf{mihtig} (adj.). mighty; powerful. the powerful; the mighty. Substantival adjective. Stem mihtg- before vowel endings and mihtig- before consonantal endings; unsyncopated forms also occur.\par
% glossary-block:B01009:end
% glossary-block:B01010:begin
\needspace{3\baselineskip}
\textbf{mildheort} (adj.). merciful; compassionate.\par
% glossary-block:B01010:end
% glossary-block:B01011:begin
\needspace{3\baselineskip}
\textbf{mildheortnes} (n. f.); gen. sg. mildheortnesse; nom. pl. mildheortnessa. compassion. mercy; mercies.\par
% glossary-block:B01011:end
% glossary-block:B01012:begin
\needspace{3\baselineskip}
\textbf{mīn} (possessive adjective). my; mine. Also gen. sg.: see \textbf{ic}.\par
% glossary-block:B01012:end
% glossary-block:B01013:begin
\needspace{3\baselineskip}
\textbf{mislic} (adj.). various.\par
% glossary-block:B01013:end
% glossary-block:B01014:begin
\needspace{3\baselineskip}
\textbf{mōd} (n. n.). mind; attention. heart; mind.\par
% glossary-block:B01014:end
% glossary-block:B01015:begin
\needspace{3\baselineskip}
\textbf{mōdor} (n. f.); gen. sg. mōdor; nom. pl. mōdru. mother. Plural variant mōdra; dative singular mēder, also a late genitive variant.\par
% glossary-block:B01015:end
% glossary-block:B01016:begin
\needspace{3\baselineskip}
\textbf{mōna} (n. m.); gen. sg. mōnan; nom. pl. mōnan. moon.\par
% glossary-block:B01016:end
% glossary-block:B01017:begin
\needspace{3\baselineskip}
\textbf{mōnað} (n. m.). month. Source plural months.\par
% glossary-block:B01017:end
% glossary-block:B01018:begin
\needspace{3\baselineskip}
\textbf{morgen} (n. m.); gen. sg. morgnes; nom. pl. morgnas. morning.\par
% glossary-block:B01018:end
% glossary-block:B01019:begin
\needspace{3\baselineskip}
\textbf{mōste} — pret. sg.: see \textbf{mōtan}.\par
% glossary-block:B01019:end
% glossary-block:B01020:begin
\needspace{3\baselineskip}
\textbf{mōtan} (v. preterite-present); pres. 1/3 sg. mōt; pres. pl. mōton; pret. 1 sg. mōste. may; be allowed to. Conventional citation label; no usable infinitive is supplied. See the finite forms in Unit 6. No past-participle citation is given in the checked source paradigm.\par
% glossary-block:B01020:end
% glossary-block:B01021:begin
\needspace{3\baselineskip}
\textbf{mugan} (v. preterite-present); pres. 1/3 sg. mæg; pres. pl. magon; pret. 1 sg. meahte. can; be able to. be able; could. Supplied modal form; ordinary capability, not permission. Conventional citation label; no usable infinitive is supplied. See the finite forms in Unit 6. No past-participle citation is given in the checked source paradigm.\par
% glossary-block:B01021:end
% glossary-block:B01022:begin
\needspace{3\baselineskip}
\textbf{munan} (v.). intend. Reference in a source or explanatory note.\par
% glossary-block:B01022:end
% glossary-block:B01023:begin
\needspace{3\baselineskip}
\textbf{munt} (n. m.); gen. sg. muntes; nom. pl. muntas. mountain.\par
% glossary-block:B01023:end
% glossary-block:B01024:begin
\needspace{3\baselineskip}
\textbf{munuc} (n. m.); gen. sg. munuces; nom. pl. munucas. monk.\par
% glossary-block:B01024:end
% glossary-block:B01025:begin
\needspace{3\baselineskip}
\textbf{mūð} (n. m.); gen. sg. mūðes; nom. pl. mūðas. mouth.\par
% glossary-block:B01025:end
% glossary-block:B01026:begin
\needspace{3\baselineskip}
\textbf{myngian} (v.); past ptcp. myngod. admonish; prompt. Supplied participial source use.\par
% glossary-block:B01026:end
% glossary-block:B01027:begin
\needspace{3\baselineskip}
\textbf{mynster} (n. n.); gen. sg. mynstres; nom. pl. mynstru. monastery.\par
% glossary-block:B01027:end
\pdfbookmark[1]{N}{outline-B01028}
% glossary-block:B01028:begin
\needspace{5\baselineskip}
\subsection*{N}
% glossary-block:B01028:end
% glossary-block:B01029:begin
\textbf{nā} (adv.). not. Negates a contrasted constituent, distinct from finite ne.\par
% glossary-block:B01029:end
% glossary-block:B01030:begin
\needspace{3\baselineskip}
\textbf{nabban} (v. weak); pret. 1 sg. næfde; past ptcp. (ge)næfd; pres. 3 sg. næfð. have not.\par
% glossary-block:B01030:end
% glossary-block:B01031:begin
\needspace{3\baselineskip}
\textbf{nacod} (adj.). naked. naked people. Substantival adjective.\par
% glossary-block:B01031:end
% glossary-block:B01032:begin
\needspace{3\baselineskip}
\textbf{nafogār} (n.). auger. Reference in a source or explanatory note.\par
% glossary-block:B01032:end
% glossary-block:B01033:begin
\needspace{3\baselineskip}
\textbf{nāgan} (v.); pres. 3 sg. nāh. not have; not possess. Negative possessive verb in its supplied forms.\par
% glossary-block:B01033:end
% glossary-block:B01034:begin
\needspace{3\baselineskip}
\textbf{nāh} — pres. ind. 3 sg.: see \textbf{nāgan}.\par
% glossary-block:B01034:end
% glossary-block:B01035:begin
\needspace{3\baselineskip}
\textbf{nāht} (substantive n.). nothing; naught.\par
% glossary-block:B01035:end
% glossary-block:B01036:begin
\needspace{3\baselineskip}
\textbf{nāhwæðer} (adverb or particle). neither.\par
% glossary-block:B01036:end
% glossary-block:B01037:begin
\needspace{3\baselineskip}
\textbf{nālæs} (adv.). not at all.\par
% glossary-block:B01037:end
% glossary-block:B01038:begin
\needspace{3\baselineskip}
\textbf{nama} (n. m.); gen. sg. naman; nom. pl. naman. name.\par
% glossary-block:B01038:end
% glossary-block:B01039:begin
\needspace{3\baselineskip}
\textbf{nān} (indefinite). no; no one.\par
% glossary-block:B01039:end
% glossary-block:B01040:begin
\needspace{3\baselineskip}
\textbf{nāst} — negative pres. ind. 2 sg.: see \textbf{nytan}.\par
% glossary-block:B01040:end
% glossary-block:B01041:begin
\needspace{3\baselineskip}
\textbf{nāt} — negative pres. ind. 1 sg.; supplied source form: see \textbf{nytan}.\par
% glossary-block:B01041:end
% glossary-block:B01042:begin
\needspace{3\baselineskip}
\textbf{nāwiht} (indefinite pronoun/substantive). nothing. nothing; naught.\par
% glossary-block:B01042:end
% glossary-block:B01043:begin
\needspace{3\baselineskip}
\textbf{næddre} (n. f.). snake; adder. Supplied weak feminine plural.\par
% glossary-block:B01043:end
% glossary-block:B01044:begin
\needspace{3\baselineskip}
\textbf{nǣdl} (n. f.); gen. sg. nǣdle; nom. pl. nǣdla. needle.\par
% glossary-block:B01044:end
% glossary-block:B01045:begin
\needspace{3\baselineskip}
\textbf{nædre} (n. f.); gen. sg. nædran; nom. pl. nædran. serpent.\par
% glossary-block:B01045:end
% glossary-block:B01046:begin
\needspace{3\baselineskip}
\textbf{næfde} — pret. sg.: see \textbf{nabban}.\par
% glossary-block:B01046:end
% glossary-block:B01047:begin
\needspace{3\baselineskip}
\textbf{nǣfre} (adv.). never.\par
% glossary-block:B01047:end
% glossary-block:B01048:begin
\needspace{3\baselineskip}
\textbf{næfð} — pres. ind. 3 sg.: see \textbf{nabban}.\par
% glossary-block:B01048:end
% glossary-block:B01049:begin
\needspace{3\baselineskip}
\textbf{nǣnig} (indefinite pronoun/adjective). no; not any; none. Reference in a source or explanatory note.\par
% glossary-block:B01049:end
% glossary-block:B01050:begin
\needspace{3\baselineskip}
\textbf{nǣron} — negative pret. ind. pl.: see \textbf{wesan}.\par
% glossary-block:B01050:end
% glossary-block:B01051:begin
\needspace{3\baselineskip}
\textbf{næs} — negative pret. ind. 1/3 sg.: see \textbf{wesan}.\par
% glossary-block:B01051:end
% glossary-block:B01052:begin
\needspace{3\baselineskip}
\textbf{ne} (particle). not; simple clause negation.\par
% glossary-block:B01052:end
% glossary-block:B01053:begin
\needspace{3\baselineskip}
\textbf{nēah} (adv.). near; at hand.\par
% glossary-block:B01053:end
% glossary-block:B01054:begin
\needspace{3\baselineskip}
\textbf{nēahgebūr} (n. m.); gen. sg. nēahgebūres; nom. pl. nēahgebūras. neighbour.\par
% glossary-block:B01054:end
% glossary-block:B01055:begin
\needspace{3\baselineskip}
\textbf{nēalǣcan} (v. weak); pret. 1 sg. nēalǣcte; past ptcp. nēalǣced. approach; draw near. Dative person or place approached when the goal is expressed.\par
% glossary-block:B01055:end
% glossary-block:B01056:begin
\needspace{3\baselineskip}
\textbf{nearu} (adj.). narrow.\par
% glossary-block:B01056:end
% glossary-block:B01057:begin
\needspace{3\baselineskip}
\textbf{nele} — negative pres. ind. 3 sg.: see \textbf{willan}.\par
% glossary-block:B01057:end
% glossary-block:B01058:begin
\needspace{3\baselineskip}
\textbf{nelle} — negative pres. ind. 1 sg.: see \textbf{willan}.\par
% glossary-block:B01058:end
% glossary-block:B01059:begin
\needspace{3\baselineskip}
\textbf{nelt} — negative pres. ind. 2 sg.: see \textbf{willan}.\par
% glossary-block:B01059:end
% glossary-block:B01060:begin
\needspace{3\baselineskip}
\textbf{nēosian} (v.). seek; look for (+ genitive). genitive object; reference comparison with genēosian.\par
% glossary-block:B01060:end
% glossary-block:B01061:begin
\needspace{3\baselineskip}
\textbf{nēowolnes} (n. f.); gen. sg. nēowolnesse; nom. pl. nēowolnessa. abyss; the deep.\par
% glossary-block:B01061:end
% glossary-block:B01062:begin
\needspace{3\baselineskip}
\textbf{nęrian} (v. weak); pret. 1 sg. nęrede; past ptcp. (ge)nęred. save.\par
% glossary-block:B01062:end
% glossary-block:B01063:begin
\needspace{3\baselineskip}
\textbf{nest} (n. n.); gen. sg. nestes; nom. pl. nest. nest.\par
% glossary-block:B01063:end
% glossary-block:B01064:begin
\needspace{3\baselineskip}
\textbf{nętt} (n. n.); gen. sg. nęttes; nom. pl. nętt. net.\par
% glossary-block:B01064:end
% glossary-block:B01065:begin
\needspace{3\baselineskip}
\textbf{nīede} (adverb or particle). necessarily; of necessity.\par
% glossary-block:B01065:end
% glossary-block:B01066:begin
\needspace{3\baselineskip}
\textbf{nīehst} (adv.). at length. Whole supplied æt nīehstan.\par
% glossary-block:B01066:end
% glossary-block:B01067:begin
\needspace{3\baselineskip}
\textbf{nīeten} (n. n.); gen. sg. nīetenes; nom. pl. nīetenu. animal; beast; cattle. Cattle is the selected plural use; retain the early task distinction from dēor.\par
% glossary-block:B01067:end
% glossary-block:B01068:begin
\needspace{3\baselineskip}
\textbf{nigon} (cardinal numeral). nine. Cardinal 9.\par
% glossary-block:B01068:end
% glossary-block:B01069:begin
\needspace{3\baselineskip}
\textbf{nigoða} (ordinal). ninth.\par
% glossary-block:B01069:end
% glossary-block:B01070:begin
\needspace{3\baselineskip}
\textbf{niht} (n. f.); gen. sg. niht; nom. pl. niht. night.\par
% glossary-block:B01070:end
% glossary-block:B01071:begin
\needspace{3\baselineskip}
\textbf{niman} (v. strong); pret. 1 sg. nōm; pret. pl. nōmon; past ptcp. numen. take. Accusative person or thing taken, including food.\par
% glossary-block:B01071:end
% glossary-block:B01072:begin
\needspace{3\baselineskip}
\textbf{nis} — negative pres. ind. 3 sg.: see \textbf{wesan}.\par
% glossary-block:B01072:end
% glossary-block:B01073:begin
\needspace{3\baselineskip}
\textbf{nīwe} (adj.). new.\par
% glossary-block:B01073:end
% glossary-block:B01074:begin
\needspace{3\baselineskip}
\textbf{nōht} — citation/spelling variant: see \textbf{nāwiht}.\par
% glossary-block:B01074:end
% glossary-block:B01075:begin
\needspace{3\baselineskip}
\textbf{nolde} — negative pret. sg.: see \textbf{willan}.\par
% glossary-block:B01075:end
% glossary-block:B01076:begin
\needspace{3\baselineskip}
\textbf{norðan} (adv.). from the north.\par
% glossary-block:B01076:end
% glossary-block:B01077:begin
\needspace{3\baselineskip}
\textbf{nū} (adv.; conj.). now.\par
% glossary-block:B01077:end
% glossary-block:B01078:begin
\needspace{3\baselineskip}
\textbf{nytan} (v.). not know. Supplied negative knowing form. not know.\par
% glossary-block:B01078:end
\pdfbookmark[1]{O}{outline-B01079}
% glossary-block:B01079:begin
\needspace{5\baselineskip}
\subsection*{O}
% glossary-block:B01079:end
% glossary-block:B01080:begin
\textbf{of} (prep.). from; out of. With dative. of; from. Portion from a whole in the supplied partitive phrases; dative whole. sense for sense. Literal correspondence in the supplied translation expression. from. Starting condition of old age in the supplied rejuvenation comparison. from. Starting generation in the supplied temporal range. of; from. Origin of wisdom in the supplied begetting statement.\par
% glossary-block:B01080:end
% glossary-block:B01081:begin
\needspace{3\baselineskip}
\textbf{of} (adv.). off.\par
% glossary-block:B01081:end
% glossary-block:B01082:begin
\needspace{3\baselineskip}
\textbf{ofen} (n. m.); gen. sg. ofenes; nom. pl. ofenas. oven.\par
% glossary-block:B01082:end
% glossary-block:B01083:begin
\needspace{3\baselineskip}
\textbf{ofer} (prep.). over; across. Spatial frame; case varies with the specific construction. Do not impose a universal motion rule. over. Accusative dependent in the taught authority frame. above. Rank, with accusative. toward. In gōd ofer + accusative people. over; about. Object of rejoicing, with the assigned accusative frame. above. Rank, with explicitly supplied dative complement; retain this limit. above; more than. Preference in the supplied comparison.\par
% glossary-block:B01083:end
% glossary-block:B01084:begin
\needspace{3\baselineskip}
\textbf{oferwinnan} (v. strong); pret. 1 sg. oferwann; pret. pl. oferwunnon; past ptcp. oferwunnen. overcome; conquer.\par
% glossary-block:B01084:end
% glossary-block:B01085:begin
\needspace{3\baselineskip}
\textbf{oferworhte} — pret. sg.: see \textbf{oferwyrcean}.\par
% glossary-block:B01085:end
% glossary-block:B01086:begin
\needspace{3\baselineskip}
\textbf{oferwyrcean} (v. weak); pret. 1 sg. oferworhte; past ptcp. oferworht. overlay; cover. Accusative object, mid + dative material used to cover it.\par
% glossary-block:B01086:end
% glossary-block:B01087:begin
\needspace{3\baselineskip}
\textbf{ofslēan} (v.); past ptcp. ofslægen. slay; kill.\par
% glossary-block:B01087:end
% glossary-block:B01088:begin
\needspace{3\baselineskip}
\textbf{oft} (adv.). often.\par
% glossary-block:B01088:end
% glossary-block:B01089:begin
\needspace{3\baselineskip}
\textbf{ofðyncean} (v. impersonal); pres. 3 sg. ofðyncð; pret. 3 sg. ofðūhte; past ptcp. ofðūht. grieve; displease. Impersonal third singular with dative experiencer and a content clause. Impersonal use: the course supplies third-singular finite forms.\par
% glossary-block:B01089:end
% glossary-block:B01090:begin
\needspace{3\baselineskip}
\textbf{ofwundrian} (v.); past ptcp. ofwundrod. be astonished.\par
% glossary-block:B01090:end
% glossary-block:B01091:begin
\needspace{3\baselineskip}
\textbf{ōht} — citation/spelling variant: see \textbf{āwiht}.\par
% glossary-block:B01091:end
% glossary-block:B01092:begin
\needspace{3\baselineskip}
\textbf{olfend} (n. m.); gen. sg. olfendes; nom. pl. olfendas. camel.\par
% glossary-block:B01092:end
% glossary-block:B01093:begin
\needspace{3\baselineskip}
\textbf{on} (prep.). in; on. Dative for the taught ordinary location construction. in the beginning. Whole supplied on anginne. on. Time preposition in the stated-day expression. in. Supplied circumstance of the storm. into; in two. Supplied division expression. alive. Only in supplied on līfe. according to. Only the supplied classificatory expression with hīw. in one's control. on + dative possessed geweald. in. Supported expression in his majesty/glory. at one's side. Accusative hand in the bounded placement expression. on one's side. The bounded on … hand side expression uses accusative, including with standan. in one's name. on + dative naman, with genitive possessor. rich in. Only in supplied berende on + dative. into; onto. Accusative goal of placement in the taught frame. in; amid. Circumstances of the stated unrest. in; with respect to. Quality or domain of description. in. Accusative person believed/trusted in. during; in. Supplied lifetime expression. among; in. Dative people among whom friends are obtained. in worldly life. Supplied temporal/situational expression. during. Dative time expression. in. Circumstance of the day's heat in the supplied source. fishing. Supplied on fiscnoð after gān. in; wearing. Supplied on scēapa gegierelum, in sheep's garments. at work; engaged in. Only the exact supplied activity phrases. upon; into. Supplied chancing-upon idiom. in. Source/object of joy in the supplied frame. to. Succeeding generation in the supplied temporal range. upon. Accusative beloved in the supplied love-falls-upon idiom. in the morning. Whole supplied on morgen; accusative singular. into. Destination of sale in the supplied expression.\par
% glossary-block:B01093:end
% glossary-block:B01094:begin
\needspace{3\baselineskip}
\textbf{on} (adv.). on. Adverb in the cloak-wearing relative clause.\par
% glossary-block:B01094:end
% glossary-block:B01095:begin
\needspace{3\baselineskip}
\textbf{oncnāwan} (v. strong); pret. 1 sg. oncnēow; pret. pl. oncnēowon; past ptcp. oncnāwen; pres. 3 sg. oncnǣwð. recognize. recognize; perceive. Supplied past source form. recognize; acknowledge. recognize; know by recognition.\par
% glossary-block:B01095:end
% glossary-block:B01096:begin
\needspace{3\baselineskip}
\textbf{oncnǣwð} — pres. ind. 3 sg.: see \textbf{oncnāwan}.\par
% glossary-block:B01096:end
% glossary-block:B01097:begin
\needspace{3\baselineskip}
\textbf{ondrǣdan} (v. strong); pret. 1 sg. ondrēd; pret. pl. ondrēdon; past ptcp. ondrǣden. fear.\par
% glossary-block:B01097:end
% glossary-block:B01098:begin
\needspace{3\baselineskip}
\textbf{onfēhð} — pres. ind. 3 sg.: see \textbf{onfōn}.\par
% glossary-block:B01098:end
% glossary-block:B01099:begin
\needspace{3\baselineskip}
\textbf{onfōn} (v. strong); pret. 1 sg. onfēng; pret. pl. onfēngon; past ptcp. onfangen; pres. 3 sg. onfēhð. receive; accept. Dative thing received in this selected frame; the later accusative payment frame is separate. receive; accept. A person received into care in the supplied frame.\par
% glossary-block:B01099:end
% glossary-block:B01100:begin
\needspace{3\baselineskip}
\textbf{ongēan} (adv.; prep.; conj.). against; opposite. Accusative landmark or opponent. back; again. by the time that. Complete supplied ongēan ðætte clause.\par
% glossary-block:B01100:end
% glossary-block:B01101:begin
\needspace{3\baselineskip}
\textbf{ongeslēan} (v.). inflict slaughter; slay.\par
% glossary-block:B01101:end
% glossary-block:B01102:begin
\needspace{3\baselineskip}
\textbf{ongietan} (v. strong); pret. 1 sg. ongeat; pret. pl. ongēaton; past ptcp. ongieten. understand; perceive. Supplied infinitive after a modal. perceive; realize. understand; perceive.\par
% glossary-block:B01102:end
% glossary-block:B01103:begin
\needspace{3\baselineskip}
\textbf{onginnan} (v.); pret. pl. ongunnon. begin. Supplied preterite began. begin. begin. Supplied playing/praising action as complement.\par
% glossary-block:B01103:end
% glossary-block:B01104:begin
\needspace{3\baselineskip}
\textbf{onlīesan} (v. weak); pret. 1 sg. onlīesde; past ptcp. onlīesed. release. Accusative person or thing released.\par
% glossary-block:B01104:end
% glossary-block:B01105:begin
\needspace{3\baselineskip}
\textbf{onsęndan} (v.). send.\par
% glossary-block:B01105:end
% glossary-block:B01106:begin
\needspace{3\baselineskip}
\textbf{ontȳnan} (v.). open. Supplied preterite opened.\par
% glossary-block:B01106:end
% glossary-block:B01107:begin
\needspace{3\baselineskip}
\textbf{onwęg} (adv.). away.\par
% glossary-block:B01107:end
% glossary-block:B01108:begin
\needspace{3\baselineskip}
\textbf{onwinnan} (v.). assail; attack.\par
% glossary-block:B01108:end
% glossary-block:B01109:begin
\needspace{3\baselineskip}
\textbf{open} (adj.). open.\par
% glossary-block:B01109:end
% glossary-block:B01110:begin
\needspace{3\baselineskip}
\textbf{ōra} (n. m.). vein of metal; ore (sense uncertain). Supplied dative plural after on.\par
% glossary-block:B01110:end
% glossary-block:B01111:begin
\needspace{3\baselineskip}
\textbf{ōrmōdnes} (n. f.). despair; hopelessness.\par
% glossary-block:B01111:end
% glossary-block:B01112:begin
\needspace{3\baselineskip}
\textbf{Osburg} (proper name). Osburg.\par
% glossary-block:B01112:end
% glossary-block:B01113:begin
\needspace{3\baselineskip}
\textbf{ōð} (prep.). as far as; up to.\par
% glossary-block:B01113:end
% glossary-block:B01114:begin
\needspace{3\baselineskip}
\textbf{oð} (prep.). up to; until. With accusative.\par
% glossary-block:B01114:end
% glossary-block:B01115:begin
\needspace{3\baselineskip}
\textbf{ōðer} (adj.; ordinal). second. other; another. remaining; the rest. The rest of the forest/earth in the source.\par
% glossary-block:B01115:end
% glossary-block:B01116:begin
\needspace{3\baselineskip}
\textbf{oð-ðæt} (conj.). until.\par
% glossary-block:B01116:end
% glossary-block:B01117:begin
\needspace{3\baselineskip}
\textbf{oððe} (conj.). or.\par
% glossary-block:B01117:end
\pdfbookmark[1]{P}{outline-B01118}
% glossary-block:B01118:begin
\needspace{5\baselineskip}
\subsection*{P}
% glossary-block:B01118:end
% glossary-block:B01119:begin
\textbf{palmtwīg} (n. n.); gen. sg. palmtwīges; nom. pl. palmtwīg. palm twig; palm branch.\par
% glossary-block:B01119:end
% glossary-block:B01120:begin
\needspace{3\baselineskip}
\textbf{pęning} (n. m.); gen. sg. pęninges; nom. pl. pęningas. penny; coin.\par
% glossary-block:B01120:end
% glossary-block:B01121:begin
\needspace{3\baselineskip}
\textbf{Pentapolis} (proper place name). Pentapolis. Cyrenian region in this story.\par
% glossary-block:B01121:end
% glossary-block:B01122:begin
\needspace{3\baselineskip}
\textbf{Piht} (ethnonym). Picts. Supplied people-name forms.\par
% glossary-block:B01122:end
% glossary-block:B01123:begin
\needspace{3\baselineskip}
\textbf{plega} (n. m.). play; game.\par
% glossary-block:B01123:end
% glossary-block:B01124:begin
\needspace{3\baselineskip}
\textbf{plegian} (v. weak); pret. 1 sg. plegode; past ptcp. plegod. play. Intransitive; the reading supplies local wið participant and mid instrument phrases. act; perform. Performance context; ordinary play is a separate assigned sense.\par
% glossary-block:B01124:end
% glossary-block:B01125:begin
\needspace{3\baselineskip}
\textbf{pliht} (n. m.); gen. sg. plihtes; nom. pl. plihtas. peril; risk.\par
% glossary-block:B01125:end
% glossary-block:B01126:begin
\needspace{3\baselineskip}
\textbf{pund} (n. n.); gen. sg. pundes; nom. pl. pund. pound. Weight measure in the supplied gold quantity.\par
% glossary-block:B01126:end
\pdfbookmark[1]{R}{outline-B01127}
% glossary-block:B01127:begin
\needspace{5\baselineskip}
\subsection*{R}
% glossary-block:B01127:end
% glossary-block:B01128:begin
\textbf{rǣdan} (v. weak); pret. 1 sg. rǣdde; past ptcp. gerǣd. read. Accusative text, or absolute with no expressed object.\par
% glossary-block:B01128:end
% glossary-block:B01129:begin
\needspace{3\baselineskip}
\textbf{rēaf} (n. n.); gen. sg. rēafes; nom. pl. rēaf. raiment; clothing. clothing; garments.\par
% glossary-block:B01129:end
% glossary-block:B01130:begin
\needspace{3\baselineskip}
\textbf{rēcelēasian} (v.). neglect; despise.\par
% glossary-block:B01130:end
% glossary-block:B01131:begin
\needspace{3\baselineskip}
\textbf{regn} (n. m.); gen. sg. regnes; nom. pl. regnas. rain.\par
% glossary-block:B01131:end
% glossary-block:B01132:begin
\needspace{3\baselineskip}
\textbf{ręst} (n. f.); gen. sg. ręste; nom. pl. ręsta. rest; repose.\par
% glossary-block:B01132:end
% glossary-block:B01133:begin
\needspace{3\baselineskip}
\textbf{rēðnes} (n. f.). fury; fierceness.\par
% glossary-block:B01133:end
% glossary-block:B01134:begin
\needspace{3\baselineskip}
\textbf{rēwð} — pres. ind. 3 sg.: see \textbf{rōwan}.\par
% glossary-block:B01134:end
% glossary-block:B01135:begin
\needspace{3\baselineskip}
\textbf{rīce} (n. n.); gen. sg. rīces; nom. pl. rīcu. kingdom.\par
% glossary-block:B01135:end
% glossary-block:B01136:begin
\needspace{3\baselineskip}
\textbf{rīce} (adj.). rich; wealthy. Includes substantival rich people in the assigned frame.\par
% glossary-block:B01136:end
% glossary-block:B01137:begin
\needspace{3\baselineskip}
\textbf{rīcsian} (v. weak); pret. 1 sg. rīcsode; past ptcp. rīcsod. rule; have dominion. Used absolutely in the taught frame; no new object government.\par
% glossary-block:B01137:end
% glossary-block:B01138:begin
\needspace{3\baselineskip}
\textbf{rīdan} (v. strong); pret. 1 sg. rād; pret. pl. ridon; past ptcp. riden. ride. Intransitive in the introduced frame.\par
% glossary-block:B01138:end
% glossary-block:B01139:begin
\needspace{3\baselineskip}
\textbf{riht} (adj.). right; just.\par
% glossary-block:B01139:end
% glossary-block:B01140:begin
\needspace{3\baselineskip}
\textbf{rihtlīce} (adv.); comp. rihtlicor. In comparative rihtlicor: more correctly.\par
% glossary-block:B01140:end
% glossary-block:B01141:begin
\needspace{3\baselineskip}
\textbf{rihtlicor} — comparative adverb; exact source spelling: see \textbf{rihtlīce}.\par
% glossary-block:B01141:end
% glossary-block:B01142:begin
\needspace{3\baselineskip}
\textbf{rihtwīs} (adj.). righteous. righteous people. Substantival adjective.\par
% glossary-block:B01142:end
% glossary-block:B01143:begin
\needspace{3\baselineskip}
\textbf{rihtwīsnes} (n. f.); gen. sg. rihtwīsnesse; nom. pl. rihtwīsnessa. righteousness.\par
% glossary-block:B01143:end
% glossary-block:B01144:begin
\needspace{3\baselineskip}
\textbf{rīnan} (v.). rain. Supplied impersonal source clause with past rīnde.\par
% glossary-block:B01144:end
% glossary-block:B01145:begin
\needspace{3\baselineskip}
\textbf{rīpan} (v.). reap.\par
% glossary-block:B01145:end
% glossary-block:B01146:begin
\needspace{3\baselineskip}
\textbf{rīpe} (adj.). ripe.\par
% glossary-block:B01146:end
% glossary-block:B01147:begin
\needspace{3\baselineskip}
\textbf{rīsan} (v. strong); pret. 1 sg. rās; pret. pl. rison; past ptcp. risen. rise. rise up. Intransitive rīsan, optionally with ūp; ongēan + accusative marks opposition.\par
% glossary-block:B01147:end
% glossary-block:B01148:begin
\needspace{3\baselineskip}
\textbf{rōd} (n. f.); gen. sg. rōde; nom. pl. rōda. cross.\par
% glossary-block:B01148:end
% glossary-block:B01149:begin
\needspace{3\baselineskip}
\textbf{Rōmane} (ethnonym). Romans.\par
% glossary-block:B01149:end
% glossary-block:B01150:begin
\needspace{3\baselineskip}
\textbf{Rōmanisc} (adj.). Roman.\par
% glossary-block:B01150:end
% glossary-block:B01151:begin
\needspace{3\baselineskip}
\textbf{Rome} (proper name). Rome.\par
% glossary-block:B01151:end
% glossary-block:B01152:begin
\needspace{3\baselineskip}
\textbf{rōwan} (v. strong); pret. 1 sg. rēow; pret. pl. rēowon; past ptcp. rōwen; pres. 3 sg. rēwð. row. Intransitive; the assigned route is expressed with ofer.\par
% glossary-block:B01152:end
\pdfbookmark[1]{S}{outline-B01153}
% glossary-block:B01153:begin
\needspace{5\baselineskip}
\subsection*{S}
% glossary-block:B01153:end
% glossary-block:B01154:begin
\textbf{Saba} (proper name). Sheba.\par
% glossary-block:B01154:end
% glossary-block:B01155:begin
\needspace{3\baselineskip}
\textbf{sacc} (n. m.); gen. sg. sacces; nom. pl. saccas. sack.\par
% glossary-block:B01155:end
% glossary-block:B01156:begin
\needspace{3\baselineskip}
\textbf{sāde} — pret. sg., source variant: see \textbf{sęcgan}.\par
% glossary-block:B01156:end
% glossary-block:B01157:begin
\needspace{3\baselineskip}
\textbf{sagað} — pres. ind. 3 sg., variant: see \textbf{sęcgan}.\par
% glossary-block:B01157:end
% glossary-block:B01158:begin
\needspace{3\baselineskip}
\textbf{Salomon} (proper name); gen. sg. Salomones. Solomon.\par
% glossary-block:B01158:end
% glossary-block:B01159:begin
\needspace{3\baselineskip}
\textbf{samod} (adv.). together.\par
% glossary-block:B01159:end
% glossary-block:B01160:begin
\needspace{3\baselineskip}
\textbf{Sanctus} (foreign adjective). Saint; holy. Latin word in the source title.\par
% glossary-block:B01160:end
% glossary-block:B01161:begin
\needspace{3\baselineskip}
\textbf{sang} (n. m.); gen. sg. sanges; nom. pl. sangas. song.\par
% glossary-block:B01161:end
% glossary-block:B01162:begin
\needspace{3\baselineskip}
\textbf{sārnes} (n. f.); gen. sg. sārnesse. grief.\par
% glossary-block:B01162:end
% glossary-block:B01163:begin
\needspace{3\baselineskip}
\textbf{sāwol} (n. f.); gen. sg. sāwle; nom. pl. sāwla. soul.\par
% glossary-block:B01163:end
% glossary-block:B01164:begin
\needspace{3\baselineskip}
\textbf{sǣ} (n. m.); gen. sg. sǣs; nom. pl. sǣs; source f. gen. sg. sǣ. sea. Also feminine in the supplied source passage.\par
% glossary-block:B01164:end
% glossary-block:B01165:begin
\needspace{3\baselineskip}
\textbf{sǣd} (n. n.); gen. sg. sǣdes; nom. pl. sǣd. seed. seed; descendants. Collective descendants, distinct from literal seed.\par
% glossary-block:B01165:end
% glossary-block:B01166:begin
\needspace{3\baselineskip}
\textbf{sǣde} — dat. n. sg.: see \textbf{sǣd}; pret. sg., variant: see \textbf{sęcgan}.\par
% glossary-block:B01166:end
% glossary-block:B01167:begin
\needspace{3\baselineskip}
\textbf{sǣdtīma} (n. m.); gen. sg. sǣdtīman; nom. pl. sǣdtīman. seedtime; sowing season.\par
% glossary-block:B01167:end
% glossary-block:B01168:begin
\needspace{3\baselineskip}
\textbf{sægde} — pret. sg.: see \textbf{sęcgan}.\par
% glossary-block:B01168:end
% glossary-block:B01169:begin
\needspace{3\baselineskip}
\textbf{sægð} — pres. ind. 3 sg.: see \textbf{sęcgan}.\par
% glossary-block:B01169:end
% glossary-block:B01170:begin
\needspace{3\baselineskip}
\textbf{sǣlic} (adj.). marine; of the sea.\par
% glossary-block:B01170:end
% glossary-block:B01171:begin
\needspace{3\baselineskip}
\textbf{sceal} — pres. ind. 1/3 sg.: see \textbf{sculan}.\par
% glossary-block:B01171:end
% glossary-block:B01172:begin
\needspace{3\baselineskip}
\textbf{scēap} (n. n.); gen. sg. scēapes; nom. pl. scēap. sheep.\par
% glossary-block:B01172:end
% glossary-block:B01173:begin
\needspace{3\baselineskip}
\textbf{scēaphierde} (n. m.); gen. sg. scēaphierdes; nom. pl. scēaphierdas. shepherd.\par
% glossary-block:B01173:end
% glossary-block:B01174:begin
\needspace{3\baselineskip}
\textbf{scēawere} (n. m.); gen. sg. scēaweres; nom. pl. scēaweras. spy.\par
% glossary-block:B01174:end
% glossary-block:B01175:begin
\needspace{3\baselineskip}
\textbf{scēawian} (v. weak); pret. 1 sg. scēawode; past ptcp. scēawod. watch; observe. Accusative person or event watched, including a journey.\par
% glossary-block:B01175:end
% glossary-block:B01176:begin
\needspace{3\baselineskip}
\textbf{scęncan} (v. weak); pret. 1 sg. scęncte; past ptcp. scęnced. pour out; give to drink. Accusative drink, dative recipient.\par
% glossary-block:B01176:end
% glossary-block:B01177:begin
\needspace{3\baselineskip}
\textbf{sceolde} — pret. sg.: see \textbf{sculan}.\par
% glossary-block:B01177:end
% glossary-block:B01178:begin
\needspace{3\baselineskip}
\textbf{scēowyrhta} (n. m.); gen. sg. scēowyrhtan; nom. pl. scēowyrhtan. shoemaker.\par
% glossary-block:B01178:end
% glossary-block:B01179:begin
\needspace{3\baselineskip}
\textbf{scęððan} (v. weak); pret. 1 sg. scęðede; past ptcp. scęðed. hurt; injure. Dative person harmed.\par
% glossary-block:B01179:end
% glossary-block:B01180:begin
\needspace{3\baselineskip}
\textbf{sciccels} (n. m.). cloak.\par
% glossary-block:B01180:end
% glossary-block:B01181:begin
\needspace{3\baselineskip}
\textbf{scieppan} (v. strong); pret. 1 sg. scēop; pret. pl. scēopon; past ptcp. sceapen. create; shape.\par
% glossary-block:B01181:end
% glossary-block:B01182:begin
\needspace{3\baselineskip}
\textbf{scieppend} (n. m.). Creator.\par
% glossary-block:B01182:end
% glossary-block:B01183:begin
\needspace{3\baselineskip}
\textbf{scīete} (n. f.). linen cloth.\par
% glossary-block:B01183:end
% glossary-block:B01184:begin
\needspace{3\baselineskip}
\textbf{scīnan} (v. strong); pret. 1 sg. scān; pret. pl. scinon; past ptcp. scinen. shine. shining.\par
% glossary-block:B01184:end
% glossary-block:B01185:begin
\needspace{3\baselineskip}
\textbf{scip} (n. n.); gen. sg. scipes; nom. pl. scipu. ship.\par
% glossary-block:B01185:end
% glossary-block:B01186:begin
\needspace{3\baselineskip}
\textbf{sciphęre} (n. m.); gen. sg. sciphęres; nom. pl. sciphęras. fleet; naval force.\par
% glossary-block:B01186:end
% glossary-block:B01187:begin
\needspace{3\baselineskip}
\textbf{scolde} — pret. sg., variant: see \textbf{sculan}.\par
% glossary-block:B01187:end
% glossary-block:B01188:begin
\needspace{3\baselineskip}
\textbf{Scottas} (ethnonym). Scots.\par
% glossary-block:B01188:end
% glossary-block:B01189:begin
\needspace{3\baselineskip}
\textbf{scrȳdan} (v. weak); pret. 1 sg. scrȳdde; past ptcp. scrȳded. clothe; dress. Accusative person clothed in this introduced frame. dress oneself. Supplied reflexive use.\par
% glossary-block:B01189:end
% glossary-block:B01190:begin
\needspace{3\baselineskip}
\textbf{sculan} (v. preterite-present); pres. 1/3 sg. sceal; pres. pl. sculon; pret. 1 sg. sceolde. must; be obliged to. should; be to. Supplied question about what events were to mean; retain its reported/interpretive force. shall; be yet to come. Supplied gīet sculon fīf, with be explicitly omitted. Conventional citation label; no usable infinitive is supplied. See the finite forms in Unit 6. No past-participle citation is given in the checked source paradigm.\par
% glossary-block:B01190:end
% glossary-block:B01191:begin
\needspace{3\baselineskip}
\textbf{scyldig} (adj.). guilty; culpable.\par
% glossary-block:B01191:end
% glossary-block:B01192:begin
\needspace{3\baselineskip}
\textbf{sē} — nom. m sg.: see \textbf{se, sēo, ðæt}.\par
% glossary-block:B01192:end
% glossary-block:B01193:begin
\needspace{3\baselineskip}
\textbf{sealde} — pret. sg.: see \textbf{sęllan}.\par
% glossary-block:B01193:end
% glossary-block:B01194:begin
\needspace{3\baselineskip}
\textbf{sealf} (n. f.); gen. sg. sealfe. ointment.\par
% glossary-block:B01194:end
% glossary-block:B01195:begin
\needspace{3\baselineskip}
\textbf{sēamere} (n. m.); gen. sg. sēameres; nom. pl. sēameras. tailor.\par
% glossary-block:B01195:end
% glossary-block:B01196:begin
\needspace{3\baselineskip}
\textbf{sēcan} (v. weak); pret. 1 sg. sōhte; past ptcp. (ge)sōht. seek. Accusative person or thing sought.\par
% glossary-block:B01196:end
% glossary-block:B01197:begin
\needspace{3\baselineskip}
\textbf{sęcgan} (v. weak); pret. 1 sg. sægde; past ptcp. (ge)sægd; pres. 3 sg. sægð. say. tell; say. Dative addressee and a reported statement or question as content. say. Supplied source preterite. say; tell. call; declare to be. Object with the stated agreeing predicate.\par
% glossary-block:B01197:end
% glossary-block:B01198:begin
\needspace{3\baselineskip}
\textbf{sęgð} — pres. ind. 3 sg., variant: see \textbf{sęcgan}.\par
% glossary-block:B01198:end
% glossary-block:B01199:begin
\needspace{3\baselineskip}
\textbf{sęlen} (n. f.). bounty; bestowal.\par
% glossary-block:B01199:end
% glossary-block:B01200:begin
\needspace{3\baselineskip}
\textbf{self} (adj.). self; the same person or thing. Agreeing emphatic/reflexive self; ordinary reflexive personal pronouns do not always require it.\par
% glossary-block:B01200:end
% glossary-block:B01201:begin
\needspace{3\baselineskip}
\textbf{sęllan} (v. weak); pret. 1 sg. sealde; past ptcp. (ge)seald. give. give to drink. Supplied dative person + plain drincan, without tō. sell. Accusative thing sold.\par
% glossary-block:B01201:end
% glossary-block:B01202:begin
\needspace{3\baselineskip}
\textbf{sellic} (adj.). strange; remarkable.\par
% glossary-block:B01202:end
% glossary-block:B01203:begin
\needspace{3\baselineskip}
\textbf{sęmninga} (adv.). suddenly.\par
% glossary-block:B01203:end
% glossary-block:B01204:begin
\needspace{3\baselineskip}
\textbf{sęndan} (v. weak); pret. 1 sg. sęnde; past ptcp. (ge)sęnd(ed). send. Accusative person or thing sent, dative recipient.\par
% glossary-block:B01204:end
% glossary-block:B01205:begin
\needspace{3\baselineskip}
\textbf{sengan} (v.). singe. Pronunciation example.\par
% glossary-block:B01205:end
% glossary-block:B01206:begin
\needspace{3\baselineskip}
\textbf{sēo} — nom. f sg.: see \textbf{se, sēo, ðæt}; pres. ind. 1 sg.; pres. subj. sg.: see \textbf{sēon}.\par
% glossary-block:B01206:end
% glossary-block:B01207:begin
\needspace{3\baselineskip}
\textbf{seofon} (cardinal numeral). seven. Cardinal 7.\par
% glossary-block:B01207:end
% glossary-block:B01208:begin
\needspace{3\baselineskip}
\textbf{seofoða} (ordinal). seventh.\par
% glossary-block:B01208:end
% glossary-block:B01209:begin
\needspace{3\baselineskip}
\textbf{seolfor} (n. n.); gen. sg. seolfres; nom. pl. seolfor. silver. Metal noun.\par
% glossary-block:B01209:end
% glossary-block:B01210:begin
\needspace{3\baselineskip}
\textbf{seolfren} (adj.). silver; made of silver. Material adjective, distinct from the noun seolfor.\par
% glossary-block:B01210:end
% glossary-block:B01211:begin
\needspace{3\baselineskip}
\textbf{seolh} (n. m.); gen. sg. sēoles; nom. pl. sēolas. seal.\par
% glossary-block:B01211:end
% glossary-block:B01212:begin
\needspace{3\baselineskip}
\textbf{sēon} (v. strong); pret. 1 sg. seah; pret. pl. sāwon; past ptcp. sewen; pres. 3 sg. siehð. see.\par
% glossary-block:B01212:end
% glossary-block:B01213:begin
\needspace{3\baselineskip}
\textbf{sēoðan} (v. strong); pret. 1 sg. sēað; pret. pl. sudon; past ptcp. soden; pres. 3 sg. sīeð. boil.\par
% glossary-block:B01213:end
% glossary-block:B01214:begin
\needspace{3\baselineskip}
\textbf{se, sēo, ðæt} (article; demonstrative; relative pron.). the. that; those. those. Supplied substantival dative plural referring to waters. what; that which. who; which; that. Inflected relative; case follows its clause role. so that. Only in supplied tō ðon ðætte result connection. after that. Only in temporal æfter ðon. so; to that degree. Only in supplied tō ðan … ðæt.\par
% glossary-block:B01214:end
% glossary-block:B01215:begin
\needspace{3\baselineskip}
\textbf{setl} (n. n.); gen. sg. setles; nom. pl. setl. seat; throne.\par
% glossary-block:B01215:end
% glossary-block:B01216:begin
\needspace{3\baselineskip}
\textbf{sęttan} (v. weak); pret. 1 sg. sętte; past ptcp. (ge)sętt. set; place. appoint; establish. With accusative entity/decree in the assigned frame.\par
% glossary-block:B01216:end
% glossary-block:B01217:begin
\needspace{3\baselineskip}
\textbf{se-ðe} (relative). which; that. Supplied se-ðe relative complex referring to the storm.\par
% glossary-block:B01217:end
% glossary-block:B01218:begin
\needspace{3\baselineskip}
\textbf{sibb} (n. f.); gen. sg. sibbe; nom. pl. sibba. peace.\par
% glossary-block:B01218:end
% glossary-block:B01219:begin
\needspace{3\baselineskip}
\textbf{sīde} (n. f.); gen. sg. sīdan; nom. pl. sīdan. silk.\par
% glossary-block:B01219:end
% glossary-block:B01220:begin
\needspace{3\baselineskip}
\textbf{sīe} — pres. subj. sg.: see \textbf{wesan}.\par
% glossary-block:B01220:end
% glossary-block:B01221:begin
\needspace{3\baselineskip}
\textbf{siehð} — pres. ind. 3 sg.: see \textbf{sēon}.\par
% glossary-block:B01221:end
% glossary-block:B01222:begin
\needspace{3\baselineskip}
\textbf{sīeð} — pres. ind. 3 sg.: see \textbf{sēoðan}.\par
% glossary-block:B01222:end
% glossary-block:B01223:begin
\needspace{3\baselineskip}
\textbf{siex} (cardinal numeral). six. Cardinal 6.\par
% glossary-block:B01223:end
% glossary-block:B01224:begin
\needspace{3\baselineskip}
\textbf{siexta} (ordinal numeral). sixth.\par
% glossary-block:B01224:end
% glossary-block:B01225:begin
\needspace{3\baselineskip}
\textbf{siextig} (cardinal numeral). sixty. Within the supplied 365-days expression.\par
% glossary-block:B01225:end
% glossary-block:B01226:begin
\needspace{3\baselineskip}
\textbf{sige} (n. m.). victory.\par
% glossary-block:B01226:end
% glossary-block:B01227:begin
\needspace{3\baselineskip}
\textbf{simle} (adv.). always.\par
% glossary-block:B01227:end
% glossary-block:B01228:begin
\needspace{3\baselineskip}
\textbf{sīn} (possessive adjective). his; her; its (own). reflexive possessive; mostly poetic in Wright; reference only.\par
% glossary-block:B01228:end
% glossary-block:B01229:begin
\needspace{3\baselineskip}
\textbf{sind} — pres. ind. pl.: see \textbf{wesan}.\par
% glossary-block:B01229:end
% glossary-block:B01230:begin
\needspace{3\baselineskip}
\textbf{sindon} — pres. ind. pl., variant: see \textbf{wesan}.\par
% glossary-block:B01230:end
% glossary-block:B01231:begin
\needspace{3\baselineskip}
\textbf{singāl} (adj.). continual; constant.\par
% glossary-block:B01231:end
% glossary-block:B01232:begin
\needspace{3\baselineskip}
\textbf{singan} (v. strong); pret. 1 sg. sang; pret. pl. sungon; past ptcp. sungen. sing. Intransitive in the introduced frame. sing. Source genitive plural present participle, of people singing.\par
% glossary-block:B01232:end
% glossary-block:B01233:begin
\needspace{3\baselineskip}
\textbf{sint} — pres. ind. pl., variant: see \textbf{wesan}.\par
% glossary-block:B01233:end
% glossary-block:B01234:begin
\needspace{3\baselineskip}
\textbf{sittan} (v. strong); pret. 1 sg. sæt; pret. pl. sǣton; past ptcp. seten. sit.\par
% glossary-block:B01234:end
% glossary-block:B01235:begin
\needspace{3\baselineskip}
\textbf{sīð} (n. m.); gen. sg. sīðes; nom. pl. sīðas. time; occasion.\par
% glossary-block:B01235:end
% glossary-block:B01236:begin
\needspace{3\baselineskip}
\textbf{sīðfæt} (n. m.); gen. sg. sīðfætes; nom. pl. sīðfatas. journey; passage.\par
% glossary-block:B01236:end
% glossary-block:B01237:begin
\needspace{3\baselineskip}
\textbf{siððan} (adv.; conj.). since; after; afterwards. Conjunction since/after versus adverb afterwards.\par
% glossary-block:B01237:end
% glossary-block:B01238:begin
\needspace{3\baselineskip}
\textbf{slēan} (v. strong); pret. 1 sg. slōg; pret. pl. slōgon; past ptcp. slægen; pres. 3 sg. sliehð. strike; smite. strike; smite. Objectless striking in the assigned hammer sentences. slay; kill. Accusative person killed. make a signal. Only in supplied slōg tācen. kill. Accusative person; the bounded negative command uses selected short-vowel imperative sleah. strike. Supplied past source form.\par
% glossary-block:B01238:end
% glossary-block:B01239:begin
\needspace{3\baselineskip}
\textbf{slęcg} (n. f.); gen. sg. slęcge; nom. pl. slęcga. hammer.\par
% glossary-block:B01239:end
% glossary-block:B01240:begin
\needspace{3\baselineskip}
\textbf{sliehð} — pres. ind. 3 sg.: see \textbf{slēan}.\par
% glossary-block:B01240:end
% glossary-block:B01241:begin
\needspace{3\baselineskip}
\textbf{smēagan} (v.). ponder; consider. Supplied preterite source form.\par
% glossary-block:B01241:end
% glossary-block:B01242:begin
\needspace{3\baselineskip}
\textbf{smierwan} (v. weak); pret. 1 sg. smierede; past ptcp. smiered. anoint. Accusative body part or person, mid + dative material.\par
% glossary-block:B01242:end
% glossary-block:B01243:begin
\needspace{3\baselineskip}
\textbf{smið} (n. m.); gen. sg. smiðes; nom. pl. smiðas. blacksmith.\par
% glossary-block:B01243:end
% glossary-block:B01244:begin
\needspace{3\baselineskip}
\textbf{smiððe} (n. f.); gen. sg. smiððan; nom. pl. smiððan. smithy.\par
% glossary-block:B01244:end
% glossary-block:B01245:begin
\needspace{3\baselineskip}
\textbf{smylte} (adj.); supplied superl. acc. m. sg. smyltestan. calm; smooth.\par
% glossary-block:B01245:end
% glossary-block:B01246:begin
\needspace{3\baselineskip}
\textbf{smyltest} (adj.). very calm; calmest. Supplied superlative adjective describing the sea. Superlative of smylte; source form smyltestan.\par
% glossary-block:B01246:end
% glossary-block:B01247:begin
\needspace{3\baselineskip}
\textbf{smyltnes} (n. f.); gen. sg. smyltnesse. calm; calmness.\par
% glossary-block:B01247:end
% glossary-block:B01248:begin
\needspace{3\baselineskip}
\textbf{snāw} (n. m.). snow.\par
% glossary-block:B01248:end
% glossary-block:B01249:begin
\needspace{3\baselineskip}
\textbf{snelnes} (n. f.). skill; agility.\par
% glossary-block:B01249:end
% glossary-block:B01250:begin
\needspace{3\baselineskip}
\textbf{snīðan} (v. strong); pret. 1 sg. snāð; pret. pl. snidon; past ptcp. sniden. cut.\par
% glossary-block:B01250:end
% glossary-block:B01251:begin
\needspace{3\baselineskip}
\textbf{sōhte} — pret. sg.: see \textbf{sēcan}.\par
% glossary-block:B01251:end
% glossary-block:B01252:begin
\needspace{3\baselineskip}
\textbf{sōna} (adv.). as soon as. In sōna swā. immediately; at once.\par
% glossary-block:B01252:end
% glossary-block:B01253:begin
\needspace{3\baselineskip}
\textbf{sorg} (n. f.). distress; anxiety; trouble; sorrow.\par
% glossary-block:B01253:end
% glossary-block:B01254:begin
\needspace{3\baselineskip}
\textbf{sōð} (adj.). true.\par
% glossary-block:B01254:end
% glossary-block:B01255:begin
\needspace{3\baselineskip}
\textbf{sōðfæstnes} (n. f.); gen. sg. sōðfæstnesse; nom. pl. sōðfæstnessa. truth.\par
% glossary-block:B01255:end
% glossary-block:B01256:begin
\needspace{3\baselineskip}
\textbf{sōðlice} (adv.). now; but. Discourse connective corresponding to source Latin autem, distinct from truly. truly; indeed.\par
% glossary-block:B01256:end
% glossary-block:B01257:begin
\needspace{3\baselineskip}
\textbf{spǣtan} (v. weak); pret. 1 sg. spǣtte; past ptcp. spǣtt. spit.\par
% glossary-block:B01257:end
% glossary-block:B01258:begin
\needspace{3\baselineskip}
\textbf{spell} (n. n.). account; discourse.\par
% glossary-block:B01258:end
% glossary-block:B01259:begin
\needspace{3\baselineskip}
\textbf{sprǣc} (n. f.); gen. sg. sprǣce; nom. pl. sprǣca. saying; speech.\par
% glossary-block:B01259:end
% glossary-block:B01260:begin
\needspace{3\baselineskip}
\textbf{sprecan} (v. strong); pret. 1 sg. spræc; pret. pl. sprǣcon; past ptcp. sprecen; pres. 3 sg. spricð. speak. speak with. mid + dative interlocutor in the taught production frame.\par
% glossary-block:B01260:end
% glossary-block:B01261:begin
\needspace{3\baselineskip}
\textbf{spricð} — pres. ind. 3 sg.: see \textbf{sprecan}.\par
% glossary-block:B01261:end
% glossary-block:B01262:begin
\needspace{3\baselineskip}
\textbf{stān} (n. m.); gen. sg. stānes; nom. pl. stānas. stone.\par
% glossary-block:B01262:end
% glossary-block:B01263:begin
\needspace{3\baselineskip}
\textbf{standan} (v. strong); pret. 1 sg. stōd; pret. pl. stōdon; past ptcp. standen; pres. 3 sg. stęnt. oppose; stand against. Supplied standan ongēan frame, distinct from simple stand. stand.\par
% glossary-block:B01263:end
% glossary-block:B01264:begin
\needspace{3\baselineskip}
\textbf{stānen} (adj.). stone; made of stone. Material adjective.\par
% glossary-block:B01264:end
% glossary-block:B01265:begin
\needspace{3\baselineskip}
\textbf{stānweall} (n. m.); gen. sg. stānwealles; nom. pl. stānweallas. stone wall.\par
% glossary-block:B01265:end
% glossary-block:B01266:begin
\needspace{3\baselineskip}
\textbf{stefn} (n. f.); gen. sg. stefne; nom. pl. stefna. voice.\par
% glossary-block:B01266:end
% glossary-block:B01267:begin
\needspace{3\baselineskip}
\textbf{stelan} (v.). steal.\par
% glossary-block:B01267:end
% glossary-block:B01268:begin
\needspace{3\baselineskip}
\textbf{stęnt} — pres. ind. 3 sg.: see \textbf{standan}.\par
% glossary-block:B01268:end
% glossary-block:B01269:begin
\needspace{3\baselineskip}
\textbf{steorra} (n. m.); gen. sg. steorran; nom. pl. steorran. star.\par
% glossary-block:B01269:end
% glossary-block:B01270:begin
\needspace{3\baselineskip}
\textbf{stille} (adj.). still; quiet.\par
% glossary-block:B01270:end
% glossary-block:B01271:begin
\needspace{3\baselineskip}
\textbf{stilnes} (n. f.). stillness; calm. Supplied genitive singular. stillness; quietness.\par
% glossary-block:B01271:end
% glossary-block:B01272:begin
\needspace{3\baselineskip}
\textbf{storm} (n. m.). storm.\par
% glossary-block:B01272:end
% glossary-block:B01273:begin
\needspace{3\baselineskip}
\textbf{strand} (n. m.); gen. sg. strandes; nom. pl. strandas. strand; sea-shore.\par
% glossary-block:B01273:end
% glossary-block:B01274:begin
\needspace{3\baselineskip}
\textbf{strang} (adj.). strong. Masculine singular predicate in the source. strong. Attributive adjective with wind in the source.\par
% glossary-block:B01274:end
% glossary-block:B01275:begin
\needspace{3\baselineskip}
\textbf{strǣt} (n. f.); gen. sg. strǣte; nom. pl. strǣta. street. public place. Street used as a gathering place in the assigned prompt.\par
% glossary-block:B01275:end
% glossary-block:B01276:begin
\needspace{3\baselineskip}
\textbf{strēaw} (n. n.). straw.\par
% glossary-block:B01276:end
% glossary-block:B01277:begin
\needspace{3\baselineskip}
\textbf{stręngðu} (n. f.). strength.\par
% glossary-block:B01277:end
% glossary-block:B01278:begin
\needspace{3\baselineskip}
\textbf{strengu} (n. f.); gen. sg. strenge; nom. pl. strenga. strength. Recognition class-reference noun; preserve the source spelling and do not infer a plural count meaning.\par
% glossary-block:B01278:end
% glossary-block:B01279:begin
\needspace{3\baselineskip}
\textbf{styccemǣlum} (adv.). little by little.\par
% glossary-block:B01279:end
% glossary-block:B01280:begin
\needspace{3\baselineskip}
\textbf{styriendlic} (adj.). moving.\par
% glossary-block:B01280:end
% glossary-block:B01281:begin
\needspace{3\baselineskip}
\textbf{sulhscear} (n. m.); gen. sg. sulhsceares; nom. pl. sulhscearas. plowshare.\par
% glossary-block:B01281:end
% glossary-block:B01282:begin
\needspace{3\baselineskip}
\textbf{sum} (indefinite adjective/pronoun). a certain; some. some. Members selected from a group; genitive whole.\par
% glossary-block:B01282:end
% glossary-block:B01283:begin
\needspace{3\baselineskip}
\textbf{sumer} (n. m.); gen. sg. sumeres; nom. pl. sumeras. summer.\par
% glossary-block:B01283:end
% glossary-block:B01284:begin
\needspace{3\baselineskip}
\textbf{sund} (n.). swimming.\par
% glossary-block:B01284:end
% glossary-block:B01285:begin
\needspace{3\baselineskip}
\textbf{sunne} (n. f.); gen. sg. sunnan; nom. pl. sunnan. sun.\par
% glossary-block:B01285:end
% glossary-block:B01286:begin
\needspace{3\baselineskip}
\textbf{sunu} (n. m.); gen. sg. suna; nom. pl. suna. son.\par
% glossary-block:B01286:end
% glossary-block:B01287:begin
\needspace{3\baselineskip}
\textbf{sūðan} (adv.). in the south; from the south. Retain the supplied spatial expression.\par
% glossary-block:B01287:end
% glossary-block:B01288:begin
\needspace{3\baselineskip}
\textbf{sūðerne} (adj.). southern.\par
% glossary-block:B01288:end
% glossary-block:B01289:begin
\needspace{3\baselineskip}
\textbf{sūðwesterne} (adj.). southwestern.\par
% glossary-block:B01289:end
% glossary-block:B01290:begin
\needspace{3\baselineskip}
\textbf{swā} (adverb; conjunction). so; thus. as. Only in sōna swā. as long as; inasmuch as. Whole swā lange swā phrase; the assigned biblical expression need not imply duration. as; in the same way as. as. First degree-marker in equal comparison. as. Introduces the comparison standard. as; like. Predicative comparison. the … the. Proportional comparison, not the definite article. whoever; whatever. swā hwā/hwæt/hwæðer swā; the internal pronoun's case and type must fit the clause.\par
% glossary-block:B01290:end
% glossary-block:B01291:begin
\needspace{3\baselineskip}
\textbf{swā-hwæt-swā} (pron.). whatsoever; whatever.\par
% glossary-block:B01291:end
% glossary-block:B01292:begin
\needspace{3\baselineskip}
\textbf{swā-swā} (conj.). just as. about; approximately. Approximation in the supplied source expression.\par
% glossary-block:B01292:end
% glossary-block:B01293:begin
\needspace{3\baselineskip}
\textbf{swēg} (n. m.). music; sound.\par
% glossary-block:B01293:end
% glossary-block:B01294:begin
\needspace{3\baselineskip}
\textbf{swēgcræft} (n. m.). musical skill.\par
% glossary-block:B01294:end
% glossary-block:B01295:begin
\needspace{3\baselineskip}
\textbf{sweostor} (n. f.); gen. sg. sweostor; nom. pl. sweostor. sister.\par
% glossary-block:B01295:end
% glossary-block:B01296:begin
\needspace{3\baselineskip}
\textbf{sweotole} (adv.); superl. sweotolost. In superlative sweotolost: most clearly.\par
% glossary-block:B01296:end
% glossary-block:B01297:begin
\needspace{3\baselineskip}
\textbf{sweotollice} (adv.). clearly.\par
% glossary-block:B01297:end
% glossary-block:B01298:begin
\needspace{3\baselineskip}
\textbf{sweotolost} — superlative adverb; exact source spelling: see \textbf{sweotole}.\par
% glossary-block:B01298:end
% glossary-block:B01299:begin
\needspace{3\baselineskip}
\textbf{swęrian} (v. strong); pret. 1 sg. swōr; pret. pl. swōron; past ptcp. sworen. swear.\par
% glossary-block:B01299:end
% glossary-block:B01300:begin
\needspace{3\baselineskip}
\textbf{swēte} (adj.); comp. swētra; superl. swētest. sweet.\par
% glossary-block:B01300:end
% glossary-block:B01301:begin
\needspace{3\baselineskip}
\textbf{swētest} (adj.). sweetest. Superlative of swēte.\par
% glossary-block:B01301:end
% glossary-block:B01302:begin
\needspace{3\baselineskip}
\textbf{swētnes} (n. f.). sweetness.\par
% glossary-block:B01302:end
% glossary-block:B01303:begin
\needspace{3\baselineskip}
\textbf{swētra} (adj.). sweeter. Comparative of swēte.\par
% glossary-block:B01303:end
% glossary-block:B01304:begin
\needspace{3\baselineskip}
\textbf{swift} (adj.). swift.\par
% glossary-block:B01304:end
% glossary-block:B01305:begin
\needspace{3\baselineskip}
\textbf{swiftnes} (n. f.). swiftness.\par
% glossary-block:B01305:end
% glossary-block:B01306:begin
\needspace{3\baselineskip}
\textbf{swīge} (n. f.). silence.\par
% glossary-block:B01306:end
% glossary-block:B01307:begin
\needspace{3\baselineskip}
\textbf{swigian} (v.). be silent. Supplied source sense.\par
% glossary-block:B01307:end
% glossary-block:B01308:begin
\needspace{3\baselineskip}
\textbf{swilc} (indefinite adjective). such.\par
% glossary-block:B01308:end
% glossary-block:B01309:begin
\needspace{3\baselineskip}
\textbf{swilce} (adv.). likewise. Adverb, distinct from inflected swilc such. about; approximately. as if. Supplied apparent-rejuvenation comparison.\par
% glossary-block:B01309:end
% glossary-block:B01310:begin
\needspace{3\baselineskip}
\textbf{swimman} (v. strong); pret. 1 sg. swamm; pret. pl. swummon; past ptcp. swummen. swim.\par
% glossary-block:B01310:end
% glossary-block:B01311:begin
\needspace{3\baselineskip}
\textbf{swingan} (v.). swing; brandish. Supplied past source form.\par
% glossary-block:B01311:end
% glossary-block:B01312:begin
\needspace{3\baselineskip}
\textbf{swīð} (adjective citation); comp. swīðra. strong.\par
% glossary-block:B01312:end
% glossary-block:B01313:begin
\needspace{3\baselineskip}
\textbf{swīðe} (adv.); comp. swīðor; superl. swīðost. very. rather; more. Preferential comparative swīðor. most; very greatly. Superlative swīðost; recognition reference, not independent production. more. Comparative degree of an action/feeling, not a count of objects.\par
% glossary-block:B01313:end
% glossary-block:B01314:begin
\needspace{3\baselineskip}
\textbf{swīðlice} (adv.); comp. swīðlicor. In comparative swīðlicor: more violently; more severely.\par
% glossary-block:B01314:end
% glossary-block:B01315:begin
\needspace{3\baselineskip}
\textbf{swīðlicor} — comparative adverb; exact source spelling: see \textbf{swīðlice}.\par
% glossary-block:B01315:end
% glossary-block:B01316:begin
\needspace{3\baselineskip}
\textbf{swīðor} — comparative adverb: see \textbf{swīðe}.\par
% glossary-block:B01316:end
% glossary-block:B01317:begin
\needspace{3\baselineskip}
\textbf{swīðost} — superlative adverb: see \textbf{swīðe}.\par
% glossary-block:B01317:end
% glossary-block:B01318:begin
\needspace{3\baselineskip}
\textbf{swīðra} (adj.). right. Right-hand/side expression, distinct from stronger. stronger; fiercer. Comparative degree, distinct from right-side sense. In stronger/fiercer, comparative of swīð; right-hand/side is separate.\par
% glossary-block:B01318:end
% glossary-block:B01319:begin
\needspace{3\baselineskip}
\textbf{synfull} (adj.). sinful.\par
% glossary-block:B01319:end
% glossary-block:B01320:begin
\needspace{3\baselineskip}
\textbf{synn} (n. f.); gen. sg. synne; nom. pl. synna. sin.\par
% glossary-block:B01320:end
\pdfbookmark[1]{T}{outline-B01321}
% glossary-block:B01321:begin
\needspace{5\baselineskip}
\subsection*{T}
% glossary-block:B01321:end
% glossary-block:B01322:begin
\textbf{tabele} (n. f.); gen. sg. tabelan; nom. pl. tabelan. tablet; table. Tables of the law, not an unrestricted furniture meaning.\par
% glossary-block:B01322:end
% glossary-block:B01323:begin
\needspace{3\baselineskip}
\textbf{tācen} (n. n.); gen. sg. tācnes; nom. pl. tācnu. sign.\par
% glossary-block:B01323:end
% glossary-block:B01324:begin
\needspace{3\baselineskip}
\textbf{tǣcan} (v.). teach; show.\par
% glossary-block:B01324:end
% glossary-block:B01325:begin
\needspace{3\baselineskip}
\textbf{tǣlan} (v. weak); pret. 1 sg. tǣlde; past ptcp. tǣled. blame; censure.\par
% glossary-block:B01325:end
% glossary-block:B01326:begin
\needspace{3\baselineskip}
\textbf{tēar} (n. m.); gen. sg. tēares; nom. pl. tēaras. tear.\par
% glossary-block:B01326:end
% glossary-block:B01327:begin
\needspace{3\baselineskip}
\textbf{tempel} (n. n.); gen. sg. temples; nom. pl. templu. temple.\par
% glossary-block:B01327:end
% glossary-block:B01328:begin
\needspace{3\baselineskip}
\textbf{tēon} (v. strong); pret. 1 sg. tēah; pret. pl. tugon; past ptcp. togen; pres. 3 sg. tīehð. draw; pull. bring forth; produce. Supplied third-plural present subjunctive, with waters as subject. bring forth; produce. Accusative thing produced with particle forð; use the assigned tēon forms.\par
% glossary-block:B01328:end
% glossary-block:B01329:begin
\needspace{3\baselineskip}
\textbf{tēð} — dat. sg.; nom./acc. pl.: see \textbf{tōð}.\par
% glossary-block:B01329:end
% glossary-block:B01330:begin
\needspace{3\baselineskip}
\textbf{tīd} (n. f.); gen. sg. tīde; nom. pl. tīda. season; time. hour. Time measure in the supplied intervals.\par
% glossary-block:B01330:end
% glossary-block:B01331:begin
\needspace{3\baselineskip}
\textbf{tīehð} — pres. ind. 3 sg.: see \textbf{tēon}.\par
% glossary-block:B01331:end
% glossary-block:B01332:begin
\needspace{3\baselineskip}
\textbf{tīen} (cardinal numeral). ten. Cardinal 10.\par
% glossary-block:B01332:end
% glossary-block:B01333:begin
\needspace{3\baselineskip}
\textbf{tilian} (v.). obtain; provide. Genitive thing procured; dative reflexive beneficiary in the assigned frame.\par
% glossary-block:B01333:end
% glossary-block:B01334:begin
\needspace{3\baselineskip}
\textbf{tin} (n. n.); gen. sg. tines; nom. pl. tinu. tin.\par
% glossary-block:B01334:end
% glossary-block:B01335:begin
\needspace{3\baselineskip}
\textbf{tō} (prep.). to; in order to. With the inflected infinitive; ordinary modal complements use a plain infinitive. for; to serve as. tō + dative function; a separate dative beneficiary may accompany it. to. Endpoint/result death. to. Dative destination. at; in response to. Only the exact at-the-prayers expression. to; toward. Person to whom knees/prayer are directed in the supplied phrase. so that. Only in supplied tō ðon ðætte result connection. for aid; for protection. Whole supplied purpose expressions. to. Postposed in the supplied cwæð him tō speech phrase. at; to. Bounded supplied root expression with genitive dependent and dative root. to. The person addressed in speech. to; into. Result in supplied tō lore, tō sorge or tō nāhte. why; for what purpose. Only in tō hwon. into; to. Resulting language/literary form. toward. Conduct toward a person in the supplied expression. in accordance with. A will as the goal of guidance in the supplied prayer. for; to the benefit of. for ever. Whole fixed tō worulde. for. Care of a body part in the supplied expression. as wife. Only the supplied tō wīfe role phrase. to. Complement of belonging in the supplied frame. to. Resulting condition of youth in the supplied comparison. so; to that degree. Whole supplied tō ðan … ðæt result construction. toward. Target of a look in the supplied frame.\par
% glossary-block:B01335:end
% glossary-block:B01336:begin
\needspace{3\baselineskip}
\textbf{tōberstan} (v. strong); pret. 1 sg. tōbærst; pret. pl. tōburston; past ptcp. tōborsten; pres. 3 sg. tōbierst. break up; go to pieces.\par
% glossary-block:B01336:end
% glossary-block:B01337:begin
\needspace{3\baselineskip}
\textbf{tōbierst} — pres. ind. 3 sg.: see \textbf{tōberstan}.\par
% glossary-block:B01337:end
% glossary-block:B01338:begin
\needspace{3\baselineskip}
\textbf{tōcuman} (v.). come; arrive.\par
% glossary-block:B01338:end
% glossary-block:B01339:begin
\needspace{3\baselineskip}
\textbf{tō-dæg} (adv.). today.\par
% glossary-block:B01339:end
% glossary-block:B01340:begin
\needspace{3\baselineskip}
\textbf{tōdæg} (adv.). today.\par
% glossary-block:B01340:end
% glossary-block:B01341:begin
\needspace{3\baselineskip}
\textbf{tōdǣlen / tōdǣlden (source forms)} (verb; supplied inflected forms). separate; distinguish. Source forms tōdǣlen/tōdǣlden are retained; no silently invented citation spelling. Only the cited source forms are supplied here.\par
% glossary-block:B01341:end
% glossary-block:B01342:begin
\needspace{3\baselineskip}
\textbf{tōforan} (prep.). before; above all.\par
% glossary-block:B01342:end
% glossary-block:B01343:begin
\needspace{3\baselineskip}
\textbf{tōgædere} (adv.). together.\par
% glossary-block:B01343:end
% glossary-block:B01344:begin
\needspace{3\baselineskip}
\textbf{tōgēanes} (postposition). towards; to meet.\par
% glossary-block:B01344:end
% glossary-block:B01345:begin
\needspace{3\baselineskip}
\textbf{tōgēiecan} (v.). add.\par
% glossary-block:B01345:end
% glossary-block:B01346:begin
\needspace{3\baselineskip}
\textbf{tōgelǣdan} (v. weak); pret. 1 sg. tōgelǣdde; past ptcp. tōgelǣded. bring. Accusative thing brought, dative recipient.\par
% glossary-block:B01346:end
% glossary-block:B01347:begin
\needspace{3\baselineskip}
\textbf{tōhopa} (n. m.). hope.\par
% glossary-block:B01347:end
% glossary-block:B01348:begin
\needspace{3\baselineskip}
\textbf{tōmiddes} (prep.). in the midst of. With dative landmark.\par
% glossary-block:B01348:end
% glossary-block:B01349:begin
\needspace{3\baselineskip}
\textbf{tō-morgen} (adv.). tomorrow.\par
% glossary-block:B01349:end
% glossary-block:B01350:begin
\needspace{3\baselineskip}
\textbf{top} (n. m.). ball?. Tentative source interpretation; retain the question mark and uncertainty.\par
% glossary-block:B01350:end
% glossary-block:B01351:begin
\needspace{3\baselineskip}
\textbf{torr} (n. m.); gen. sg. torres; nom. pl. torras. tower. crag. Probable contextual sense in Cook's note, distinct from earlier tower.\par
% glossary-block:B01351:end
% glossary-block:B01352:begin
\needspace{3\baselineskip}
\textbf{tōscǣt} — pres. ind. 3 sg.: see \textbf{tōsceadan}.\par
% glossary-block:B01352:end
% glossary-block:B01353:begin
\needspace{3\baselineskip}
\textbf{tōsceadan} (v.); pres. 3 sg. tōscǣt. separate; divide.\par
% glossary-block:B01353:end
% glossary-block:B01354:begin
\needspace{3\baselineskip}
\textbf{tō-slītan} (v.). tear. Supplied past source form.\par
% glossary-block:B01354:end
% glossary-block:B01355:begin
\needspace{3\baselineskip}
\textbf{tōtwǣman} (v. weak); pret. 1 sg. tōtwǣmde; past ptcp. tōtwǣmed. divide; separate. Waters divided from waters in the supplied frame. separate. Accusative beings separated; fram + dative other party.\par
% glossary-block:B01355:end
% glossary-block:B01356:begin
\needspace{3\baselineskip}
\textbf{tōð} (n. m.); gen. sg. tōðes; nom. pl. tēð. tooth.\par
% glossary-block:B01356:end
% glossary-block:B01357:begin
\needspace{3\baselineskip}
\textbf{tōweorpan} (v. strong); pret. 1 sg. tōwearp; pret. pl. tōwurpon; past ptcp. tōworpen; pres. 3 sg. tōwierpð. break; overthrow. Break a commandment in the assigned use.\par
% glossary-block:B01357:end
% glossary-block:B01358:begin
\needspace{3\baselineskip}
\textbf{tōwierpð} — pres. ind. 3 sg.: see \textbf{tōweorpan}.\par
% glossary-block:B01358:end
% glossary-block:B01359:begin
\needspace{3\baselineskip}
\textbf{trēo} (n. n.); gen. sg. trēowes; nom. pl. trēowu. tree.\par
% glossary-block:B01359:end
% glossary-block:B01360:begin
\needspace{3\baselineskip}
\textbf{trēownes} (n. f.). trust; confidence.\par
% glossary-block:B01360:end
% glossary-block:B01361:begin
\needspace{3\baselineskip}
\textbf{trēowwyrhta} (n. m.); gen. sg. trēowwyrhtan; nom. pl. trēowwyrhtan. carpenter.\par
% glossary-block:B01361:end
% glossary-block:B01362:begin
\needspace{3\baselineskip}
\textbf{tū} — nom./acc. neut. pl.: see \textbf{twēgen}.\par
% glossary-block:B01362:end
% glossary-block:B01363:begin
\needspace{3\baselineskip}
\textbf{twā} — nom./acc. fem./neut. pl.: see \textbf{twēgen}.\par
% glossary-block:B01363:end
% glossary-block:B01364:begin
\needspace{3\baselineskip}
\textbf{twām} — dat. pl.: see \textbf{twēgen}.\par
% glossary-block:B01364:end
% glossary-block:B01365:begin
\needspace{3\baselineskip}
\textbf{twǣm} — dat. pl., variant: see \textbf{twēgen}.\par
% glossary-block:B01365:end
% glossary-block:B01366:begin
\needspace{3\baselineskip}
\textbf{twēgen} (cardinal numeral). two. two; in two. Two parts in the supplied division expression.\par
% glossary-block:B01366:end
% glossary-block:B01367:begin
\needspace{3\baselineskip}
\textbf{twęlf} (cardinal numeral). twelve. Cardinal 12.\par
% glossary-block:B01367:end
% glossary-block:B01368:begin
\needspace{3\baselineskip}
\textbf{twēntig} (numeral). twenty.\par
% glossary-block:B01368:end
% glossary-block:B01369:begin
\needspace{3\baselineskip}
\textbf{twīfeald} (adj.). double; twofold.\par
% glossary-block:B01369:end
\pdfbookmark[1]{Þ / Ð}{outline-B01370}
% glossary-block:B01370:begin
\needspace{5\baselineskip}
\subsection*{Þ / Ð}
% glossary-block:B01370:end
% glossary-block:B01371:begin
\textbf{ðā} (adv.). then. Also acc. f sg.; nom. pl.; acc. pl.: see \textbf{se, sēo, ðæt}.\par
% glossary-block:B01371:end
% glossary-block:B01372:begin
\needspace{3\baselineskip}
\textbf{ðā} (conj.). when; factual past time. Also acc. f sg.; nom. pl.; acc. pl.: see \textbf{se, sēo, ðæt}.\par
% glossary-block:B01372:end
% glossary-block:B01373:begin
\needspace{3\baselineskip}
\textbf{ðā-hwæðre} (conj.). nevertheless; yet. Whole supplied connective.\par
% glossary-block:B01373:end
% glossary-block:B01374:begin
\needspace{3\baselineskip}
\textbf{ðā-hwīle-ðe} (conj.). while.\par
% glossary-block:B01374:end
% glossary-block:B01375:begin
\needspace{3\baselineskip}
\textbf{ðām} — dat. m./n. sg.; dat. pl., variant: see \textbf{se, sēo, ðæt}.\par
% glossary-block:B01375:end
% glossary-block:B01376:begin
\needspace{3\baselineskip}
\textbf{ðan} — in supplied tō ðan, to that degree: see \textbf{se, sēo, ðæt}.\par
% glossary-block:B01376:end
% glossary-block:B01377:begin
\needspace{3\baselineskip}
\textbf{ðanan} (adv.). thence; from there.\par
% glossary-block:B01377:end
% glossary-block:B01378:begin
\needspace{3\baselineskip}
\textbf{ðancung} (n. f.); gen. sg. ðancunge. thanksgiving.\par
% glossary-block:B01378:end
% glossary-block:B01379:begin
\needspace{3\baselineskip}
\textbf{ðāra} — gen. pl.: see \textbf{se, sēo, ðæt}.\par
% glossary-block:B01379:end
% glossary-block:B01380:begin
\needspace{3\baselineskip}
\textbf{ðās} — acc. f sg.; nom. pl.; acc. pl.: see \textbf{ðes, ðēos, ðis}.\par
% glossary-block:B01380:end
% glossary-block:B01381:begin
\needspace{3\baselineskip}
\textbf{ðæce} (n. f.). roof.\par
% glossary-block:B01381:end
% glossary-block:B01382:begin
\needspace{3\baselineskip}
\textbf{ðǣm} — dat. m sg.; dat. n sg.; dat. pl.: see \textbf{se, sēo, ðæt}.\par
% glossary-block:B01382:end
% glossary-block:B01383:begin
\needspace{3\baselineskip}
\textbf{ðǣr} (adv.). there.\par
% glossary-block:B01383:end
% glossary-block:B01384:begin
\needspace{3\baselineskip}
\textbf{ðǣra} — gen. pl., variant: see \textbf{se, sēo, ðæt}.\par
% glossary-block:B01384:end
% glossary-block:B01385:begin
\needspace{3\baselineskip}
\textbf{ðǣre} — gen. f sg.; dat. f sg.: see \textbf{se, sēo, ðæt}.\par
% glossary-block:B01385:end
% glossary-block:B01386:begin
\needspace{3\baselineskip}
\textbf{ðǣron} (adverb or particle). therein; in it.\par
% glossary-block:B01386:end
% glossary-block:B01387:begin
\needspace{3\baselineskip}
\textbf{ðǣr-tō-ēacan} (adverbial expression). in addition to that; besides.\par
% glossary-block:B01387:end
% glossary-block:B01388:begin
\needspace{3\baselineskip}
\textbf{ðæs} — gen. m sg.; gen. n sg.: see \textbf{se, sēo, ðæt}.\par
% glossary-block:B01388:end
% glossary-block:B01389:begin
\needspace{3\baselineskip}
\textbf{ðæs-ðe} (conj.). as soon as. Complete source construction with sōna. as. Only in supplied ðæs-ðe mē ðyncð.\par
% glossary-block:B01389:end
% glossary-block:B01390:begin
\needspace{3\baselineskip}
\textbf{ðæt} (conj.). that. Content conjunction, distinct from demonstrative/relative that. so that; in order that. Purpose conjunction with the assigned subjunctive construction. that; so that. Supplied expanded spelling; distinguish content from result by the source clause. that; so that. Asserted result, distinct from intended purpose. that. Cook leaves the source linkage uncertain; do not impose one content/result classification. Also nom. n sg.; acc. n sg.: see \textbf{se, sēo, ðæt}.\par
% glossary-block:B01390:end
% glossary-block:B01391:begin
\needspace{3\baselineskip}
\textbf{ðæt-ðæt} (pron.). that which; what.\par
% glossary-block:B01391:end
% glossary-block:B01392:begin
\needspace{3\baselineskip}
\textbf{ðe} (relative particle; clause marker). that; which; who. which; that. Subject-relative referring to the waters. after. Only as the clause marker in the whole supplied æfter ðǣm ðe expression. clause marker. Only in supplied ðēah ðe, although; ðe alone is not an equivalent of although.\par
% glossary-block:B01392:end
% glossary-block:B01393:begin
\needspace{3\baselineskip}
\textbf{ðē} — dat. sg.; acc. sg.: see \textbf{ðū}.\par
% glossary-block:B01393:end
% glossary-block:B01394:begin
\needspace{3\baselineskip}
\textbf{ðēah} (conj.; adv.). although; even though. nevertheless; yet. Independent adverb, distinct from concessive conjunction although.\par
% glossary-block:B01394:end
% glossary-block:B01395:begin
\needspace{3\baselineskip}
\textbf{ðearf} (n. f.); gen. sg. ðearfe. benefit; good; need. Source contextual nouns; retain the local expression. Also pres. ind. 1/3 sg.: see \textbf{ðurfan}.\par
% glossary-block:B01395:end
% glossary-block:B01396:begin
\needspace{3\baselineskip}
\textbf{ðearfa} (n. m.); gen. sg. ðearfan; nom. pl. ðearfan. poor person; needy person. Noun referring to a person.\par
% glossary-block:B01396:end
% glossary-block:B01397:begin
\needspace{3\baselineskip}
\textbf{ðearle} (adv.). greatly; exceedingly.\par
% glossary-block:B01397:end
% glossary-block:B01398:begin
\needspace{3\baselineskip}
\textbf{ðearlice} (adv.). violently; fiercely.\par
% glossary-block:B01398:end
% glossary-block:B01399:begin
\needspace{3\baselineskip}
\textbf{ðec} — acc. sg., variant: see \textbf{ðū}.\par
% glossary-block:B01399:end
% glossary-block:B01400:begin
\needspace{3\baselineskip}
\textbf{ðegn} (n. m.); gen. sg. ðegnes; nom. pl. ðegnas. vassal; retainer; thane.\par
% glossary-block:B01400:end
% glossary-block:B01401:begin
\needspace{3\baselineskip}
\textbf{ðęncan} (v. weak); pret. 1 sg. ðōhte; past ptcp. (ge)ðōht. think. Supplied source preterite. think.\par
% glossary-block:B01401:end
% glossary-block:B01402:begin
\needspace{3\baselineskip}
\textbf{ðēnian} (v.). serve; attend. Supplied preterite source form.\par
% glossary-block:B01402:end
% glossary-block:B01403:begin
\needspace{3\baselineskip}
\textbf{ðēnung} (n. f.); gen. sg. ðēnunge; nom. pl. ðēnunga. service; attendance.\par
% glossary-block:B01403:end
% glossary-block:B01404:begin
\needspace{3\baselineskip}
\textbf{ðēod} (n. f.); gen. sg. ðēode; nom. pl. ðēoda. people; nation. region; country.\par
% glossary-block:B01404:end
% glossary-block:B01405:begin
\needspace{3\baselineskip}
\textbf{ðēos} — nom. f sg.: see \textbf{ðes, ðēos, ðis}.\par
% glossary-block:B01405:end
% glossary-block:B01406:begin
\needspace{3\baselineskip}
\textbf{ðēow} (adj.). the unfree; enslaved people. Substantival adjective in the source comparison.\par
% glossary-block:B01406:end
% glossary-block:B01407:begin
\needspace{3\baselineskip}
\textbf{ðēow} (n. m.); gen. sg. ðēowes; nom. pl. ðēowas. servant; enslaved person.\par
% glossary-block:B01407:end
% glossary-block:B01408:begin
\needspace{3\baselineskip}
\textbf{ðēowa} (n. m.); gen. sg. ðēowan; nom. pl. ðēowan. servant.\par
% glossary-block:B01408:end
% glossary-block:B01409:begin
\needspace{3\baselineskip}
\textbf{ðēowa manna} (noun phrase). enslaved men; servants. Whole supplied phrase ðēowa manna.\par
% glossary-block:B01409:end
% glossary-block:B01410:begin
\needspace{3\baselineskip}
\textbf{ðēowdōm} (n. m.); gen. sg. ðēowdōmes; nom. pl. ðēowdōmas. service.\par
% glossary-block:B01410:end
% glossary-block:B01411:begin
\needspace{3\baselineskip}
\textbf{ðēowian} (v.). enslave; reduce to servitude.\par
% glossary-block:B01411:end
% glossary-block:B01412:begin
\needspace{3\baselineskip}
\textbf{ðes, ðēos, ðis} (demonstrative). this; these.\par
% glossary-block:B01412:end
% glossary-block:B01413:begin
\needspace{3\baselineskip}
\textbf{ðider} (adv.). thither; to that place.\par
% glossary-block:B01413:end
% glossary-block:B01414:begin
\needspace{3\baselineskip}
\textbf{ðīestru} (n. n.); nom. pl. ðīestru. darkness. Plural only.\par
% glossary-block:B01414:end
% glossary-block:B01415:begin
\needspace{3\baselineskip}
\textbf{ðīn} (possessive adjective). your; yours. One possessor; agrees with the possessed noun. Also gen. sg.: see \textbf{ðū}.\par
% glossary-block:B01415:end
% glossary-block:B01416:begin
\needspace{3\baselineskip}
\textbf{ðīnen} (n. f.); gen. sg. ðīnenne; nom. pl. ðīnenna. handmaiden.\par
% glossary-block:B01416:end
% glossary-block:B01417:begin
\needspace{3\baselineskip}
\textbf{ðing} (n. n.); gen. sg. ðinges; nom. pl. ðing. thing. Ordinary singular/plural noun. thing; fact; matter. Matter known/discussed, distinct from a generic physical object.\par
% glossary-block:B01417:end
% glossary-block:B01418:begin
\needspace{3\baselineskip}
\textbf{þing} — citation/spelling variant: see \textbf{ðing}.\par
% glossary-block:B01418:end
% glossary-block:B01419:begin
\needspace{3\baselineskip}
\textbf{ðis} — nom. n sg.; acc. n sg.: see \textbf{ðes, ðēos, ðis}.\par
% glossary-block:B01419:end
% glossary-block:B01420:begin
\needspace{3\baselineskip}
\textbf{ðisne} — acc. m sg.: see \textbf{ðes, ðēos, ðis}.\par
% glossary-block:B01420:end
% glossary-block:B01421:begin
\needspace{3\baselineskip}
\textbf{ðissa} — gen. pl.: see \textbf{ðes, ðēos, ðis}.\par
% glossary-block:B01421:end
% glossary-block:B01422:begin
\needspace{3\baselineskip}
\textbf{ðisse} — gen. f sg.; dat. f sg.: see \textbf{ðes, ðēos, ðis}.\par
% glossary-block:B01422:end
% glossary-block:B01423:begin
\needspace{3\baselineskip}
\textbf{ðisses} — gen. m sg.; gen. n sg.: see \textbf{ðes, ðēos, ðis}.\par
% glossary-block:B01423:end
% glossary-block:B01424:begin
\needspace{3\baselineskip}
\textbf{ðissum} — dat. m sg.; dat. n sg.; dat. pl.: see \textbf{ðes, ðēos, ðis}.\par
% glossary-block:B01424:end
% glossary-block:B01425:begin
\needspace{3\baselineskip}
\textbf{ðisum} — dat. n pl.; dat. n sg.: see \textbf{ðes, ðēos, ðis}.\par
% glossary-block:B01425:end
% glossary-block:B01426:begin
\needspace{3\baselineskip}
\textbf{ðōhte} — pret. sg.: see \textbf{ðęncan}.\par
% glossary-block:B01426:end
% glossary-block:B01427:begin
\needspace{3\baselineskip}
\textbf{ðolian} (v. weak); pret. 1 sg. ðolode; past ptcp. ðolod. suffer; endure. Accusative hardship endured.\par
% glossary-block:B01427:end
% glossary-block:B01428:begin
\needspace{3\baselineskip}
\textbf{ðon} — instr. m./n. sg.; also supplied dat. n. sg.: see \textbf{se, sēo, ðæt}.\par
% glossary-block:B01428:end
% glossary-block:B01429:begin
\needspace{3\baselineskip}
\textbf{ðone} — acc. m sg.: see \textbf{se, sēo, ðæt}.\par
% glossary-block:B01429:end
% glossary-block:B01430:begin
\needspace{3\baselineskip}
\textbf{ðonne} (conj.; adv.). than. then. Sequence of future judgment events. then; at that point.\par
% glossary-block:B01430:end
% glossary-block:B01431:begin
\needspace{3\baselineskip}
\textbf{ðonne} (conj.). when; recurring or future time.\par
% glossary-block:B01431:end
% glossary-block:B01432:begin
\needspace{3\baselineskip}
\textbf{þonne} — citation/spelling variant: see \textbf{ðonne}.\par
% glossary-block:B01432:end
% glossary-block:B01433:begin
\needspace{3\baselineskip}
\textbf{ðorfte} — pret. sg.: see \textbf{ðurfan}.\par
% glossary-block:B01433:end
% glossary-block:B01434:begin
\needspace{3\baselineskip}
\textbf{ðorn} (n. m.); gen. sg. ðornes; nom. pl. ðornas. thorn.\par
% glossary-block:B01434:end
% glossary-block:B01435:begin
\needspace{3\baselineskip}
\textbf{ðoðor} (n. m.); gen. sg. ðoðores; nom. pl. ðoðoras. ball.\par
% glossary-block:B01435:end
% glossary-block:B01436:begin
\needspace{3\baselineskip}
\textbf{ðrēan} (v.). chastise; strike.\par
% glossary-block:B01436:end
% glossary-block:B01437:begin
\needspace{3\baselineskip}
\textbf{þrēo} — nom./acc. fem./neut. pl., spelling variant: see \textbf{ðrīe}.\par
% glossary-block:B01437:end
% glossary-block:B01438:begin
\needspace{3\baselineskip}
\textbf{ðrēo} — nom./acc. fem./neut. pl.: see \textbf{ðrīe}.\par
% glossary-block:B01438:end
% glossary-block:B01439:begin
\needspace{3\baselineskip}
\textbf{ðridda} (ordinal numeral). third.\par
% glossary-block:B01439:end
% glossary-block:B01440:begin
\needspace{3\baselineskip}
\textbf{ðrīe} (cardinal numeral). three.\par
% glossary-block:B01440:end
% glossary-block:B01441:begin
\needspace{3\baselineskip}
\textbf{ðrītig} (numeral). thirty.\par
% glossary-block:B01441:end
% glossary-block:B01442:begin
\needspace{3\baselineskip}
\textbf{ðrīwa} (adv.). thrice; three times.\par
% glossary-block:B01442:end
% glossary-block:B01443:begin
\needspace{3\baselineskip}
\textbf{ðrymsetl} (n. n.); gen. sg. ðrymsetles; nom. pl. ðrymsetl. throne; throne of glory.\par
% glossary-block:B01443:end
% glossary-block:B01444:begin
\needspace{3\baselineskip}
\textbf{ðū} (personal pron.). you. Personal singular/dual/plural and governed cases; dual subjects take plural verb forms. of you. Personal genitive citation identifying singular/dual/plural; not by itself an attributive possessive lesson.\par
% glossary-block:B01444:end
% glossary-block:B01445:begin
\needspace{3\baselineskip}
\textbf{ðūhte} — pret. sg.: see \textbf{ðyncan}.\par
% glossary-block:B01445:end
% glossary-block:B01446:begin
\needspace{3\baselineskip}
\textbf{ðurfan} (v. preterite-present); pres. 1/3 sg. ðearf; pres. pl. ðurfon; pret. 1 sg. ðorfte. need to. No past-participle citation is given in the checked source paradigm.\par
% glossary-block:B01446:end
% glossary-block:B01447:begin
\needspace{3\baselineskip}
\textbf{þurh} — citation/spelling variant: see \textbf{ðurh, þurh}.\par
% glossary-block:B01447:end
% glossary-block:B01448:begin
\needspace{3\baselineskip}
\textbf{ðurh, þurh} (prep.). through. Accusative landmark or medium. through; by. Mediated action through a person in the supplied frame. by. Accusative oath standard. through; by means of. Supplied service/evidence expressions, including Ðurh hwæt.\par
% glossary-block:B01448:end
% glossary-block:B01449:begin
\needspace{3\baselineskip}
\textbf{ðurhwunian} (v. weak); pret. 1 sg. ðurhwunode; past ptcp. ðurhwunod. remain; endure; last.\par
% glossary-block:B01449:end
% glossary-block:B01450:begin
\needspace{3\baselineskip}
\textbf{ðurstig} (adj.). thirsty.\par
% glossary-block:B01450:end
% glossary-block:B01451:begin
\needspace{3\baselineskip}
\textbf{ðus} (adv.). thus.\par
% glossary-block:B01451:end
% glossary-block:B01452:begin
\needspace{3\baselineskip}
\textbf{ðūsend} (numeral). thousand. Within the complete supplied high-number expression.\par
% glossary-block:B01452:end
% glossary-block:B01453:begin
\needspace{3\baselineskip}
\textbf{þūsend} — citation/spelling variant: see \textbf{ðūsend}.\par
% glossary-block:B01453:end
% glossary-block:B01454:begin
\needspace{3\baselineskip}
\textbf{ðwēal} (n. n.). bath; bathing place.\par
% glossary-block:B01454:end
% glossary-block:B01455:begin
\needspace{3\baselineskip}
\textbf{ðwēan} (v. strong); pret. 1 sg. ðwōh; pret. pl. ðwōgon; past ptcp. ðwegen; pres. 3 sg. ðwiehð. wash. Accusative body part, mid + dative liquid.\par
% glossary-block:B01455:end
% glossary-block:B01456:begin
\needspace{3\baselineskip}
\textbf{ðwiehð} — pres. ind. 3 sg.: see \textbf{ðwēan}.\par
% glossary-block:B01456:end
% glossary-block:B01457:begin
\needspace{3\baselineskip}
\textbf{ðȳ} — instr. m sg.; instr. n sg.: see \textbf{se, sēo, ðæt}.\par
% glossary-block:B01457:end
% glossary-block:B01458:begin
\needspace{3\baselineskip}
\textbf{ðȳ-lǣs-ðe} (conj.). lest; so that not.\par
% glossary-block:B01458:end
% glossary-block:B01459:begin
\needspace{3\baselineskip}
\textbf{ðyncan} (v. impersonal); pres. 3 sg. ðyncð; pret. 3 sg. ðūhte; past ptcp. geðūht. seem. Dative experiencer; the thing that seems is nominative. seem. Nominative thing that seems, dative experiencer and agreeing predicate. seem. Supplied past source form. Impersonal use: the course supplies third-singular finite forms.\par
% glossary-block:B01459:end
% glossary-block:B01460:begin
\needspace{3\baselineskip}
\textbf{ðyrel} (n. n.); gen. sg. ðyreles; nom. pl. ðyrelu. hole.\par
% glossary-block:B01460:end
% glossary-block:B01461:begin
\needspace{3\baselineskip}
\textbf{ðȳs} — instr. m sg.; instr. n sg.: see \textbf{ðes, ðēos, ðis}.\par
% glossary-block:B01461:end
\pdfbookmark[1]{U}{outline-B01462}
% glossary-block:B01462:begin
\needspace{5\baselineskip}
\subsection*{U}
% glossary-block:B01462:end
% glossary-block:B01463:begin
\textbf{unāsęcgende} (adj.). indescribable.\par
% glossary-block:B01463:end
% glossary-block:B01464:begin
\needspace{3\baselineskip}
\textbf{unæðele} (adj.). people not noble. Substantival adjective in the supplied contrast.\par
% glossary-block:B01464:end
% glossary-block:B01465:begin
\needspace{3\baselineskip}
\textbf{unc} — dat. dual; acc. dual: see \textbf{ic}.\par
% glossary-block:B01465:end
% glossary-block:B01466:begin
\needspace{3\baselineskip}
\textbf{uncer} (possessive adjective). our; ours. Two possessors; agreeing possessive. Also gen. dual: see \textbf{ic}.\par
% glossary-block:B01466:end
% glossary-block:B01467:begin
\needspace{3\baselineskip}
\textbf{under} (prep.). under. Dative landmark in the taught static frame. under. Directional accusative in the assigned entry-under-the-roof frame.\par
% glossary-block:B01467:end
% glossary-block:B01468:begin
\needspace{3\baselineskip}
\textbf{underfōn} (v.); past ptcp. underfangen. receive; take charge of.\par
% glossary-block:B01468:end
% glossary-block:B01469:begin
\needspace{3\baselineskip}
\textbf{underðēodan} (v.). add; subjoin. Supplied passive/participial addition to speech.\par
% glossary-block:B01469:end
% glossary-block:B01470:begin
\needspace{3\baselineskip}
\textbf{underðēodnes} (n. f.). subjection; submission.\par
% glossary-block:B01470:end
% glossary-block:B01471:begin
\needspace{3\baselineskip}
\textbf{undēð} — pres. ind. 3 sg.: see \textbf{undōn}.\par
% glossary-block:B01471:end
% glossary-block:B01472:begin
\needspace{3\baselineskip}
\textbf{undōn} (v. irregular); pres. 1 sg. undō; pres. 3 sg. undēð; pres. pl. undōð; pret. 1 sg. undyde; pret. pl. undydon; past ptcp. undōn. undo; open.\par
% glossary-block:B01472:end
% glossary-block:B01473:begin
\needspace{3\baselineskip}
\textbf{undyde} — pret. sg.: see \textbf{undōn}.\par
% glossary-block:B01473:end
% glossary-block:B01474:begin
\needspace{3\baselineskip}
\textbf{unēaðe} (adv.). with difficulty; barely.\par
% glossary-block:B01474:end
% glossary-block:B01475:begin
\needspace{3\baselineskip}
\textbf{ungecnāwen} (past participle; supplied predicate); past ptcp. ungecnāwen. unknown.\par
% glossary-block:B01475:end
% glossary-block:B01476:begin
\needspace{3\baselineskip}
\textbf{ungesewenlic} (adj.). invisible.\par
% glossary-block:B01476:end
% glossary-block:B01477:begin
\needspace{3\baselineskip}
\textbf{ungeðancfull} (adj.). unthankful.\par
% glossary-block:B01477:end
% glossary-block:B01478:begin
\needspace{3\baselineskip}
\textbf{ungewunelic} (adj.). unusual; unfamiliar.\par
% glossary-block:B01478:end
% glossary-block:B01479:begin
\needspace{3\baselineskip}
\textbf{unnan} (v. preterite-present); pres. 1/3 sg. ann; pres. pl. unnon; pret. 1 sg. ūðe; past ptcp. (ge)unnen. grant; allow. Dative person, genitive thing granted.\par
% glossary-block:B01479:end
% glossary-block:B01480:begin
\needspace{3\baselineskip}
\textbf{unrihtwīs} (adj.). unrighteous. Adjective, distinct from the substantive unrighteous people. unrighteous people. Substantival adjective.\par
% glossary-block:B01480:end
% glossary-block:B01481:begin
\needspace{3\baselineskip}
\textbf{unrihtwīsnes} (n. f.); gen. sg. unrihtwīsnesse; nom. pl. unrihtwīsnessa. unrighteousness.\par
% glossary-block:B01481:end
% glossary-block:B01482:begin
\needspace{3\baselineskip}
\textbf{unrīm} (n. n.). abundance; uncounted quantity. Supplied quantity with genitive of the substance.\par
% glossary-block:B01482:end
% glossary-block:B01483:begin
\needspace{3\baselineskip}
\textbf{unscrȳdan} (v.). undress; take off. Supplied source form; retain its object/reflexive roles.\par
% glossary-block:B01483:end
% glossary-block:B01484:begin
\needspace{3\baselineskip}
\textbf{unstilnes} (n. f.). unrest; disorder.\par
% glossary-block:B01484:end
% glossary-block:B01485:begin
\needspace{3\baselineskip}
\textbf{untrum} (adj.). sick. sick people. Substantival adjective.\par
% glossary-block:B01485:end
% glossary-block:B01486:begin
\needspace{3\baselineskip}
\textbf{untrymnes} (n. f.); gen. sg. untrymnesse; nom. pl. untrymnessa. infirmity.\par
% glossary-block:B01486:end
% glossary-block:B01487:begin
\needspace{3\baselineskip}
\textbf{ūp} (adv.). up.\par
% glossary-block:B01487:end
% glossary-block:B01488:begin
\needspace{3\baselineskip}
\textbf{ūp-ā-hębban} (v. strong); pret. 1 sg. ūpāhōf; pret. pl. ūpāhōfon; past ptcp. ūpāhafen; pres. 3 sg. ūpāhęfð. lift up. Supplied source form. lift up. Supplied first-plural preterite. raise; exalt.\par
% glossary-block:B01488:end
% glossary-block:B01489:begin
\needspace{3\baselineskip}
\textbf{ūpāhęfð} — pres. ind. 3 sg.: see \textbf{ūp-ā-hębban}.\par
% glossary-block:B01489:end
% glossary-block:B01490:begin
\needspace{3\baselineskip}
\textbf{ūpāstīgan} (v.). rise up; ascend.\par
% glossary-block:B01490:end
% glossary-block:B01491:begin
\needspace{3\baselineskip}
\textbf{ūpēode} — pret. sg.: see \textbf{ūpgān}.\par
% glossary-block:B01491:end
% glossary-block:B01492:begin
\needspace{3\baselineskip}
\textbf{ūpgān} (v.). go up; come up. From the water in the supplied expression.\par
% glossary-block:B01492:end
% glossary-block:B01493:begin
\needspace{3\baselineskip}
\textbf{ūpgang} (n. m.). rising; sunrise.\par
% glossary-block:B01493:end
% glossary-block:B01494:begin
\needspace{3\baselineskip}
\textbf{upon} (prep.). on; onto.\par
% glossary-block:B01494:end
% glossary-block:B01495:begin
\needspace{3\baselineskip}
\textbf{uppan} (prep.). upon; on top of.\par
% glossary-block:B01495:end
% glossary-block:B01496:begin
\needspace{3\baselineskip}
\textbf{ūre} (possessive adjective). our; ours. Plural possessors; agreeing possessive. Also gen. pl.: see \textbf{ic}.\par
% glossary-block:B01496:end
% glossary-block:B01497:begin
\needspace{3\baselineskip}
\textbf{urnon} — pret. ind. pl.: see \textbf{iernan}.\par
% glossary-block:B01497:end
% glossary-block:B01498:begin
\needspace{3\baselineskip}
\textbf{ūs} — dat. pl.; acc. pl.: see \textbf{ic}.\par
% glossary-block:B01498:end
% glossary-block:B01499:begin
\needspace{3\baselineskip}
\textbf{ūser} — gen. pl., variant: see \textbf{ic}.\par
% glossary-block:B01499:end
% glossary-block:B01500:begin
\needspace{3\baselineskip}
\textbf{ūsic} — acc. pl., variant: see \textbf{ic}.\par
% glossary-block:B01500:end
% glossary-block:B01501:begin
\needspace{3\baselineskip}
\textbf{ūt} (adv.). out.\par
% glossary-block:B01501:end
% glossary-block:B01502:begin
\needspace{3\baselineskip}
\textbf{ūte} (adv.). outside.\par
% glossary-block:B01502:end
% glossary-block:B01503:begin
\needspace{3\baselineskip}
\textbf{ūtēode} — pret. sg.: see \textbf{ūtgān}.\par
% glossary-block:B01503:end
% glossary-block:B01504:begin
\needspace{3\baselineskip}
\textbf{ūtgān} (v.). go out. Supplied past source form. go out.\par
% glossary-block:B01504:end
% glossary-block:B01505:begin
\needspace{3\baselineskip}
\textbf{ūtgangende} (v.). come out. Supplied present participle; source permits its interpretation with earlier action. Only the cited source forms are supplied here.\par
% glossary-block:B01505:end
% glossary-block:B01506:begin
\needspace{3\baselineskip}
\textbf{ūðe} — pret. sg.: see \textbf{unnan}.\par
% glossary-block:B01506:end
\pdfbookmark[1]{W}{outline-B01507}
% glossary-block:B01507:begin
\needspace{5\baselineskip}
\subsection*{W}
% glossary-block:B01507:end
% glossary-block:B01508:begin
\textbf{wafian} (v.). be amazed. Supplied source form.\par
% glossary-block:B01508:end
% glossary-block:B01509:begin
\needspace{3\baselineskip}
\textbf{wāt} — pres. ind. 1/3 sg.: see \textbf{witan}.\par
% glossary-block:B01509:end
% glossary-block:B01510:begin
\needspace{3\baselineskip}
\textbf{wǣd} (n. f.); gen. sg. wǣde; nom. pl. wǣda. garment.\par
% glossary-block:B01510:end
% glossary-block:B01511:begin
\needspace{3\baselineskip}
\textbf{wǣdla} (n. m.); gen. sg. wǣdlan; nom. pl. wǣdlan. poor or destitute man.\par
% glossary-block:B01511:end
% glossary-block:B01512:begin
\needspace{3\baselineskip}
\textbf{wǣfels} (n. m.). cloak.\par
% glossary-block:B01512:end
% glossary-block:B01513:begin
\needspace{3\baselineskip}
\textbf{wæl} (n. n.). slaughter.\par
% glossary-block:B01513:end
% glossary-block:B01514:begin
\needspace{3\baselineskip}
\textbf{wǣpen} (n. n.); gen. sg. wǣpnes; nom. pl. wǣpnu. weapon.\par
% glossary-block:B01514:end
% glossary-block:B01515:begin
\needspace{3\baselineskip}
\textbf{wǣron} — pret. ind. pl.: see \textbf{wesan}.\par
% glossary-block:B01515:end
% glossary-block:B01516:begin
\needspace{3\baselineskip}
\textbf{wæs} — pret. ind. 1/3 sg.: see \textbf{wesan}.\par
% glossary-block:B01516:end
% glossary-block:B01517:begin
\needspace{3\baselineskip}
\textbf{wæstm} (n. m.); gen. sg. wæstmes; nom. pl. wæstmas. fruit; produce.\par
% glossary-block:B01517:end
% glossary-block:B01518:begin
\needspace{3\baselineskip}
\textbf{wæstmbǣre} (adj.). fruitful.\par
% glossary-block:B01518:end
% glossary-block:B01519:begin
\needspace{3\baselineskip}
\textbf{wæter} (n. n.); gen. sg. wæteres; nom. pl. wæteru. water.\par
% glossary-block:B01519:end
% glossary-block:B01520:begin
\needspace{3\baselineskip}
\textbf{wē} — nom. pl.: see \textbf{ic}.\par
% glossary-block:B01520:end
% glossary-block:B01521:begin
\needspace{3\baselineskip}
\textbf{wealdend} (n. m.). Ruler.\par
% glossary-block:B01521:end
% glossary-block:B01522:begin
\needspace{3\baselineskip}
\textbf{wealhstōd} (n. m.). interpreter; translator.\par
% glossary-block:B01522:end
% glossary-block:B01523:begin
\needspace{3\baselineskip}
\textbf{weall} (n. m.); gen. sg. wealles; nom. pl. weallas. wall.\par
% glossary-block:B01523:end
% glossary-block:B01524:begin
\needspace{3\baselineskip}
\textbf{weallan} (v.). surge; boil. Supplied present participle used adjectivally of the sea.\par
% glossary-block:B01524:end
% glossary-block:B01525:begin
\needspace{3\baselineskip}
\textbf{weax} (n.). wax. Pronunciation example.\par
% glossary-block:B01525:end
% glossary-block:B01526:begin
\needspace{3\baselineskip}
\textbf{weaxan} (v. strong); pret. 1 sg. wēox; pret. pl. wēoxon; past ptcp. weaxen; pres. 3 sg. wiext, wiexð. grow.\par
% glossary-block:B01526:end
% glossary-block:B01527:begin
\needspace{3\baselineskip}
\textbf{węcg} (n. m.). metal; ingot. Supplied genitive plural dependent on veins.\par
% glossary-block:B01527:end
% glossary-block:B01528:begin
\needspace{3\baselineskip}
\textbf{weg} (n. m.); gen. sg. weges; nom. pl. wegas. way; path. way; means of access. Figurative access to the Father in the declaration; distinct from physical road.\par
% glossary-block:B01528:end
% glossary-block:B01529:begin
\needspace{3\baselineskip}
\textbf{wel} (adv.); comp. bęt; superl. bętst. well. Comparative bęt: better. Superlative bętst: best.\par
% glossary-block:B01529:end
% glossary-block:B01530:begin
\needspace{3\baselineskip}
\textbf{wela} (n. m.). riches; wealth.\par
% glossary-block:B01530:end
% glossary-block:B01531:begin
\needspace{3\baselineskip}
\textbf{welig} (adj.). rich; wealthy. rich in; abounding in. on + dative thing abounded in.\par
% glossary-block:B01531:end
% glossary-block:B01532:begin
\needspace{3\baselineskip}
\textbf{wēn} (n. f.); gen. sg. wēne. expectation; chance. wēn is ðæt + subjunctive supplies the taught perhaps construction.\par
% glossary-block:B01532:end
% glossary-block:B01533:begin
\needspace{3\baselineskip}
\textbf{wēnan} (v. weak); pret. 1 sg. wēnde; past ptcp. wēned. suppose; expect; think. Genitive thing expected in the selected nominal frame; content-clause uses are separately taught.\par
% glossary-block:B01533:end
% glossary-block:B01534:begin
\needspace{3\baselineskip}
\textbf{węndan} (v.). translate. turn; return. Supplied source form.\par
% glossary-block:B01534:end
% glossary-block:B01535:begin
\needspace{3\baselineskip}
\textbf{weorc} (n. n.); gen. sg. weorces; nom. pl. weorc. work.\par
% glossary-block:B01535:end
% glossary-block:B01536:begin
\needspace{3\baselineskip}
\textbf{weorod} (n. n.). host; army.\par
% glossary-block:B01536:end
% glossary-block:B01537:begin
\needspace{3\baselineskip}
\textbf{weorpan} (v. strong); pret. 1 sg. wearp; pret. pl. wurpon; past ptcp. worpen; pres. 3 sg. wierpð. throw. Accusative thing thrown.\par
% glossary-block:B01537:end
% glossary-block:B01538:begin
\needspace{3\baselineskip}
\textbf{weorð} (n. n.); gen. sg. weorðes. price; value.\par
% glossary-block:B01538:end
% glossary-block:B01539:begin
\needspace{3\baselineskip}
\textbf{weorðan} (v. strong); pret. 1 sg. wearð; pret. pl. wurdon; past ptcp. worden; pres. 3 sg. wierð. become; be. Auxiliary of the assigned passive event with a past participle. become. Nominative subject and agreeing nominative predicate. arise; come about. Intransitive with a nominative subject, including the arising of a calm. come about; result. Supplied resultative expression.\par
% glossary-block:B01539:end
% glossary-block:B01540:begin
\needspace{3\baselineskip}
\textbf{weorðscipe} (n. m.); gen. sg. weorðscipes. honour; dignity.\par
% glossary-block:B01540:end
% glossary-block:B01541:begin
\needspace{3\baselineskip}
\textbf{wēpan} (v. strong); pret. 1 sg. wēop; pret. pl. wēopon; past ptcp. wōpen. weep. weeping; tearful. Supplied participial/adjectival use.\par
% glossary-block:B01541:end
% glossary-block:B01542:begin
\needspace{3\baselineskip}
\textbf{wer} (n. m.); nom. pl. weras. man. Supplied genitive singular. husband.\par
% glossary-block:B01542:end
% glossary-block:B01543:begin
\needspace{3\baselineskip}
\textbf{wesan} (v. irregular); pres. 1 sg. eom; pres. 3 sg. is; pres. pl. sindon; pret. 1 sg. wæs; pret. pl. wǣron; bēon-series pres. 1 sg. bēo; bēon-series pres. 3 sg. bið; bēon-series pres. pl. bēoð; negative pret. 1 sg. næs; negative pret. pl. nǣron; negative pres. 3 sg. nis. be. Copular present/past identification or state; other auxiliary constructions are separate senses. be. Passive auxiliary with the assigned participle. be. Auxiliary with intransitive past participle in the specified compound constructions. be. Auxiliary with a present participle in the supplied construction; not a universal English-progressive equivalent. be. Supplied plural present copula. seem. Only in the supplied impersonal expression it seemed to the king. not be. Supplied negative nǣre in reported supposition. would have come. Supplied past-subjunctive wesan + cumen in the counterfactual result. The checked source paradigm gives no past participle.\par
% glossary-block:B01543:end
% glossary-block:B01544:begin
\needspace{3\baselineskip}
\textbf{westan} (adv.). from the west.\par
% glossary-block:B01544:end
% glossary-block:B01545:begin
\needspace{3\baselineskip}
\textbf{wīc} (n. n.); gen. sg. wīces; nom. pl. wīc. dwelling; abode.\par
% glossary-block:B01545:end
% glossary-block:B01546:begin
\needspace{3\baselineskip}
\textbf{wician} (v. weak); pret. 1 sg. wicode; past ptcp. wicod. sojourn; lodge. Intransitive; mid + dative person with whom one lodges.\par
% glossary-block:B01546:end
% glossary-block:B01547:begin
\needspace{3\baselineskip}
\textbf{wīd} (adjective citation). wide.\par
% glossary-block:B01547:end
% glossary-block:B01548:begin
\needspace{3\baselineskip}
\textbf{wīde} (adv.); comp. wīdor; superl. wīdost. widely; far and wide. Comparative wīdor: more widely. Superlative wīdost: most widely.\par
% glossary-block:B01548:end
% glossary-block:B01549:begin
\needspace{3\baselineskip}
\textbf{wīdor} — comparative adverb: see \textbf{wīde}.\par
% glossary-block:B01549:end
% glossary-block:B01550:begin
\needspace{3\baselineskip}
\textbf{wīdost} — superlative adverb: see \textbf{wīde}.\par
% glossary-block:B01550:end
% glossary-block:B01551:begin
\needspace{3\baselineskip}
\textbf{wierpð} — pres. ind. 3 sg.: see \textbf{weorpan}.\par
% glossary-block:B01551:end
% glossary-block:B01552:begin
\needspace{3\baselineskip}
\textbf{wierrest} (adj.). worst. Superlative of yfel.\par
% glossary-block:B01552:end
% glossary-block:B01553:begin
\needspace{3\baselineskip}
\textbf{wiersa} (adj.). worse. Comparative of yfel.\par
% glossary-block:B01553:end
% glossary-block:B01554:begin
\needspace{3\baselineskip}
\textbf{wierð} — pres. ind. 3 sg.: see \textbf{weorðan}.\par
% glossary-block:B01554:end
% glossary-block:B01555:begin
\needspace{3\baselineskip}
\textbf{wierðe} (adj.). worthy.\par
% glossary-block:B01555:end
% glossary-block:B01556:begin
\needspace{3\baselineskip}
\textbf{wiext} — pres. ind. 3 sg.: see \textbf{weaxan}.\par
% glossary-block:B01556:end
% glossary-block:B01557:begin
\needspace{3\baselineskip}
\textbf{wiexð} — pres. ind. 3 sg.: see \textbf{weaxan}.\par
% glossary-block:B01557:end
% glossary-block:B01558:begin
\needspace{3\baselineskip}
\textbf{wīf} (n. n.); gen. sg. wīfes; nom. pl. wīf. wife. woman.\par
% glossary-block:B01558:end
% glossary-block:B01559:begin
\needspace{3\baselineskip}
\textbf{wilde} (adj.). wild. Reference in a source or explanatory note.\par
% glossary-block:B01559:end
% glossary-block:B01560:begin
\needspace{3\baselineskip}
\textbf{willa} (n. m.); nom. pl. willan. will; disposition.\par
% glossary-block:B01560:end
% glossary-block:B01561:begin
\needspace{3\baselineskip}
\textbf{willan} (v. irregular); pres. 1/3 sg. wille; pres. pl. willað; pret. 1 sg. wolde; pret. pl. woldon; negative pres. 3 sg. nele; negative pret. 1 sg. nolde; negative pret. pl. noldon; negative pres. 1 sg. nelle; negative pres. pl. nellað. wish; want to. be unwilling; refuse. With ne or a licensed negative contraction; followed by plain infinitive. might. Only in the exact supplied purpose construction with watching/seeing infinitives. The checked source paradigm gives no past participle.\par
% glossary-block:B01561:end
% glossary-block:B01562:begin
\needspace{3\baselineskip}
\textbf{wīn} (n. n.); gen. sg. wīnes; nom. pl. wīn. wine.\par
% glossary-block:B01562:end
% glossary-block:B01563:begin
\needspace{3\baselineskip}
\textbf{wind} (n. m.); gen. sg. windes; nom. pl. windas. wind.\par
% glossary-block:B01563:end
% glossary-block:B01564:begin
\needspace{3\baselineskip}
\textbf{windan} (v. strong); pret. 1 sg. wand; pret. pl. wundon; past ptcp. wunden. wind; plait.\par
% glossary-block:B01564:end
% glossary-block:B01565:begin
\needspace{3\baselineskip}
\textbf{wīngeard} (n. m.); gen. sg. wīngeardes; nom. pl. wīngeardas. vineyard.\par
% glossary-block:B01565:end
% glossary-block:B01566:begin
\needspace{3\baselineskip}
\textbf{winnan} (v. strong); pret. 1 sg. wann; pret. pl. wunnon; past ptcp. wunnen. struggle; toil.\par
% glossary-block:B01566:end
% glossary-block:B01567:begin
\needspace{3\baselineskip}
\textbf{winstre} (adj.). left. Supplied side phrase.\par
% glossary-block:B01567:end
% glossary-block:B01568:begin
\needspace{3\baselineskip}
\textbf{winter} (n. m.); gen. sg. wintres; nom. pl. wintras. winter.\par
% glossary-block:B01568:end
% glossary-block:B01569:begin
\needspace{3\baselineskip}
\textbf{winterlic} (adj.). winter; of winter.\par
% glossary-block:B01569:end
% glossary-block:B01570:begin
\needspace{3\baselineskip}
\textbf{wīs} (adj.); comp. wīsra. wise. wise man; wise one. Substantival direct address.\par
% glossary-block:B01570:end
% glossary-block:B01571:begin
\needspace{3\baselineskip}
\textbf{wīsdōm} (n. m.); gen. sg. wīsdōmes. wisdom.\par
% glossary-block:B01571:end
% glossary-block:B01572:begin
\needspace{3\baselineskip}
\textbf{wīsra} (adj.). wiser. Comparative of wīs.\par
% glossary-block:B01572:end
% glossary-block:B01573:begin
\needspace{3\baselineskip}
\textbf{wisse} — pret. sg., variant: see \textbf{witan}.\par
% glossary-block:B01573:end
% glossary-block:B01574:begin
\needspace{3\baselineskip}
\textbf{wiste} — pret. sg.: see \textbf{witan}.\par
% glossary-block:B01574:end
% glossary-block:B01575:begin
\needspace{3\baselineskip}
\textbf{wit} — nom. dual: see \textbf{ic}.\par
% glossary-block:B01575:end
% glossary-block:B01576:begin
\needspace{3\baselineskip}
\textbf{witan} (v. preterite-present); pres. 1/3 sg. wāt; pres. pl. witon; pret. 1 sg. wiste; past ptcp. (ge)witen. know. know. A stated fact or the answer to a question as complement. know. Accusative person in this taught frame; cunnan has a separate acquaintance sense.\par
% glossary-block:B01576:end
% glossary-block:B01577:begin
\needspace{3\baselineskip}
\textbf{wītan} (v.). blame; censure.\par
% glossary-block:B01577:end
% glossary-block:B01578:begin
\needspace{3\baselineskip}
\textbf{wīte} (n. n.). punishment; penalty.\par
% glossary-block:B01578:end
% glossary-block:B01579:begin
\needspace{3\baselineskip}
\textbf{wītga} (n. m.); gen. sg. wītgan; nom. pl. wītgan. prophet.\par
% glossary-block:B01579:end
% glossary-block:B01580:begin
\needspace{3\baselineskip}
\textbf{witodlice} (adv.). now; and. Discourse connective.\par
% glossary-block:B01580:end
% glossary-block:B01581:begin
\needspace{3\baselineskip}
\textbf{wið} (prep.). against; opposed to. Accusative opponent. by; beside; at. Accusative landmark in the taught proximity phrase. in return for. With dative; distinct from opponent/hostility. with; against. Dative fellow players in the exact supplied playing frame. with. Supplied conversational counterpart; do not infer hostility. toward. Conduct toward a person in the supplied local frame.\par
% glossary-block:B01581:end
% glossary-block:B01582:begin
\needspace{3\baselineskip}
\textbf{wiðerwinna} (n. m.). adversaries.\par
% glossary-block:B01582:end
% glossary-block:B01583:begin
\needspace{3\baselineskip}
\textbf{wiðstandan} (v. strong); pret. 1 sg. wiðstōd; pret. pl. wiðstōdon; past ptcp. wiðstanden; pres. 3 sg. wiðstęnt. withstand; resist.\par
% glossary-block:B01583:end
% glossary-block:B01584:begin
\needspace{3\baselineskip}
\textbf{wiðstęnt} — pres. ind. 3 sg.: see \textbf{wiðstandan}.\par
% glossary-block:B01584:end
% glossary-block:B01585:begin
\needspace{3\baselineskip}
\textbf{wlanc} (adj.). proud; wealthy. Contextually the rich man addressed.\par
% glossary-block:B01585:end
% glossary-block:B01586:begin
\needspace{3\baselineskip}
\textbf{wolcen} (n. n.); gen. sg. wolcnes; nom. pl. wolcnu. cloud.\par
% glossary-block:B01586:end
% glossary-block:B01587:begin
\needspace{3\baselineskip}
\textbf{wolde} — pret. sg.: see \textbf{willan}.\par
% glossary-block:B01587:end
% glossary-block:B01588:begin
\needspace{3\baselineskip}
\textbf{word} (n. n.); gen. sg. wordes; nom. pl. word. word. speech; utterance.\par
% glossary-block:B01588:end
% glossary-block:B01589:begin
\needspace{3\baselineskip}
\textbf{worhte} — pret. sg.: see \textbf{wyrcean}.\par
% glossary-block:B01589:end
% glossary-block:B01590:begin
\needspace{3\baselineskip}
\textbf{woruld} (n. f.); gen. sg. worulde; nom. pl. worulda. world.\par
% glossary-block:B01590:end
% glossary-block:B01591:begin
\needspace{3\baselineskip}
\textbf{woruldbisgu} (n. f.). worldly occupations; cares.\par
% glossary-block:B01591:end
% glossary-block:B01592:begin
\needspace{3\baselineskip}
\textbf{woruldcræft} (n. m.). worldly occupation.\par
% glossary-block:B01592:end
% glossary-block:B01593:begin
\needspace{3\baselineskip}
\textbf{woruldgeðyngð} (n. f.). worldly honour; worldly dignity.\par
% glossary-block:B01593:end
% glossary-block:B01594:begin
\needspace{3\baselineskip}
\textbf{woruldlīf} (n. n.). worldly life.\par
% glossary-block:B01594:end
% glossary-block:B01595:begin
\needspace{3\baselineskip}
\textbf{wrēon} (v. strong); pret. 1 sg. wrāh; pret. pl. wrigon; past ptcp. wrigen; pres. 3 sg. wrīhð. cover; clothe.\par
% glossary-block:B01595:end
% glossary-block:B01596:begin
\needspace{3\baselineskip}
\textbf{wrīhð} — pres. ind. 3 sg.: see \textbf{wrēon}.\par
% glossary-block:B01596:end
% glossary-block:B01597:begin
\needspace{3\baselineskip}
\textbf{wrītan} (v. strong); pret. 1 sg. wrāt; pret. pl. writon; past ptcp. writen. write. Accusative text written.\par
% glossary-block:B01597:end
% glossary-block:B01598:begin
\needspace{3\baselineskip}
\textbf{wudu} (n. m.). wood; forest. Collective vegetation, not a new material sense.\par
% glossary-block:B01598:end
% glossary-block:B01599:begin
\needspace{3\baselineskip}
\textbf{wudubearu} (n. m.). forest; grove.\par
% glossary-block:B01599:end
% glossary-block:B01600:begin
\needspace{3\baselineskip}
\textbf{wuldor} (n. n.); gen. sg. wuldres; nom. pl. wuldru. glory.\par
% glossary-block:B01600:end
% glossary-block:B01601:begin
\needspace{3\baselineskip}
\textbf{wuldrian} (v. weak); pret. 1 sg. wuldrode; past ptcp. wuldrod. glorify. Accusative person; an accusative reflexive expresses glorifying oneself.\par
% glossary-block:B01601:end
% glossary-block:B01602:begin
\needspace{3\baselineskip}
\textbf{wulf} (n. m.); gen. sg. wulfes; nom. pl. wulfas. wolf.\par
% glossary-block:B01602:end
% glossary-block:B01603:begin
\needspace{3\baselineskip}
\textbf{wund} (n. f.); gen. sg. wunde; nom. pl. wunda. wound; sore.\par
% glossary-block:B01603:end
% glossary-block:B01604:begin
\needspace{3\baselineskip}
\textbf{wundor} (n. n.); gen. sg. wundres; nom. pl. wundor. wonder.\par
% glossary-block:B01604:end
% glossary-block:B01605:begin
\needspace{3\baselineskip}
\textbf{wundrian} (v.). wonder. Supplied source form.\par
% glossary-block:B01605:end
% glossary-block:B01606:begin
\needspace{3\baselineskip}
\textbf{wunian} (v. weak); pret. 1 sg. wunode; past ptcp. wunod. dwell; remain.\par
% glossary-block:B01606:end
% glossary-block:B01607:begin
\needspace{3\baselineskip}
\textbf{wynsum} (adj.). pleasant; winsome.\par
% glossary-block:B01607:end
% glossary-block:B01608:begin
\needspace{3\baselineskip}
\textbf{wyrcean} (v. weak); pret. 1 sg. worhte; past ptcp. worht. make; build. Supplied unprefixed form; preserve explicit task controls. yield; produce. Accusative fruit in the taught/source expression. work; perform. Accusative work or deed.\par
% glossary-block:B01608:end
% glossary-block:B01609:begin
\needspace{3\baselineskip}
\textbf{wyrhta} (n. m.); gen. sg. wyrhtan; nom. pl. wyrhtan. workman; craftsman. Maker; Creator. Divine title, distinct from ordinary workman.\par
% glossary-block:B01609:end
% glossary-block:B01610:begin
\needspace{3\baselineskip}
\textbf{wyrt} (n. f.); gen. sg. wyrte; nom. pl. wyrta. herb; plant.\par
% glossary-block:B01610:end
% glossary-block:B01611:begin
\needspace{3\baselineskip}
\textbf{wyrtgemang} (n. n.); gen. sg. wyrtgemanges; nom. pl. wyrtgemang. spice; mixture of spices.\par
% glossary-block:B01611:end
% glossary-block:B01612:begin
\needspace{3\baselineskip}
\textbf{wyrttruma} (n. m.); gen. sg. wyrttruman; nom. pl. wyrttruman. root.\par
% glossary-block:B01612:end
\pdfbookmark[1]{Y}{outline-B01613}
% glossary-block:B01613:begin
\needspace{5\baselineskip}
\subsection*{Y}
% glossary-block:B01613:end
% glossary-block:B01614:begin
\textbf{yfel} (adj.); comp. wiersa; superl. wierrest. bad; evil. evil; evil people. Substantival adjective in the assigned use. evil. Substantive noun in fram yfele.\par
% glossary-block:B01614:end
% glossary-block:B01615:begin
\needspace{3\baselineskip}
\textbf{yfele} (adv.). badly; wrongly.\par
% glossary-block:B01615:end
% glossary-block:B01616:begin
\needspace{3\baselineskip}
\textbf{ymbe} (prep.). about. Approximate time, with accusative.\par
% glossary-block:B01616:end
% glossary-block:B01617:begin
\needspace{3\baselineskip}
\textbf{ymbhwyrft} (n. m.); gen. sg. ymbhwyrftes; nom. pl. ymbhwyrftas. circuit. Whole on ymbhwyrfte expression means all around; supplied dependent landmark.\par
% glossary-block:B01617:end
% glossary-block:B01618:begin
\needspace{3\baselineskip}
\textbf{ymbsealde} — pret. sg.: see \textbf{ymbsęllan}.\par
% glossary-block:B01618:end
% glossary-block:B01619:begin
\needspace{3\baselineskip}
\textbf{ymbsęllan} (v.); past ptcp. ymbseald. surround; envelop. Supplied source form. surround; beset. Supplied source form.\par
% glossary-block:B01619:end
% glossary-block:B01620:begin
\needspace{3\baselineskip}
\textbf{ymbsittan} (v.). sit around. Supplied participle referring to people sitting around.\par
% glossary-block:B01620:end
% glossary-block:B01621:begin
\needspace{3\baselineskip}
\textbf{ȳst} (n. f.); gen. sg. ȳste; nom. pl. ȳsta. storm.\par
% glossary-block:B01621:end
% glossary-block:B01622:begin
\needspace{3\baselineskip}
\textbf{ȳð} (n. f.). wave; billow. Supplied dative plural after fram.\par
% glossary-block:B01622:end
\end{multicols*}
\clearpage
\pdfbookmark[0]{Quoted errors and reconstructions}{outline-B01623}
% glossary-block:B01623:begin
\needspace{5\baselineskip}
\section*{Quoted errors and reconstructions}
% glossary-block:B01623:end
% glossary-block:B01624:begin
These notices preserve what the book quotes. A spelling rejected in a particular exercise is not thereby declared impossible in every historical context.\par
% glossary-block:B01624:end
% glossary-block:B01625:begin
\textbf{āgyltan}. Weak verb: sin; fail in duty. In Cook 34.9 this is an erroneous hint for reward; see āgyldan (strong III), whose contracted third singular is āgylt. [C06].\par
% glossary-block:B01625:end
% glossary-block:B01626:begin
\textbf{dægas}. Deliberate learner error; requested plural is dagas.\par
% glossary-block:B01626:end
% glossary-block:B01627:begin
\textbf{fiscnað}. Unselected respelling of the source word; the course retains fiscnoð.\par
% glossary-block:B01627:end
% glossary-block:B01628:begin
\textbf{fōnende}. Rejected mechanical present participle; selected form is fōnde.\par
% glossary-block:B01628:end
% glossary-block:B01629:begin
\textbf{hwælas}. Deliberate learner error; requested plural is hwalas.\par
% glossary-block:B01629:end
% glossary-block:B01630:begin
\textbf{mynsteru}. Unselected answer in the course exercise; see the selected plural under mynster. This does not rule out every historical use of the spelling.\par
% glossary-block:B01630:end
% glossary-block:B01631:begin
\textbf{sǣwa}. Reconstructed genitive plural of sǣ, marked with an asterisk in the book; it is not tested as an attested form.\par
% glossary-block:B01631:end
% glossary-block:B01632:begin
\textbf{swētrra}. Rejected unsimplified comparative genitive plural; selected form is swētra.\par
% glossary-block:B01632:end
% glossary-block:B01633:begin
\textbf{unnað}. Rejected manufactured indicative plural; selected present plural is unnon.\par
% glossary-block:B01633:end
% glossary-block:B01634:begin
\textbf{wastmas}. Deliberate learner error; requested plural is wæstmas.\par
% glossary-block:B01634:end
% glossary-block:B01635:begin
\textbf{wǣstm}. Cook’s erroneous quantity mark in the cited hint; see wæstm, with short æ. [C01].\par
% glossary-block:B01635:end
% glossary-block:B01636:begin
\textbf{wīsrra}. Rejected unsimplified comparative genitive plural; selected form is wīsra.\par
% glossary-block:B01636:end
% glossary-block:B01637:begin
\textbf{wundas}. Rejected masculine-pattern plural of the feminine wund in this context; selected plural is wunda.\par
% glossary-block:B01637:end
\end{document}
